\documentclass[fs9,seriesFre]{korfarm-base}
\usepackage{korfarm-frege}

% ============================================================
% passage-note.tex (v11)
% 지문 박스 + 첫 페이지 우측 메모란 (동적 높이)
% 정책: PAGE_BREAK_POLICY.md §3, §4 참조
%
% 변경 (v11):
%  · 지문 박스 폭 확대: enlarge right by=-3.7cm  (메모 3.4cm + gap 3mm)
%  · 메모란 동적 높이: frame.north east ~ frame.south east 사이 자동 채움
%  · 메모 줄선: clip 영역으로 박스 내부만 채움 (박스 높이 모르고도 작동)
%  · 문단 들여쓰기 활성화: tcolorbox 내부 \parindent=1em, \noindent 제거
%
% 사용:
%   \kfPassageNote{제목}{지문 본문}      → 메모란 포함
%   \kfPassageNoteNT{지문 본문}          → 메모란 포함, 제목 없음
%   \kfPassageFull{제목}{본문}           → 메모란 없음 (활동 내 보기)
% ============================================================

\makeatletter

% 메모란 규격 (cm 고정)
\newlength{\kf@memoW}
\setlength{\kf@memoW}{3.4cm}
\newlength{\kf@memoGap}
\setlength{\kf@memoGap}{3mm}

% --- 메인: 제목 있는 버전 ---
\newcommand{\kf@passagenote@with}[2]{%
  \par\ifkfPassageNewPage\ifdim\pagetotal>0.4\textheight\clearpage\fi\fi\smallskip\needspace{6\baselineskip}%
  \begin{tcolorbox}[
    enhanced, breakable,
    width=\dimexpr\linewidth-4cm\relax,
    colback=kfPassageBg, colframe=kfPrimary!70,
    boxrule=0.6pt, arc=3pt,
    left=\kfBoxPad, right=\kfBoxPad, top=\dimexpr\kfBoxPad+8pt\relax, bottom=\kfBoxPad,
    title={\small\bfseries\color{kfPrimary} #1},
    coltitle=kfPrimary,
    colbacktitle=kfPrimaryLight!60,
    attach boxed title to top left={yshift=-2.4mm, xshift=8pt},
    boxed title style={colframe=kfPrimary!50, boxrule=0.5pt, arc=2pt},
    parbox=false,
    overlay unbroken and first={%
      \begin{scope}
        \draw[kfRule, line width=0.4pt, fill=kfNoteBg, rounded corners=2pt]
          ([xshift=\kf@memoGap]frame.north east) rectangle
          ([xshift=\kf@memoGap+\kf@memoW]frame.south east);
        \node[anchor=north east, font=\footnotesize\bfseries\color{kfMute}, inner sep=3pt]
          at ([xshift=\kf@memoGap+\kf@memoW, yshift=-2pt]frame.north east) {메모};
        \node[anchor=north east, fill=kfPrimary!15, text=kfPrimary,
              font=\scriptsize\bfseries, rounded corners=2pt,
              inner xsep=4pt, inner ysep=2pt]
          at ([xshift=\kf@memoGap+\kf@memoW-3pt, yshift=-19pt]frame.north east) {\kf@memocat};
        \begin{scope}
          \path[clip]
            ([xshift=\kf@memoGap]frame.north east) rectangle
            ([xshift=\kf@memoGap+\kf@memoW]frame.south east);
          \foreach \i in {1,...,40}{%
            \draw[kfRule, line width=0.25pt]
              ([xshift=\kf@memoGap+2mm, yshift=-7mm - \i*5mm]frame.north east) --
              ([xshift=\kf@memoGap+\kf@memoW-2mm, yshift=-7mm - \i*5mm]frame.north east);%
          }
        \end{scope}
      \end{scope}
    }
  ]
    \setlength{\parindent}{1em}%
    \hspace*{1em}\ignorespaces#2%
  \end{tcolorbox}\par\smallskip
}

% --- 제목 없는 버전 ---
\newcommand{\kf@passagenote@notitle}[1]{%
  \par\ifkfPassageNewPage\ifdim\pagetotal>0.4\textheight\clearpage\fi\fi\smallskip\needspace{6\baselineskip}%
  \begin{tcolorbox}[
    enhanced, breakable,
    width=\dimexpr\linewidth-4cm\relax,
    colback=kfPassageBg, colframe=kfPrimary!70,
    boxrule=0.6pt, arc=3pt,
    left=\kfBoxPad, right=\kfBoxPad, top=\dimexpr\kfBoxPad+8pt\relax, bottom=\kfBoxPad,
    parbox=false,
    overlay unbroken and first={%
      \begin{scope}
        \draw[kfRule, line width=0.4pt, fill=kfNoteBg, rounded corners=2pt]
          ([xshift=\kf@memoGap]frame.north east) rectangle
          ([xshift=\kf@memoGap+\kf@memoW]frame.south east);
        \node[anchor=north east, font=\footnotesize\bfseries\color{kfMute}, inner sep=3pt]
          at ([xshift=\kf@memoGap+\kf@memoW, yshift=-2pt]frame.north east) {메모};
        \node[anchor=north east, fill=kfPrimary!15, text=kfPrimary,
              font=\scriptsize\bfseries, rounded corners=2pt,
              inner xsep=4pt, inner ysep=2pt]
          at ([xshift=\kf@memoGap+\kf@memoW-3pt, yshift=-19pt]frame.north east) {\kf@memocat};
        \begin{scope}
          \path[clip]
            ([xshift=\kf@memoGap]frame.north east) rectangle
            ([xshift=\kf@memoGap+\kf@memoW]frame.south east);
          \foreach \i in {1,...,40}{%
            \draw[kfRule, line width=0.25pt]
              ([xshift=\kf@memoGap+2mm, yshift=-7mm - \i*5mm]frame.north east) --
              ([xshift=\kf@memoGap+\kf@memoW-2mm, yshift=-7mm - \i*5mm]frame.north east);%
          }
        \end{scope}
      \end{scope}
    }
  ]
    \setlength{\parindent}{1em}%
    \hspace*{1em}\ignorespaces#1%
  \end{tcolorbox}\par\smallskip
}

\newcommand{\kfPassageNote}[2]{%
  \def\kf@tmptitle{#1}%
  \ifx\kf@tmptitle\@empty
    \kf@passagenote@notitle{#2}%
  \else
    \kf@passagenote@with{#1}{#2}%
  \fi
}
\newcommand{\kfPassageNoteNT}[1]{\kf@passagenote@notitle{#1}}
\makeatother

% ============================================================
% 풀폭 지문 박스 (메모란 없음) — 활동 내 보기 박스
% PAGE_BREAK_POLICY.md §3
% ============================================================
\makeatletter
\newcommand{\kf@passagefull@with}[2]{%
  \par\ifkfPassageNewPage\ifdim\pagetotal>0.4\textheight\clearpage\fi\fi\smallskip\needspace{6\baselineskip}%
  \begin{tcolorbox}[
    enhanced, breakable,
    colback=kfPassageBg, colframe=kfPrimary!70,
    boxrule=0.6pt, arc=3pt,
    left=\kfBoxPad, right=\kfBoxPad, top=\dimexpr\kfBoxPad+8pt\relax, bottom=\kfBoxPad,
    title={\small\bfseries\color{kfPrimary} #1},
    coltitle=kfPrimary,
    colbacktitle=kfPrimaryLight!60,
    attach boxed title to top left={yshift=-2.4mm, xshift=8pt},
    boxed title style={colframe=kfPrimary!50, boxrule=0.5pt, arc=2pt},
    parbox=false
  ]
    \setlength{\parindent}{1em}%
    \hspace*{1em}\ignorespaces#2%
  \end{tcolorbox}\par\smallskip
}
\newcommand{\kf@passagefull@notitle}[1]{%
  \par\ifkfPassageNewPage\ifdim\pagetotal>0.4\textheight\clearpage\fi\fi\smallskip\needspace{6\baselineskip}%
  \begin{tcolorbox}[
    enhanced, breakable,
    colback=kfPassageBg, colframe=kfPrimary!70,
    boxrule=0.6pt, arc=3pt,
    left=\kfBoxPad, right=\kfBoxPad, top=\dimexpr\kfBoxPad+8pt\relax, bottom=\kfBoxPad,
    parbox=false
  ]
    \setlength{\parindent}{1em}%
    \hspace*{1em}\ignorespaces#1%
  \end{tcolorbox}\par\smallskip
}
\newcommand{\kfPassageFull}[2]{%
  \def\kf@tmptitle{#1}%
  \ifx\kf@tmptitle\@empty
    \kf@passagefull@notitle{#2}%
  \else
    \kf@passagefull@with{#1}{#2}%
  \fi
}
\newcommand{\kfPassageFullNT}[1]{\kf@passagefull@notitle{#1}}
\makeatother

% ============================================================
% 작은 문장 박스 (문장 독해용) — PAGE_BREAK_POLICY.md §2.5
% ============================================================
\newcommand{\kfSentenceBox}[2]{%
  \par\smallskip\needspace{3\baselineskip}%
  \noindent\begin{tcolorbox}[
    enhanced, breakable=false,
    colback=kfPassageBg, colframe=kfMute,
    boxrule=0.4pt, arc=2pt,
    left=8pt, right=6pt, top=4pt, bottom=4pt,
    title={\footnotesize\bfseries\color{kfPrimary} #1},
    coltitle=kfPrimary, colbacktitle=kfPaper,
    attach boxed title to top left={yshift=-1.6mm, xshift=4pt},
    boxed title style={colframe=kfMute, boxrule=0.3pt, arc=1pt}
  ]%
    \setstretch{1.3}#2%
  \end{tcolorbox}\par\smallskip%
}

% ============================================================
% 지문 본문 아래 안내/정리 박스 (v16 신규)
% 사용: \kfPassageHint{한 줄 정리}{본문 안내 텍스트}
% ============================================================
\newcommand{\kfPassageHint}[2]{%
  \par\smallskip\needspace{4\baselineskip}%
  \begin{tcolorbox}[
    enhanced, breakable=false,
    colback=kfPrimary!5, colframe=kfPrimary!40,
    boxrule=0.4pt, arc=2pt,
    left=8pt, right=6pt, top=4pt, bottom=4pt,
    title={\footnotesize\bfseries\color{kfPrimary} #1},
    coltitle=kfPrimary, colbacktitle=kfPaper,
    attach boxed title to top left={yshift=-1.5mm, xshift=4pt},
    boxed title style={colframe=kfPrimary!40, boxrule=0.3pt, arc=1pt}
  ]%
    \small\setstretch{1.3}#2%
  \end{tcolorbox}\par\smallskip%
}

% ============================================================
% 환경 버전 (호환성) — 메모란 포함
% ============================================================
\makeatletter
\newenvironment{kfPassageNoteEnv}[1]%
  {\par\needspace{6\baselineskip}%
   \def\kf@psrc{#1}%
   \begin{tcolorbox}[
     enhanced, breakable,
     width=\dimexpr\linewidth-4cm\relax,
     colback=kfPassageBg, colframe=kfPrimary!70,
     boxrule=0.6pt, arc=3pt,
     left=\kfBoxPad, right=\kfBoxPad, top=\dimexpr\kfBoxPad+8pt\relax, bottom=\kfBoxPad,
     title={\small\bfseries\color{kfPrimary} \kf@psrc},
     coltitle=kfPrimary, colbacktitle=kfPrimaryLight!60,
     attach boxed title to top left={yshift=-2.4mm, xshift=8pt},
     boxed title style={colframe=kfPrimary!50, boxrule=0.5pt, arc=2pt},
     parbox=false,
     overlay unbroken and first={%
       \begin{scope}
         \draw[kfRule, line width=0.4pt, fill=kfNoteBg, rounded corners=2pt]
           ([xshift=\kf@memoGap]frame.north east) rectangle
           ([xshift=\kf@memoGap+\kf@memoW]frame.south east);
         \node[anchor=north east, font=\footnotesize\bfseries\color{kfMute}, inner sep=3pt]
           at ([xshift=\kf@memoGap+\kf@memoW, yshift=-2pt]frame.north east) {메모};
         \begin{scope}
           \path[clip]
             ([xshift=\kf@memoGap]frame.north east) rectangle
             ([xshift=\kf@memoGap+\kf@memoW]frame.south east);
           \foreach \i in {1,...,40}{%
             \draw[kfRule, line width=0.25pt]
               ([xshift=\kf@memoGap+2mm, yshift=-7mm - \i*5mm]frame.north east) --
               ([xshift=\kf@memoGap+\kf@memoW-2mm, yshift=-7mm - \i*5mm]frame.north east);%
           }
         \end{scope}
       \end{scope}
     }
   ]%
   \setlength{\parindent}{1em}%
  }%
  {\end{tcolorbox}\par\smallskip}
\makeatother

% ============================================================
% condition-box.tex
% <조건> 박스 컴포넌트 (tcolorbox)
% 사용:
%   \begin{kfCondition}
%     \item 첫째 조건
%     \item 둘째 조건
%   \end{kfCondition}
% ============================================================

\newenvironment{kfCondition}%
  {\par\smallskip\needspace{4\baselineskip}%
   \begin{tcolorbox}[
     enhanced, breakable=false,
     colback=kfPrimaryLight!40, colframe=kfPrimary,
     boxrule=0.5pt, arc=2pt,
     left=\kfBoxPad, right=\kfBoxPad, top=\kfBoxPad, bottom=\kfBoxPad,
     title={\small\bfseries\color{kfPrimary} \,조건\,},
     coltitle=kfPrimary, colbacktitle=kfPaper,
     attach boxed title to top left={yshift=-2mm, xshift=6pt},
     boxed title style={colframe=kfPrimary, boxrule=0.4pt, arc=1pt}
   ]%
   \begin{enumerate}[leftmargin=16pt, itemsep=1pt, topsep=1pt,
                    label=\textbf{\arabic*.}]}%
  {\end{enumerate}\end{tcolorbox}\par\smallskip}

% ============================================================
% evidence-box.tex
% <보기> 박스 컴포넌트
% 사용:
%   \begin{kfEvidence}
%     보기 본문 ...
%   \end{kfEvidence}
% ============================================================

\newenvironment{kfEvidence}%
  {\par\smallskip\needspace{4\baselineskip}%
   \begin{tcolorbox}[
     enhanced, breakable=false,
     colback=kfPaper, colframe=kfMute,
     boxrule=0.5pt, arc=2pt,
     left=\kfBoxPad, right=\kfBoxPad, top=\kfBoxPad, bottom=\kfBoxPad,
     title={\small\bfseries\color{kfInk} \,보기\,},
     coltitle=kfInk, colbacktitle=kfPrimaryLight!50,
     attach boxed title to top left={yshift=-2mm, xshift=6pt},
     boxed title style={colframe=kfMute, boxrule=0.4pt, arc=1pt}
   ]%
   \leavevmode\mbox{}\ignorespaces%
  }%
  {\end{tcolorbox}\par\smallskip}

% ============================================================
% choice-list.tex
% 5선지 (① ② ③ ④ ⑤) 자동 들여쓰기 컴포넌트
% 사용:
%   \begin{kfChoices}
%     \kfChoice 첫째 선택지
%     \kfChoice 둘째 선택지
%     \kfChoice 셋째 선택지
%     \kfChoice 넷째 선택지
%     \kfChoice 다섯째 선택지
%   \end{kfChoices}
% ============================================================

\newcounter{kfChoiceCnt}
\newcommand{\kfChoiceMark}[1]{%
  \ifcase#1\relax%
  \or ①\or ②\or ③\or ④\or ⑤\or ⑥\or ⑦\or ⑧\or ⑨%
  \fi
}

% 환경: 자동 번호
\newenvironment{kfChoices}%
  {\setcounter{kfChoiceCnt}{0}%
   \par\smallskip%
   \begin{list}{}{%
     \setlength{\leftmargin}{20pt}%
     \setlength{\itemindent}{0pt}%
     \setlength{\labelwidth}{18pt}%
     \setlength{\labelsep}{6pt}%
     \setlength{\itemsep}{2pt}%
     \setlength{\topsep}{1pt}%
     \setlength{\parsep}{0pt}%
     \renewcommand{\makelabel}[1]{##1\hss}%
   }}%
  {\end{list}\par\smallskip}

\newcommand{\kfChoice}{%
  \stepcounter{kfChoiceCnt}%
  \item[\textbf{\kfChoiceMark{\value{kfChoiceCnt}}}]\ignorespaces%
}

% --- 한 줄 5선지 (짧은 선택지용) ---
\newcommand{\kfChoicesInline}[5]{%
  \par\smallskip%
  \noindent\textbf{①}~#1\quad\textbf{②}~#2\quad\textbf{③}~#3\quad\textbf{④}~#4\quad\textbf{⑤}~#5%
  \par\smallskip%
}

% ============================================================
% answer-note.tex
% 답안 작성란 (노트 스타일, 줄친 영역)
% 사용:
%   \kfAnswerLines{5}        % 5줄 답안란
%   \kfAnswerNote{서술형 답안란}{8} % 라벨+8줄
% ============================================================

% 단일 줄 (필요 시 인라인)
\newcommand{\kfAnswerLine}{%
  \par\noindent\textcolor{kfRule}{\rule{\linewidth}{0.4pt}}\par
}

% N줄 답안란 (간단)
\newcommand{\kfAnswerLines}[1]{%
  \par\smallskip%
  \begin{minipage}{\linewidth}%
  \foreach \i in {1,...,#1}{%
    \textcolor{kfRule}{\rule{\linewidth}{0.4pt}}\\[6pt]%
  }%
  \end{minipage}%
  \par\smallskip%
}

% 라벨이 있는 노트형 답안란
\newcommand{\kfAnswerNote}[2]{%
  \par\smallskip\needspace{#2\baselineskip}%
  \begin{tcolorbox}[
    enhanced, breakable=false,
    colback=kfPaper, colframe=kfRule,
    boxrule=0.4pt, arc=2pt,
    left=\kfBoxPad, right=\kfBoxPad, top=\kfBoxPad, bottom=\kfBoxPad,
    title={\footnotesize\bfseries\color{kfMute} #1},
    coltitle=kfMute, colbacktitle=kfPaper,
    attach boxed title to top left={yshift=-1.5mm, xshift=6pt},
    boxed title style={colframe=kfRule, boxrule=0.3pt, arc=1pt}
  ]%
  \foreach \i in {1,...,#2}{%
    \textcolor{kfRule}{\rule{\linewidth}{0.3pt}}\\[6pt]%
  }%
  \end{tcolorbox}\par\smallskip%
}

% 짧은 빈칸 (단답형)
\newcommand{\kfBlank}[1][3cm]{\underline{\hspace{#1}}}
% 넓은 빈칸 (서술 한 줄)
\newcommand{\kfBlankWide}[1][6cm]{\underline{\hspace{#1}}}
% 괄호 빈칸
\newcommand{\kfBlankParen}{(\,\hspace{1cm}\,)}
% O/X 빈칸
\newcommand{\kfBlankOX}{(\,\hspace{1.5cm}\,)}

% ============================================================
% question-block.tex
% 문제 블록 (번호 배지 + 발문 + 보기/조건 + 선택지 + 답안란)
% 모든 구성을 하나의 minipage로 묶어 단/페이지 단절 금지
%
% 사용:
%   \begin{kfQBlock}{1}{객관식}
%     \kfQStem{다음 글에 대한 설명으로 가장 적절한 것은?}
%     \begin{kfEvidence} ... \end{kfEvidence}    % 선택
%     \begin{kfCondition} ... \end{kfCondition}  % 선택
%     \begin{kfChoices}
%       \kfChoice ...
%     \end{kfChoices}
%     \kfAnswerLines{2}                            % 선택
%   \end{kfQBlock}
% ============================================================

% 무단절 minipage로 묶고, 발문/보기/조건/선택지/답안란을 자유롭게 배치
% 사용자 피드백 #4: [객관식]/[서술형] 등 텍스트 라벨 제거 — 번호 배지만 표시
% #2(유형 라벨)는 호환성 인자로만 받고 출력하지 않음.
\newenvironment{kfQBlock}[2]%
  {\par\smallskip\needspace{6\baselineskip}%
   \noindent\begin{minipage}{\linewidth}%
   \samepage%
   \par\noindent
   \kfQNum{#1}\hspace{4pt}%
   \begin{minipage}[t]{\dimexpr\linewidth-30pt\relax}%
  }%
  {\end{minipage}\end{minipage}\par\smallskip}

% 문제 발문
\newcommand{\kfQStem}[1]{%
  \par\noindent#1\par\vspace{2pt}%
}

% 짧은 소문항 라벨 (1), (2), (3)
\newcounter{kfSubQ}
\newcommand{\kfSubQReset}{\setcounter{kfSubQ}{0}}
\newcommand{\kfSubQ}{%
  \stepcounter{kfSubQ}%
  \par\noindent\textbf{(\arabic{kfSubQ})}~\ignorespaces%
}

% ============================================================
% activity-table.tex
% 활동: 정리표 / 내용 확인표 (마크다운 표 → tabularx 변환)
% 사용:
%   % 일반 tabular (직접 컬럼 명시)
%   \begin{kfActTable}{|l|X|X|}
%     \kfTHead{항목} & \kfTHead{설명} & \kfTHead{학생 작성} \\\hline
%     ① & ... & \kfTBlank \\\hline
%   \end{kfActTable}
%
%   % adjustbox 자동 축소 (정해진 폭으로 강제)
%   \kfActTableWrap{
%     ... (\begin{tabular}...\end{tabular} 코드) ...
%   }
% ============================================================

% 사용 시 본문 마지막 행은 \\\hline 으로 끝나야 함.
% v17.1: 흰색 mask로 Canva 배경 본문 침범 차단
\newenvironment{kfActTable}[1]%
  {\par\smallskip\needspace{4\baselineskip}%
   \renewcommand{\arraystretch}{1.3}%
   \arrayrulecolor{kfRule}%
   \begin{tcolorbox}[blanker, colback=white, left=2pt, right=2pt, top=2pt, bottom=2pt]%
   \noindent\begin{adjustbox}{max width=\linewidth}%
   \begin{tabular}{#1}%
   \hline}%
  {\end{tabular}%
   \end{adjustbox}%
   \end{tcolorbox}%
   \arrayrulecolor{black}%
   \par\smallskip}

% 헤더 행 헬퍼
\newcommand{\kfTHead}[1]{\textbf{\color{kfPrimary} #1}}
\newcommand{\kfTHeadRow}[1]{\rowcolor{kfPrimaryLight!50}\textbf{\color{kfPrimary} #1}}

% 빈 셀 (학생 작성용)
\newcommand{\kfTBlank}{\rule[-1.6em]{0pt}{3.4em}}
\newcommand{\kfTBlankSm}{\rule[-1.0em]{0pt}{2.4em}}

% adjustbox 강제 축소 래퍼 — Python에서 변환된 raw tabular 블록을 감쌀 때 사용
\newcommand{\kfActTableWrap}[1]{%
  \par\smallskip\needspace{4\baselineskip}%
  \begin{tcolorbox}[blanker, colback=white, left=2pt, right=2pt, top=2pt, bottom=2pt]%
  \noindent\begin{adjustbox}{max width=\linewidth, center}%
  #1%
  \end{adjustbox}%
  \end{tcolorbox}\par\smallskip%
}

% ============================================================
% activity-vocab.tex
% 어휘 학습 표 (3열: 번호 - 뜻 - 빈칸)
% 사용:
%   \begin{kfVocab}
%     \kfVocabItem{1}{눈으로 보고 마음으로 깨달음}
%     \kfVocabItem{2}{사물의 이치를 따져 헤아림}
%   \end{kfVocab}
% ============================================================

% adjustbox로 폭 강제 + 일반 tabular + p{} 열 사용
\newenvironment{kfVocab}%
  {\par\smallskip\needspace{4\baselineskip}%
   \renewcommand{\arraystretch}{1.35}%
   \arrayrulecolor{kfRule}%
   \noindent\begin{adjustbox}{max width=\linewidth, center}%
   \begin{tabular}{|>{\centering\arraybackslash}m{1.0cm}|m{8.5cm}|>{\centering\arraybackslash}m{4.0cm}|}%
   \hline
   \rowcolor{kfPrimaryLight!50}%
   \multicolumn{1}{|c|}{\textbf{\color{kfPrimary}번호}} &
   \multicolumn{1}{c|}{\textbf{\color{kfPrimary}뜻풀이}} &
   \multicolumn{1}{c|}{\textbf{\color{kfPrimary}어휘 (정답 작성)}} \\\hline
   \ignorespaces}%
  {\end{tabular}%
   \end{adjustbox}%
   \arrayrulecolor{black}%
   \par\smallskip}

\newcommand{\kfVocabItem}[2]{%
  \textbf{#1} & #2 & \rule[-1.0em]{0pt}{2.4em}\\\hline%
}

% ============================================================
% activity-blank.tex
% 빈칸 채우기 활동
% 사용:
%   \begin{kfBlankList}
%     \kfBlankItem 첫 번째 문장에서 (\kfBlank)은(는) ...
%     \kfBlankItem 둘째 문장에서 ...
%   \end{kfBlankList}
% ============================================================

\newenvironment{kfBlankList}%
  {\par\smallskip%
   \begin{enumerate}[leftmargin=20pt, itemsep=4pt, topsep=2pt,
                    label=\textbf{\arabic*.}, labelsep=6pt]}%
  {\end{enumerate}\par\smallskip}

\newcommand{\kfBlankItem}{\item}

% 텍스트 안에 박스형 빈칸 (단어 1개 채우기)
\newcommand{\kfBlankBox}[1][2.4cm]{%
  \tikz[baseline=-0.5ex]{%
    \draw[kfRule, line width=0.4pt] (0,0) -- ++(#1,0);%
  }%
}

% ============================================================
% activity-ox.tex (v15)
% O/X 표기 활동 — 그래픽 디자인
% 사용:
%   \begin{kfOXList}
%     \kfOXItem 글쓴이는 자연을 예찬하고 있다.\kfOX
%     \kfOXItem 1연과 4연은 수미상관 구조이다.\kfOX
%   \end{kfOXList}
%
% \kfOX = 우측에 [O] [X] 두 박스 (학생이 하나 동그라미)
% ============================================================

\newenvironment{kfOXList}%
  {\par\smallskip%
   \begin{enumerate}[leftmargin=20pt, itemsep=5pt, topsep=2pt,
                    label=\textbf{\arabic*.}, labelsep=6pt]}%
  {\end{enumerate}\par\smallskip}

\newcommand{\kfOXItem}{\item}

% v16: O/X 그래픽 박스 키움 (18×14 → 24×20pt) + 영역색 강조
\newcommand{\kfOX}{%
  \hfill\nolinebreak
  \tikz[baseline=-0.9ex]{%
    \node[draw=kfPrimary, fill=kfPrimary!5, line width=0.6pt, rounded corners=3pt,
          minimum width=24pt, minimum height=20pt, inner sep=0pt,
          font=\normalsize\bfseries\color{kfPrimary}] (ob) {O};%
    \node[draw=kfMute, fill=kfMute!5, line width=0.6pt, rounded corners=3pt,
          minimum width=24pt, minimum height=20pt, inner sep=0pt,
          font=\normalsize\bfseries\color{kfMute},
          right=4pt of ob] (xb) {X};%
  }%
}

% ============================================================
% activity-connect.tex
% 좌우 선 긋기 활동 (2단 양식 + 중앙 여백)
% 사용:
%   \begin{kfConnect}
%     \kfPair{사르트르}{실존은 본질에 앞선다}
%     \kfPair{니체}{초인 사상}
%     \kfPair{하이데거}{현존재}
%   \end{kfConnect}
% ============================================================

\newenvironment{kfConnect}%
  {\par\smallskip%
   \renewcommand{\arraystretch}{1.6}%
   \begin{tabular}{@{}>{\raggedleft\arraybackslash}m{4.5cm}
                     @{\hspace{2.5cm}}
                     >{\raggedright\arraybackslash}m{6.5cm}@{}}}%
  {\end{tabular}\par\smallskip}

\newcommand{\kfPair}[2]{%
  \tikz[baseline]{\node[draw=kfRule, fill=kfPrimaryLight!30, rounded corners=2pt,
        inner sep=4pt, font=\small]{#1};} \tikz[baseline]{\node[circle, draw=kfMute, fill=kfPaper, inner sep=1.5pt]{};}
  &
  \tikz[baseline]{\node[circle, draw=kfMute, fill=kfPaper, inner sep=1.5pt]{};} \tikz[baseline]{\node[draw=kfRule, fill=kfNoteBg, rounded corners=2pt,
        inner sep=4pt, font=\small]{#2};} \\
}

% ============================================================
% activity-writing.tex
% 글쓰기 활동 (큰 노트 영역)
% 사용:
%   \kfWriting{독후감}{12}     % 라벨 + 12줄 큰 노트
%   \kfWritingPrompt{프롬프트 텍스트}
% ============================================================

\newcommand{\kfWritingPrompt}[1]{%
  \par\smallskip%
  \begin{tcolorbox}[
    enhanced, breakable=false,
    colback=kfPrimaryLight!30, colframe=kfPrimary,
    boxrule=0.4pt, arc=3pt,
    left=\kfBoxPad, right=\kfBoxPad, top=\kfBoxPad, bottom=\kfBoxPad,
  ]%
  \small\textbf{\color{kfPrimary}글쓰기 안내}\\[2pt]%
  #1
  \end{tcolorbox}\par\smallskip%
}

\newcommand{\kfWriting}[2]{%
  \par\smallskip\needspace{\numexpr#2+2\relax\baselineskip}%
  \begin{tcolorbox}[
    enhanced, breakable=false,
    colback=kfPaper, colframe=kfRule,
    boxrule=0.4pt, arc=2pt,
    left=\kfBoxPad, right=\kfBoxPad, top=\kfBoxPad, bottom=\kfBoxPad,
    title={\footnotesize\bfseries\color{kfMute} #1},
    coltitle=kfMute, colbacktitle=kfPaper,
    attach boxed title to top left={yshift=-1.5mm, xshift=6pt},
    boxed title style={colframe=kfRule, boxrule=0.3pt, arc=1pt}
  ]%
  \foreach \i in {1,...,#2}{%
    \textcolor{kfRule}{\rule{\linewidth}{0.3pt}}\\[7pt]%
  }%
  \end{tcolorbox}\par\smallskip%
}

% ============================================================
% activity-structure.tex
% 구조도 (트리 + 빈칸)
% 사용:
%   \begin{kfStructure}
%     \kfNode{0}{서론}{문제 제기}
%     \kfNode{1}{본론}{(\,\quad\,)}
%     \kfNode{1}{본론}{근거 2}
%     \kfNode{0}{결론}{주장 강화}
%   \end{kfStructure}
% ============================================================

\newenvironment{kfStructure}%
  {\par\smallskip\needspace{6\baselineskip}%
   \begin{tcolorbox}[
     enhanced, breakable=false,
     colback=kfPaper, colframe=kfRule,
     boxrule=0.4pt, arc=2pt,
     left=\kfBoxPad, right=\kfBoxPad, top=\kfBoxPad, bottom=\kfBoxPad,
     title={\footnotesize\bfseries\color{kfMute} 글의 구조},
     coltitle=kfMute, colbacktitle=kfPrimaryLight!40,
     attach boxed title to top left={yshift=-1.5mm, xshift=6pt},
     boxed title style={colframe=kfRule, boxrule=0.3pt, arc=1pt}
   ]%
   \begin{tabbing}
     \hspace{1.4cm}\=\hspace{1.4cm}\=\hspace{8cm}\kill}%
  {\end{tabbing}\end{tcolorbox}\par\smallskip}

% 노드: #1=들여쓰기 단계 (0~3), #2=라벨, #3=내용
\newcommand{\kfNode}[3]{%
  \ifcase#1\relax%
    \>\>\textbf{\color{kfPrimary} ▶ #2}\quad #3\\
  \or
    \>\>\quad\textbf{\color{kfMute} ◦ #2}\quad #3\\
  \or
    \>\>\quad\quad\textbf{\color{kfMute} - #2}\quad #3\\
  \or
    \>\>\quad\quad\quad\textbf{\color{kfMute} \textperiodcentered\ #2}\quad #3\\
  \fi
}

% 빈칸 노드
\newcommand{\kfNodeBlank}[2]{\kfNode{#1}{#2}{(\hspace{2.5cm})}}

% ============================================================
% activity-sentence-reading.tex
% 문장 독해 활동 — 문장 박스 1개 + 그에 묶인 문제 N개를 한 minipage로
% 사용:
%   \begin{kfSentenceUnit}{1}{"바람은 내 귀에 속삭이며..."}
%     \kfSentQ{(1)}{서술어 '속삭이며'의 주어는?}
%     \begin{kfChoices}
%       \kfChoice 바람은
%       ...
%     \end{kfChoices}
%   \end{kfSentenceUnit}
% ============================================================

% kfSentenceBox 매크로는 passage-note.tex에 정의됨

\newenvironment{kfSentenceUnit}[2]%
  {\par\smallskip\needspace{6\baselineskip}%
   \noindent\begin{minipage}{\linewidth}%
   \samepage%
   \kfSentenceBox{문장 #1}{#2}%
   \par\smallskip%
  }%
  {\end{minipage}\par\smallskip}

% 같은 문장에 묶인 소문항 (문장 안내 + 발문)
\newcommand{\kfSentQ}[2]{%
  \par\noindent\textbf{#1}~#2\par\smallskip%
}

% 헤더 없이 문장만 있는 묶음 (sentence_ref가 없는 일반 문항)
\newenvironment{kfSentencelessUnit}%
  {\par\smallskip\noindent\begin{minipage}{\linewidth}\samepage}%
  {\end{minipage}\par\smallskip}

% ============================================================
% chapter-cover.tex (v6)
% 챕터 진입 페이지 (표지) — 하단에 영역 목차 추가
% 사용:
%   \kfChapterCover{1}{근대성과 자아}{Chapter 1}{이번 챕터에서는 ...}
%   \begin{kfChapterAreaList}
%     \kfChapterAreaItem{어휘}{kfAreaVocab}{핵심 어휘 12개}
%     ...
%   \end{kfChapterAreaList}
%
%   영역 목차가 없을 때는 \kfChapterCoverEnd 호출
% ============================================================

\newcommand{\kfChapterCover}[4]{%
  \clearpage%
  \thispagestyle{kfBlank}%
  \kfMarkChapter{Chapter #1. #2}%
  \kfChapterCoverBg%
  % v18.5: 챕터 표지 우상단에 로고
  \begin{tikzpicture}[remember picture, overlay]
    \node[anchor=north east, inner sep=0pt] at ($(current page.north east)+(-18mm,-18mm)$) {\kfLogoMd};
  \end{tikzpicture}%
  \vspace*{20mm}%
  \noindent
  \begin{tikzpicture}[remember picture, overlay]
    % 좌측 액센트 바
    \fill[kfPrimary] (current page.west) ++(12mm,25mm) rectangle ++(5mm, -50mm);
    % 챕터 번호 배지 (이중 원, 약간 작게)
    \node[anchor=north west, inner sep=0pt] at ($(current page.north west)+(24mm,-45mm)$) {%
      \begin{tikzpicture}
        \draw[kfPrimary, line width=1pt] (0,0) circle (10mm);
        \draw[kfPrimary, line width=0.4pt] (0,0) circle (8mm);
        \node[font=\normalsize\bfseries\color{kfPrimary}] at (0,2mm) {CHAPTER};
        \node[font=\LARGE\bfseries\color{kfPrimary}] at (0,-3mm) {#1};
      \end{tikzpicture}%
    };
  \end{tikzpicture}%
  \vspace{42mm}%
  \par\noindent\hspace{32mm}%
  \begin{minipage}{0.65\linewidth}%
    {\footnotesize\color{kfMute} #3}\\[4pt]
    {\LARGE\bfseries\color{kfPrimary} #2}\\[8pt]
    {\color{kfRule}\rule{50mm}{0.5pt}}\\[6pt]
    {\small\color{kfInk} #4}
  \end{minipage}%
  \par\vspace{14mm}%
}

% v6 / v18.5 / v18.6: 챕터 표지의 영역 목차 — 흰색 반투명 박스
% v18.6: 배경 가로선/도형과 안 겹치도록 TikZ overlay로 우하단 절대 위치
\makeatletter
% 영역 항목들을 모아 둘 박스
\newsavebox{\kf@arealistbox}
\newenvironment{kfChapterAreaList}{%
  \begin{lrbox}{\kf@arealistbox}%
  \begin{minipage}{0.62\linewidth}%
    \begin{tcolorbox}[%
      enhanced, colback=white, colframe=white,
      opacityback=0.92, opacityframe=0,
      boxrule=0pt, arc=4pt,
      width=\linewidth,
      top=8pt, bottom=8pt, left=10pt, right=10pt,
    ]%
      {\footnotesize\bfseries\color{kfMute} 이 챕터의 영역}\par\vspace{4pt}%
      {\color{kfRule}\hrule height 0.3pt}\par\vspace{4pt}%
      \begin{itemize}[leftmargin=14pt, itemsep=2pt, topsep=0pt, label={\color{kfPrimary!70}\textbullet}]%
}{%
      \end{itemize}%
    \end{tcolorbox}%
  \end{minipage}%
  \end{lrbox}%
  % v18.7: 배경 상단 가로선 ~ 하단 가로선 사이의 중간 영역에 배치
  % 가로선이 내용을 가리지 않도록 박스 중심이 두 선 사이 중점에 위치
  \begin{tikzpicture}[remember picture, overlay]%
    \node[anchor=center, inner sep=0pt]
      at ($(current page.south)+(0mm,90mm)$)
      {\usebox{\kf@arealistbox}};%
  \end{tikzpicture}%
  \vfill%
  \noindent\hfill\kfSeriesBadge\hspace{12mm}%
  \clearpage%
}
\makeatother

% 영역 항목: \kfChapterAreaItem{영역명}{kfArea색}{부제}
\makeatletter
\newcommand{\kfChapterAreaItem}[3]{%
  \item {\small\bfseries\color{#2} #1}%
  \def\kf@cci@sub{#3}%
  \ifx\kf@cci@sub\@empty\else
    {\small\color{kfMute}\, \textemdash\, #3}%
  \fi
}
\makeatother

% 영역 목록 없이 cover만 (legacy)
\newcommand{\kfChapterCoverEnd}{%
  \vfill%
  \noindent\hfill\kfSeriesBadge\hspace{12mm}%
  \clearpage%
}

% ============================================================
% area-header.tex (v16)
% 영역 시작 표제 — 시리즈별 차별화 라벨 + 큰 영역명 + 부제
% PAGE_BREAK_POLICY.md §1.4
%
% v16 (디자인 고도화 Phase 1):
%   - 시리즈별 라벨 박스 추가:
%     소쉬르 = AREA START (영역색 채움 박스)
%     프레게 = SECTION OPEN (영역색 테두리)
%     러셀   = RUSSELL (영역색 라이트 박스)
%     비트   = WITTGENSTEIN (회색 라이트 박스)
%   - 라벨 위치: 영역명 위 좌측
% ============================================================

\makeatletter

% 시리즈 라벨 텍스트
\newcommand{\kf@serieslabel}{%
  \ifcase\kf@series\relax
    AREA START%
  \or
    AREA START%
  \or
    SECTION OPEN%
  \or
    RUSSELL%
  \or
    WITTGENSTEIN%
  \fi
}

% 시리즈 라벨 박스 스타일 (#1 = 영역색)
\newcommand{\kf@labelbox}[1]{%
  \ifcase\kf@series\relax
    % 0/1 = 소쉬르: 영역색 채움
    \tikz[baseline=-2pt]{\node[fill=#1, text=white,
      rounded corners=3pt, inner xsep=6pt, inner ysep=2pt,
      font=\scriptsize\bfseries]{\kf@serieslabel};}%
  \or
    % 1 = 소쉬르: 영역색 채움
    \tikz[baseline=-2pt]{\node[fill=#1, text=white,
      rounded corners=3pt, inner xsep=6pt, inner ysep=2pt,
      font=\scriptsize\bfseries]{\kf@serieslabel};}%
  \or
    % 2 = 프레게: 영역색 테두리
    \tikz[baseline=-2pt]{\node[draw=#1, line width=0.6pt, text=#1,
      rounded corners=3pt, inner xsep=6pt, inner ysep=2pt,
      font=\scriptsize\bfseries]{\kf@serieslabel};}%
  \or
    % 3 = 러셀: 영역색 라이트 박스
    \tikz[baseline=-2pt]{\node[fill=#1!12, text=#1,
      rounded corners=2pt, inner xsep=5pt, inner ysep=1.5pt,
      font=\scriptsize\bfseries]{\kf@serieslabel};}%
  \or
    % 4 = 비트: 회색 미니 박스
    \tikz[baseline=-2pt]{\node[fill=kfMute!10, text=kfMute,
      rounded corners=2pt, inner xsep=5pt, inner ysep=1.5pt,
      font=\scriptsize\bfseries]{\kf@serieslabel};}%
  \fi
}

% v17: 시리즈+영역별 Canva 배경 자동 매핑
% \kf@hasbg=1로 set되면 \kfAreaStart에서 라벨 박스 출력 생략
\def\kf@areaVocab{어휘}\def\kf@areaGrammar{문법}\def\kf@areaConcept{개념}
\def\kf@areaLit{문학}\def\kf@areaNonlit{비문학}
\newif\ifkf@bgon
% 파일 있으면 배경 적용 + boolean true. 없으면 nothing.
\newcommand{\kfBgSet}[1]{%
  \IfFileExists{#1.pdf}{\kfPageBg{#1}\kf@bgontrue}{}%
}
\newcommand{\kf@trybg}[1]{%
  % v18.4: 영역 배경 PDF 비활성화 — 큰 도형이 본문 영역을 침범해
  % "영역 헤더 후 빈 페이지" 문제를 일으킴. 라벨 박스로만 영역 표시.
  \kf@bgonfalse%
}

% Note: 러셀(rus_*)/비트(wit_*) 배경 PDF 미존재 시 \kfPageBg가 에러 → \IfFileExists로 가드
% (현재는 파일이 모두 존재한다고 가정. 부분 제공 시 cls 수정 필요)

\newcommand{\kfAreaStart}[3]{%
  \clearpage%
  \thispagestyle{fancy}%
  \kfMarkArea{#1}%
  \kf@trybg{#1}%
  \vspace*{1mm}%
  % 배경 없으면 라텍스 라벨 박스 출력 (배경 있으면 일러스트가 라벨 역할)
  \ifkf@bgon
    \vspace*{4mm}%
  \else
    \noindent\kf@labelbox{#2}\par\vspace{4pt}%
  \fi
  % 영역명 + 부제
  \noindent
  \begin{tikzpicture}
    \node[anchor=west, inner sep=0pt] (bar) at (0,0) {\color{#2}\rule{3pt}{22pt}};
    \node[anchor=base west, inner sep=0pt, xshift=8pt, yshift=2pt] (lbl) at (bar.east |- 0,0) {%
      {\LARGE\bfseries\color{#2} #1}%
    };
    \def\kf@areasub{#3}%
    \ifx\kf@areasub\@empty\else
      \node[anchor=base west, inner sep=0pt, xshift=8pt, font=\normalsize\color{kfMute}]
        at (lbl.base east) {\,\textperiodcentered\, #3};%
    \fi
  \end{tikzpicture}%
  \par\vspace{2pt}
  {\color{#2!70}\hrule height 1pt}%
  \par\vspace{1pt}%
  {\color{kfRule}\hrule height 0.3pt}%
  \par\vspace{4pt}%
  \needspace{6\baselineskip}\nopagebreak[4]%
  \global\kf@areaJustOpenedtrue% v18.5: 첫 컨텐츠 needspace 무시
}
\makeatother

% 시리즈/영역 단축 매크로
\newcommand{\kfStartVocab}[1]{\kfAreaStart{어휘}{kfAreaVocab}{#1}}
\newcommand{\kfStartGrammar}[1]{\kfAreaStart{문법}{kfAreaGrammar}{#1}}
\newcommand{\kfStartConcept}[1]{\kfAreaStart{개념}{kfAreaConcept}{#1}}
\newcommand{\kfStartLit}[1]{\kfAreaStart{문학}{kfAreaLiterature}{#1}}
\newcommand{\kfStartNonlit}[1]{\kfAreaStart{비문학}{kfAreaNonlit}{#1}}
\newcommand{\kfStartWeekly}[1]{\kfAreaStart{실력 확인}{kfAreaWeekly}{#1}}

% ============================================================
% section-header.tex
% 섹션 시작 라벨 헤더 — [지문] / [활동] / [문제] 라벨
% 좌측 액센트 바 + 라벨 + 부제 (영역 액센트 색)
%
% 사용:
%   \kfSecHeader{kfAreaLiterature}{지문}{빼앗긴 들에도 봄은 오는가 / 이상화}
%   \kfSecHeader{kfAreaConcept}{활동}{내용 확인}
%   \kfSecHeader{kfAreaNonlit}{문제}{}
%
% 헤더 직후 페이지 분할 금지: \needspace{4\baselineskip} + \nopagebreak
% ============================================================

\providecommand{\kfSecHeader}[3]{%
  \par\needspace{10\baselineskip}\filbreak%
  \vspace{3pt}%
  \noindent
  \begin{tikzpicture}
    \node[anchor=west, inner sep=0pt] (bar) at (0,0) {\color{#1}\rule{3pt}{14pt}};
    \node[anchor=west, inner sep=0pt, xshift=6pt, font=\normalsize\bfseries\color{#1}]
       (lbl) at (bar.east) {[\,#2\,]};
    \node[anchor=west, inner sep=0pt, xshift=4pt, font=\small\color{kfMute}]
       at (lbl.east) {#3};
  \end{tikzpicture}%
  \par\vspace{1pt}%
  {\color{#1!50}\hrule height 0.4pt}%
  \par\vspace{2pt}\nopagebreak%
}

% 영역별 단축 매크로
\newcommand{\kfSecPassageVocab}[1]{\kfSecHeader{kfAreaVocab}{지문}{#1}}
\newcommand{\kfSecPassageGrammar}[1]{\kfSecHeader{kfAreaGrammar}{지문}{#1}}
\newcommand{\kfSecPassageConcept}[1]{\kfSecHeader{kfAreaConcept}{지문}{#1}}
\newcommand{\kfSecPassageLit}[1]{\kfSecHeader{kfAreaLiterature}{지문}{#1}}
\newcommand{\kfSecPassageNonlit}[1]{\kfSecHeader{kfAreaNonlit}{지문}{#1}}
\newcommand{\kfSecPassageWeekly}[1]{\kfSecHeader{kfAreaWeekly}{지문}{#1}}

\newcommand{\kfSecActVocab}[1]{\kfSecHeader{kfAreaVocab}{활동}{#1}}
\newcommand{\kfSecActGrammar}[1]{\kfSecHeader{kfAreaGrammar}{활동}{#1}}
\newcommand{\kfSecActConcept}[1]{\kfSecHeader{kfAreaConcept}{활동}{#1}}
\newcommand{\kfSecActLit}[1]{\kfSecHeader{kfAreaLiterature}{활동}{#1}}
\newcommand{\kfSecActNonlit}[1]{\kfSecHeader{kfAreaNonlit}{활동}{#1}}
\newcommand{\kfSecActWeekly}[1]{\kfSecHeader{kfAreaWeekly}{활동}{#1}}

\newcommand{\kfSecQVocab}[1]{\kfSecHeader{kfAreaVocab}{문제}{#1}}
\newcommand{\kfSecQGrammar}[1]{\kfSecHeader{kfAreaGrammar}{문제}{#1}}
\newcommand{\kfSecQConcept}[1]{\kfSecHeader{kfAreaConcept}{문제}{#1}}
\newcommand{\kfSecQLit}[1]{\kfSecHeader{kfAreaLiterature}{문제}{#1}}
\newcommand{\kfSecQNonlit}[1]{\kfSecHeader{kfAreaNonlit}{문제}{#1}}
\newcommand{\kfSecQWeekly}[1]{\kfSecHeader{kfAreaWeekly}{문제}{#1}}


\begin{document}

\thispagestyle{empty}
\vspace*{40mm}
\begin{center}
{\Huge\bfseries\color{kfPrimary} 프레게3}\\[10pt]
{\Large\color{kfMute} 학생용 교재 — 제 1 권}\\[20pt]
{\small\color{kfMute} (챕터 1 \textendash{} 4)}
\end{center}
\vfill
\hfill\kfSeriesBadge\hspace{15mm}
\clearpage
\kfChapterCover{1}{프레게3 (챕터1)}{Chapter 1}{이번 챕터의 학습 내용을 시작합니다.}

\kfStartConcept{개념 영역 — 다른 것에 빗대어 표현하기}


\kfSecHeader{kfAreaConcept}{개념}{다른 것에 빗대어 표현하기}

비유법은 표현하고 싶은 것을 다른 것에 빗대어 더 생생하고 재미있게 나타내는 방법이에요. 이때 진짜 말하고 싶은 대상을 원관념, 빗대어 표현하는 대상을 보조관념이라고 해요. 원관념과 보조관념은 서로 비슷한 점이 있기 때문에 비교할 수 있답니다.

직유법은 '같이', '처럼', '-듯이', '-인 양' 등을 사용하여 원관념과 보조관념을 직접적으로 연결하는 비유법이에요. "가르마 같은 논길을 따라"에서 논길을 가르마에 비유해서 일직선으로 뻗은 모습을 표현하거나, "내 누님같이 생긴 꽃이여"처럼 국화꽃을 누님에 빗대어 꽃의 고운 모습을 나타내요. 직유법은 비교하는 말이 직접 나와 있어서 이해하기 쉬워요.

은유법은 'A는 B이다'의 형태로 원관념과 보조관념을 간접적으로 견주어 표현하는 비유법이에요. "내 마음은 호수요"처럼 마음이 호수처럼 고요하다는 뜻이 숨어 있어서 더 깊이 생각해봐야 해요. "이것은 소리 없는 아우성"처럼 깃발을 아우성에 빗대어 힘차게 펄럭이는 모습을 표현하기도 해요.

의인법은 사람이 아닌 것을 사람처럼 표현하는 방법이에요. "풀이 눕는다"나 "말갛게 씻은 얼굴 고운 해야 솟아라"처럼 풀이나 해를 마치 사람처럼 행동하는 것으로 표현해서 친근하고 생생하게 느껴지도록 만들어요. "민주주의여"라고 부르며 민주주의를 사람처럼 대하는 것도 의인법의 한 예랍니다.

이런 비유법들을 사용하면 평범한 말보다 훨씬 더 아름답고 감동적으로 표현할 수 있어요. 독자들도 시나 소설을 읽으면서 더 재미있고 깊이 있게 느낄 수 있답니다.
\par\medskip\begin{itemize}[leftmargin=18pt, itemsep=2pt, topsep=2pt, label=\textbullet]
\item 직유법 — '가르마 같은 논길'(논길↔가르마, 연결어 '같은')
\item 은유법 — '이것은 소리 없는 아우성'(깃발↔아우성, A는 B이다)
\item 의인법 — '풀이 눕는다 / 민주주의여!'(무생물·추상물을 사람처럼)
\end{itemize}

\kfSecHeader{kfAreaConcept}{문제}{}

\begin{kfQBlock}{1}{}
\kfQStem{다음 중, ‘처럼’, ‘같이’, ‘-듯이’과 같은 연결어를 사용하여 대상을 직접 빗대어 표현한 직유법이 사용된 것은 무엇인가요?}
\begin{kfChoices}
\kfChoice 구름에 달 가듯 가는 나그네
\kfChoice 해야 솟아라.
\kfChoice 내 마음은 호수요.
\kfChoice 이것은 소리 없는 아우성.
\kfChoice 풀이 눕는다.
\end{kfChoices}
\end{kfQBlock}
\begin{kfQBlock}{2}{}
\kfQStem{다음 중, ‘A는 B이다’의 형태로 두 대상을 간접적으로 연결한 은유법이 사용된 것은 무엇인가요?}
\begin{kfChoices}
\kfChoice 가르마 같은 논길
\kfChoice 인연은 갈밭을 건너는 바람
\kfChoice 꽃잎이 지는 어느 날
\kfChoice 내 누님같이 생긴 꽃이여.
\kfChoice 밥티처럼 따스한 별
\end{kfChoices}
\end{kfQBlock}
\begin{kfQBlock}{3}{}
\kfQStem{다음 중, 사람이 아닌 것을 사람처럼 표현한 의인법이 사용된 것은 무엇인가요?}
\begin{kfChoices}
\kfChoice 모든 산맥들이 바다를 사랑해 달려갈 때에도
\kfChoice 밥티처럼 따스한 별
\kfChoice 내 누님같이 생긴 꽃이여.
\kfChoice 세상 같은 건 더러워서 버리는 것이다.
\kfChoice 내 마음은 호수요.
\end{kfChoices}
\end{kfQBlock}
\begin{kfQBlock}{4}{}
\kfQStem{김동명의 시 「내 마음은」에 나오는 “내 마음은 호수요”라는 표현에 대한 설명으로 적절하지 \uline{않은} 것은 무엇인가요?}
\begin{kfChoices}
\kfChoice 'A는 B이다' 형태의 은유법이 사용되었다.
\kfChoice '내 마음'은 원관념에 해당한다.
\kfChoice '호수'는 보조관념에 해당한다.
\kfChoice '처럼', '같이'를 사용한 직유법이 쓰였다.
\kfChoice 추상적인 마음을 구체적인 대상에 빗대어 표현하였다.
\end{kfChoices}
\end{kfQBlock}
\begin{kfQBlock}{5}{}
\kfQStem{비유법에 대한 설명으로 알맞지 \uline{않은} 것은 무엇인가요?}
\begin{kfChoices}
\kfChoice 직유법, 은유법, 의인법은 같은 방식으로 대상을 표현하는 기법이다.
\kfChoice 직유법은 은유법보다 의미를 이해하기가 한결 쉽다.
\kfChoice 의인법은 사물이 친근하고 생생하게 느껴지게 해 준다.
\kfChoice 은유법은 일반적으로 'A는 B이다'와 같은 형태로 주로 나타난다.
\kfChoice 평범한 대상을 더 생생하고 재미있게 표현할 수 있다.
\end{kfChoices}
\end{kfQBlock}
\begin{kfQBlock}{6}{}
\kfQStem{김수영의 시 「풀」에 나오는 “풀은 눕고 / 드디어 울었다”와 같이 의인법을 사용하여 얻는 효과는 무엇일까요?}
\begin{kfChoices}
\kfChoice 시의 의미를 일부러 숨겨 풀이를 어렵게 만든다.
\kfChoice 대상에 생명력을 부여해 친근하고 생생하게 느끼게 한다.
\kfChoice 시에 그려진 풍경이 모두 사실이라고 믿게 만든다.
\kfChoice 독자가 시인의 주장을 비판적으로 더 검토하게 만든다.
\kfChoice 시의 분위기를 훨씬 가볍고 유머러스하게 만든다.
\end{kfChoices}
\end{kfQBlock}
\begin{kfQBlock}{7}{}
\kfQStem{다음 시 구절과 그에 사용된 비유법의 연결이 알맞지 \uline{않은} 것은 무엇인가요?}
\begin{kfChoices}
\kfChoice 이것은 소리 없는 아우성 - 의인법
\kfChoice 내 마음은 호수요 / 은유법 사용
\kfChoice 가르마 같은 논길 / 직유법 사용
\kfChoice 밥티처럼 따스한 별 / 직유법 사용
\kfChoice 해야 솟아라 / 의인법 사용
\end{kfChoices}
\end{kfQBlock}
\begin{kfQBlock}{8}{}
\kfQStem{비유법에서, 진짜 말하고 싶은 대상인 '원관념'과 그것을 빗대기 위해 사용된 '보조관념'의 관계로 알맞은 것은 무엇인가요?}
\begin{kfChoices}
\kfChoice 원관념을 통해 보조관념의 특징을 설명한다.
\kfChoice 보조관념을 통해 원관념의 특징을 더 생생하게 드러낸다.
\kfChoice 원관념이 보조관념보다 항상 더 아름다운 대상이다.
\kfChoice 원관념과 보조관념은 아무런 관련이 없는 단어들이다.
\kfChoice 원관념과 보조관념은 항상 똑같은 의미를 가진다.
\end{kfChoices}
\end{kfQBlock}
\begin{kfQBlock}{9}{}
\kfQStem{다음 중, 사용된 비유법의 종류가 나머지 넷과 \uline{다른} 하나는 무엇인가요?}
\begin{kfChoices}
\kfChoice 내 누님같이 생긴 꽃이여
\kfChoice 삼단 같은 머리를 감았구나
\kfChoice 인연은 갈밭을 건너는 바람
\kfChoice 구름에 달 가듯이 가는 나그네
\kfChoice 그는 마치 어린애인 양 행동했다.
\end{kfChoices}
\end{kfQBlock}
\begin{kfQBlock}{10}{}
\kfQStem{박두진의 시 「해」에서, “말갛게 씻은 얼굴 고운 해야 솟아라”와 같이 해를 사람처럼 표현하여 얻는 효과는 무엇일까요?}
\begin{kfChoices}
\kfChoice ‘해’라는 존재를 더 가깝고 친근하게 느끼게 한다.
\kfChoice 해가 얼마나 뜨거운지 알려준다.
\kfChoice 해가 뜨는 과학적 원리를 설명한다.
\kfChoice 해가 뜨는 것을 막고 싶다는 시인의 마음을 표현한다.
\kfChoice 시의 분위기를 슬프게 만든다.
\end{kfChoices}
\end{kfQBlock}
\begin{kfQBlock}{1}{}
\kfQStem{표현하고 싶은 것을 다른 것에 빗대어 더 생생하고 재미있게 나타내는 방법을 무엇이라고 하나요?}
\kfAnswerLines{1}
\end{kfQBlock}
\begin{kfQBlock}{2}{}
\kfQStem{비유법에서, 진짜 말하고 싶은 원래의 대상을 무엇이라고 하나요?}
\kfAnswerLines{1}
\end{kfQBlock}
\begin{kfQBlock}{3}{}
\kfQStem{비유법에서, 빗대어 표현하는 데 사용되는 대상을 무엇이라고 하나요?}
\kfAnswerLines{1}
\end{kfQBlock}
\begin{kfQBlock}{4}{}
\kfQStem{'처럼', '같이', '-듯이'과 같은 말을 사용하여 두 대상을 직접 연결하는 비유법은 무엇인가요?}
\kfAnswerLines{1}
\end{kfQBlock}
\begin{kfQBlock}{5}{}
\kfQStem{"내 마음은 호수요"처럼, 'A는 B이다'의 형태로 두 대상을 간접적으로 연결하는 비유법은 무엇인가요?}
\kfAnswerLines{1}
\end{kfQBlock}
\begin{kfQBlock}{6}{}
\kfQStem{"풀이 눕는다"처럼, 사람이 아닌 것을 사람처럼 표현하여 친근하게 만드는 방법을 무엇이라고 하나요?}
\kfAnswerLines{1}
\end{kfQBlock}
\begin{kfQBlock}{7}{}
\kfQStem{"가르마 같은 논길"이라는 표현에 사용된 비유법은 무엇인가요?}
\kfAnswerLines{1}
\end{kfQBlock}
\begin{kfQBlock}{1}{}
\kfQStem{직유법과 은유법은 대상을 빗대어 표현하는 방식에서 어떤 차이가 있는지 서술하세요.}
\kfAnswerNote{답안}{4}
\end{kfQBlock}


\kfStartLit{문학 영역 — 작품}


\kfSecHeader{kfAreaLiterature}{지문}{작품}
\kfPassageNote{작품}{%
나는 이제 너에게도 슬픔을 주겠다. \\[4pt]
사랑보다 소중한 슬픔을 주겠다. \\[14pt]
겨울 밤 거리에서 귤 몇 개 놓고 \\[4pt]
살아온 추위와 떨고 있는 할머니에게 \\[4pt]
귤값을 깎으면서 기뻐하던 너를 위하여 \\[4pt]
나는 슬픔의 평등한 얼굴을 보여 주겠다. \\[14pt]
내가 어둠 속에서 너를 부를 때 \\[4pt]
단 한 번도 평등하게 웃어 주지 않은, \\[4pt]
가마니에 덮인 사람이 얼어 죽을 때 \\[4pt]
가마니 한 장조차 덮어 주지 않은 \\[4pt]
무관심한 너의 사랑을 위해 \\[4pt]
흘릴 줄 모르는 너의 눈물을 위해 \\[14pt]
나는 이제 너에게도 기다림을 주겠다. \\[4pt]
이 세상에 내리던 함박눈을 멈추겠다. \\[4pt]
보리밭에 내리던 봄눈들을 데리고 \\[4pt]
추위에 떠는 사람들의 슬픔에게 다녀와서 \\[4pt]
눈 그친 눈길을 너와 함께 걷겠다. \\[14pt]
슬픔의 힘에 대한 이야기를 하며 \\[4pt]
기다림의 슬픔까지 걸어가겠다.
}


\kfSecHeader{kfAreaLiterature}{문제}{}

\begin{kfQBlock}{1}{}
\kfQStem{이 시에서 말하는 ‘나’와 듣는 ‘너’는 각각 무엇을 의미할까요?}
\begin{kfChoices}
\kfChoice 나 : 겨울, 너 : 봄
\kfChoice 나 : 어른, 너 : 아이
\kfChoice 나 : 슬픔, 너 : 기쁨
\kfChoice 나 : 선생님, 너 : 학생
\kfChoice 나: 할머니, 너 : 귤 장수
\end{kfChoices}
\end{kfQBlock}
\begin{kfQBlock}{2}{}
\kfQStem{이 시에서 ‘너’는 어떤 존재로 그려지고 있나요?}
\begin{kfChoices}
\kfChoice 자신만의 즐거움을 위해 남을 생각하지 않는 이기적인 존재
\kfChoice 다른 사람의 아픔을 위로해주는 따뜻한 존재
\kfChoice 세상 모든 사람을 평등하게 대하는 정의로운 존재
\kfChoice 슬픔에 빠진 사람들을 보면 함께 슬퍼하는 공감적인 존재
\kfChoice 미래에 대한 희망을 주는 긍정적인 존재
\end{kfChoices}
\end{kfQBlock}
\begin{kfQBlock}{3}{}
\kfQStem{“사랑보다 소중한 슬픔”이라는 표현에 담긴 뜻으로 가장 알맞은 것은 무엇일까요?}
\begin{kfChoices}
\kfChoice 사랑에 실패해서 느끼는 깊은 슬픔이 가장 아름답고 소중하다는 뜻
\kfChoice 자기만 아는 사랑보다, 남의 아픔에 공감하는 슬픔이 더 가치 있다는 뜻
\kfChoice 진정한 사랑을 하려면 먼저 많은 슬픔을 충분히 겪어 보아야 한다는 뜻
\kfChoice 사랑은 결국 슬픔으로 끝날 운명이기 때문에 슬픔이 더 중요하다는 뜻
\kfChoice 슬픔이라는 감정이 기쁨이라는 감정보다 사람을 더 강렬하게 만든다는 뜻
\end{kfChoices}
\end{kfQBlock}
\begin{kfQBlock}{4}{}
\kfQStem{이 시에서 ‘함박눈’과 ‘봄눈’에 대한 설명으로 알맞은 것은 무엇일까요?}
\begin{kfChoices}
\kfChoice '함박눈'은 이기적인 기쁨을, '봄눈'은 따뜻한 위로를 상징한다.
\kfChoice 둘 다 아름다운 자연을 나타내는 긍정적인 의미이다.
\kfChoice 둘 다 가난한 사람들을 더 힘들게 만드는 부정적인 의미이다.
\kfChoice '함박눈'은 겨울을, '봄눈'은 봄을 나타내는 단순한 계절어이다.
\kfChoice '함박눈'은 슬픔을, '봄눈'은 기쁨을 상징한다.
\end{kfChoices}
\end{kfQBlock}
\begin{kfQBlock}{5}{}
\kfQStem{말하는 이가 ‘너’에게 “슬픔을 주겠다”고 말하는 진짜 이유는 무엇일까요?}
\begin{kfChoices}
\kfChoice ‘너’도 다른 사람의 아픔에 공감할 줄 아는 마음을 갖게 하려고
\kfChoice ‘너’에게 슬픔이 얼마나 춥고 외롭게 만드는 무서운 감정인지 뼈저리게 알려주기 위해서
\kfChoice ‘너’가 나보다 훨씬 더 행복하고 기뻐하는 모습을 보니 참을 수 없을 만큼 질투가 나서
\kfChoice ‘너’ 또한 나처럼 깊은 슬픔에 잠겨 헤어나오지 못하는 불행한 존재로 만들어 버리고 싶어서
\kfChoice ‘너’가 과거에 가난한 이웃들에게 저지른 씻을 수 없는 잘못들에 대해 엄중한 벌을 내리기 위해서
\end{kfChoices}
\end{kfQBlock}
\begin{kfQBlock}{6}{}
\kfQStem{말하는 이가 마지막에 꿈꾸는 세상의 모습으로 가장 알맞은 것은 무엇일까요?}
\begin{kfChoices}
\kfChoice 슬픔과 기쁨이 함께 걸으며 서로를 이해하는 세상
\kfChoice 기쁨이 슬픔에게 잘못을 빌고 영원히 모두 복종하는 세상
\kfChoice 슬픔이 기쁨을 이기고 세상 모든 것을 다스리는 세상
\kfChoice 슬픔과 기쁨이 서로 관계없이 각자 따로 살아가는 세상
\kfChoice 오직 슬픔만이 가득하여 모든 것이 꽁꽁 얼어붙은 세상
\end{kfChoices}
\end{kfQBlock}
\begin{kfQBlock}{7}{}
\kfQStem{말하는 이가 이야기하려는 ‘슬픔의 힘’이란 무엇일까요?}
\begin{kfChoices}
\kfChoice 상대방을 굵은 눈물로 끝내 굴복시키는 힘
\kfChoice 다른 사람의 아픔을 이해하고 함께 나눌 수 있는 힘
\kfChoice 모든 것을 한꺼번에 파괴할 수 있는 무서운 힘
\kfChoice 어떤 어려움도 혼자서 묵묵히 견뎌낼 수 있는 힘
\kfChoice 온 세상을 온통 슬프게 만들어 버릴 수 있는 힘
\end{kfChoices}
\end{kfQBlock}
\begin{kfQBlock}{8}{}
\kfQStem{이 시에 대한 설명으로 알맞지 \uline{않은} 것은 무엇일까요?}
\begin{kfChoices}
\kfChoice '슬픔'과 '기쁨'을 사람처럼 표현하고 있다.
\kfChoice 결국 기쁨을 추구하는 삶을 버려야 한다고 주장하고 있다.
\kfChoice ‘-겠다’라는 말을 반복하여 의지적인 태도를 보여준다.
\kfChoice 이기적인 삶의 태도에 대해 비판적으로 이야기하고 있다.
\kfChoice 어려운 이웃에 대한 따뜻한 관심과 사랑을 촉구하고 있다.
\end{kfChoices}
\end{kfQBlock}
\begin{kfQBlock}{1}{}
\kfQStem{이 시에서 말하는 ‘슬픔’과 ‘기쁨’은 우리가 보통 생각하는 것과 다릅니다. 두 감정이 어떻게 다르게 표현되었는지 서술하세요.}
\kfAnswerNote{답안}{4}
\end{kfQBlock}
\begin{kfQBlock}{2}{}
\kfQStem{‘나’가 ‘너’를 이기적이라고 생각하게 된 구체적인 행동 두 가지를 서술하세요.}
\kfAnswerNote{답안}{4}
\end{kfQBlock}

\kfSecHeader{kfAreaLiterature}{활동}{문장\_독해}
\noindent\textit{\small 다음 문장을 읽고 물음에 답하시오.}\par\medskip
\begin{kfSentenceUnit}{1}{}
\kfSentQ{1.}{다음 문장에서 서술어를 기준으로 각각의 주어는 무엇인가요?(1) 서술어 '놓고'의 주어는 무엇인가요?}
\begin{kfChoices}
\kfChoice 귤이
\kfChoice 내가
\kfChoice 네가
\kfChoice 할머니가
\end{kfChoices}
\kfSentQ{1.}{(2) 서술어 '떨고 있는'의 주어는 무엇인가요?}
\begin{kfChoices}
\kfChoice 내가
\kfChoice 네가
\kfChoice 추위가
\kfChoice 할머니가
\end{kfChoices}
\kfSentQ{1.}{(3) 서술어 '깎으면서'의 주어는 무엇인가요?}
\begin{kfChoices}
\kfChoice 귤이
\kfChoice 내가
\kfChoice 네가
\kfChoice 할머니가
\end{kfChoices}
\kfSentQ{1.}{(4) 서술어 '기뻐하던'의 주어는 무엇인가요?}
\begin{kfChoices}
\kfChoice 내가
\kfChoice 네가
\kfChoice 추위가
\kfChoice 할머니가
\end{kfChoices}
\kfSentQ{1.}{(5) 서술어 '보여 주겠다'의 주어는 무엇인가요?}
\begin{kfChoices}
\kfChoice 내가
\kfChoice 네가
\kfChoice 슬픔이
\kfChoice 할머니가
\end{kfChoices}
\end{kfSentenceUnit}

\begin{kfSentenceUnit}{2}{}
\kfSentQ{2.}{다음 문장에서 서술어를 기준으로 각각의 주어는 무엇인가요?(1) 서술어 '주겠다'의 주어는 무엇인가요?}
\begin{kfChoices}
\kfChoice 내가
\kfChoice 네가
\kfChoice 세상이
\kfChoice 기다림이
\end{kfChoices}
\kfSentQ{2.}{(2) 서술어 '내리던'의 주어는 무엇인가요?}
\begin{kfChoices}
\kfChoice 내가
\kfChoice 네가
\kfChoice 세상이
\kfChoice 함박눈이
\end{kfChoices}
\kfSentQ{2.}{(3) 서술어 '멈추겠다'의 주어는 무엇인가요?}
\begin{kfChoices}
\kfChoice 내가
\kfChoice 네가
\kfChoice 기다림이
\kfChoice 함박눈이
\end{kfChoices}
\end{kfSentenceUnit}


\kfStartNonlit{비문학 영역 — [인문] 사람은 본래 착할까, 나쁠까?}


\kfSecHeader{kfAreaNonlit}{지문}{[인문] 사람은 본래 착할까, 나쁠까?}
\kfPassageNote{[인문] 사람은 본래 착할까, 나쁠까?}{%
중국의 전국시대는 여러 나라가 서로 싸우는 혼란한 시대였다. 왕들은 더 많은 땅을 차지하려고 끊임없이 전쟁을 벌였다. 이 때문에 백성들은 매우 힘들고 고통스러운 삶을 살아야 했다. 이런 상황에서 많은 학자들이 나타나 세상을 평화롭게 만들고 백성들의 고통을 덜어주는 방법을 찾으려고 노력했다. 그 과정에서 '사람의 본성은 무엇인가'라는 문제에 대해 깊이 생각하게 되었다.

사람의 본성에 대한 생각은 크게 세 가지로 나뉜다. 첫 번째는 맹자가 주장한 성선설이다. 맹자는 사람이 태어날 때부터 착한 마음을 가지고 있다고 믿었다. 예를 들어, 우물에 빠진 아이를 보면 누구나 저절로 불쌍한 마음이 든다는 것이다. 이런 마음은 누가 가르쳐 주지 않아도 자연스럽게 나오는 것이라고 했다. 맹자는 사람이 스스로 노력하면 이런 착한 본성을 키울 수 있다고 생각했다.

두 번째는 순자가 말한 성악설이다. 순자는 사람이 본래 나쁜 마음을 가지고 있다고 주장했다. 사람은 자기가 원하는 것을 얻으려는 욕심이 많지만, 모든 사람이 원하는 것을 다 가질 수는 없다. 그래서 사람들끼리 서로 싸우게 된다는 것이다. 이런 문제를 해결하려면 나라의 법과 규칙이 필요하다고 했다.

세 번째는 고자의 성무선악설이다. 고자는 사람이 본래 착하지도 나쁘지도 않다고 생각했다. 사람의 본성은 배고플 때 음식을 먹고 싶어하는 것과 같은 자연스러운 욕구일 뿐이라고 했다. 이것을 착하다거나 나쁘다고 말할 수는 없다는 것이다.

이런 생각들은 단순히 학자들의 깊은 논의에서 끝나지 않았다. 힘을 가진 사람들은 자신에게 유리한 이론을 선택해서 정치에 이용했다. 맹자의 성선설은 권력자들이 '우리는 본래 착하니까 간섭하지 말라'는 주장을 펼 때 사용되었다. 반면 순자의 성악설은 임금이 백성들을 강하게 다스려야 한다는 이유로 쓰였다. 결국 사람의 본성에 대한 생각은 정치와 깊은 관련이 있었던 것이다.
}


\kfSecHeader{kfAreaNonlit}{문제}{}

\begin{kfQBlock}{1}{}
\kfQStem{윗글의 중심 내용으로 가장 적절한 것은 무엇인가요?}
\begin{kfChoices}
\kfChoice 인간 본성에 대한 세 가지 관점과 정치적 이용
\kfChoice 고대 중국 학자들의 생애와 업적
\kfChoice 힘을 가진 사람과 백성들의 관계 변화
\kfChoice 인간의 본성에 대한 동양과 서양의 관점 비교
\kfChoice 전국 시대의 역사적 특징
\end{kfChoices}
\end{kfQBlock}
\begin{kfQBlock}{2}{}
\kfQStem{윗글의 내용과 일치하지 않는 것은 무엇인가요?}
\begin{kfChoices}
\kfChoice 순자는 인간의 욕심을 해결하기 위해 개인의 자율적인 노력을 강조했다.
\kfChoice 맹자는 인간이 노력하면 착한 본성을 키울 수 있다고 보았다.
\kfChoice 고자는 인간의 본성이 선과 악으로 판단할 수 없는 것이라 생각했다.
\kfChoice 인간 본성에 대한 이론은 통치자들의 정치적 목적으로 사용되기도 했다.
\kfChoice 전국 시대는 전쟁으로 혼란스러운 시기였다.
\end{kfChoices}
\end{kfQBlock}
\begin{kfQBlock}{3}{}
\kfQStem{<보기>의 관점과 가장 가까운 생각을 가진 학자는 누구인가요?}
\begin{kfEvidence}인간은 이기적인 존재이지만, 꾸준한 교육과 사회 제도를 통해 도덕적인 사람으로 거듭날 수 있다. 타고난 그대로 두면 사회는 혼란에 빠질 것이다.\end{kfEvidence}
\begin{kfChoices}
\kfChoice 맹자
\kfChoice 순자
\kfChoice 고자
\kfChoice 공자
\kfChoice 노자
\end{kfChoices}
\end{kfQBlock}
\begin{kfQBlock}{4}{}
\kfQStem{다음 <보기>의 상황을 지문의 관점으로 분석할 때, 적절하지 \uline{않은} 것은 무엇인가요?}
\begin{kfEvidence}한 마을에서 공동 우물의 물이 부족해지자 사람들이 서로 더 많이 물을 가져가려고 다투기 시작했다.\end{kfEvidence}
\begin{kfChoices}
\kfChoice 맹자의 관점: 다툼이 일어나는 것은 사람의 본성이 악하기 때문이다.
\kfChoice 순자의 관점: 물 사용에 대한 명확한 규칙을 만들어 강제해야 한다.
\kfChoice 순자의 관점: 사람들의 욕심이 충돌할 때 갈등이 발생하는 것은 당연하다.
\kfChoice 고자의 관점: 물을 원하는 것은 자연스러운 욕구이므로 선악으로 판단할 일이 아니다.
\kfChoice 맹자의 관점: 사람들이 서로 양보하는 마음을 기를 수 있다면 문제가 해결될 것이다.
\end{kfChoices}
\end{kfQBlock}
\begin{kfQBlock}{5}{}
\kfQStem{㉠\textasciitilde{}㉢에 대한 설명으로 적절하지 \uline{않은} 것은 무엇인가요?}
\begin{kfEvidence}㉠ 성선설 ㉡ 성악설 ㉢ 성무선악설\end{kfEvidence}
\begin{kfChoices}
\kfChoice ㉡은 인간 사이의 다툼이 발생하는 원인을 본성에서 찾았다.
\kfChoice ㉡과 ㉢은 인간의 욕구를 긍정적으로 평가한다는 점에서 공통된다.
\kfChoice ㉠은 인간의 자연스러운 감정에서 선한 본성의 근거를 찾았다.
\kfChoice ㉢은 ㉠, ㉡과 달리 본성을 도덕적 판단의 대상에서 제외했다.
\kfChoice ㉠과 ㉡은 모두 인간의 본성을 한 가지 방향으로 규정하고 있다.
\end{kfChoices}
\end{kfQBlock}
\begin{kfQBlock}{6}{}
\kfQStem{다음 <보기>의 주장을 한 사람의 생각과 가장 가까운 이론은 무엇인가요?}
\begin{kfEvidence}"사람들은 맛있는 것을 먹고 싶어 하고, 따뜻한 옷을 입고 싶어 하는 욕심이 있습니다. 만약 이 욕심을 그대로 둔다면 사람들은 서로 좋은 것을 차지하려고 다툴 것입니다. 그러니 반드시 규칙을 정해서 잘못된 행동을 바로잡아야 합니다."\end{kfEvidence}
\begin{kfChoices}
\kfChoice 성선설
\kfChoice 성악설
\kfChoice 성무선악설
\kfChoice 자연주의 사상
\kfChoice 민주주의의 원리
\end{kfChoices}
\end{kfQBlock}
\begin{kfQBlock}{7}{}
\kfQStem{다음 <보기>의 대화에서 ㉠에 들어갈 알맞은 말은 무엇인가요?}
\begin{kfEvidence}철수: 배고플 때 밥을 먹는 건 착한 걸까, 나쁜 걸까?

영희: 그건 착한 것도 아니고 나쁜 것도 아니야. 그냥 배가 고파서 먹는 거잖아.

철수: 맞아. 고자라는 학자도 사람의 본성은 [    ㉠    ]과 같다고 했어.\end{kfEvidence}
\begin{kfChoices}
\kfChoice 악한 마음
\kfChoice 착한 본성
\kfChoice 자연스러운 욕구
\kfChoice 남을 돕는 행동
\kfChoice 규칙을 지키는 태도
\end{kfChoices}
\end{kfQBlock}
\begin{kfQBlock}{1}{}
\kfQStem{글에서 설명하는 여러 나라가 서로 싸우던 어지러운 시대를 무엇이라고 하나요?}
\kfAnswerLines{1}
\end{kfQBlock}
\begin{kfQBlock}{2}{}
\kfQStem{전국 시대의 학자들이 사람의 본성에 대해 깊이 생각하게 된 근본적인 이유는 무엇인가요?}
\kfAnswerLines{1}
\end{kfQBlock}
\begin{kfQBlock}{3}{}
\kfQStem{순자는 사람들끼리 서로 싸우게 되는 근본적인 원인이 무엇이라고 보았나요?}
\kfAnswerLines{1}
\end{kfQBlock}
\begin{kfQBlock}{4}{}
\kfQStem{힘을 가진 사람들이 '우리는 원래 착하니 간섭하지 말라'는 주장을 하기 위해 이용했던 이론은 무엇인가요?}
\kfAnswerLines{1}
\end{kfQBlock}
\begin{kfQBlock}{5}{}
\kfQStem{임금이 백성을 강하게 다스리는 통치의 근거로 사용된 이론은 무엇인가요?}
\kfAnswerLines{1}
\end{kfQBlock}
\begin{kfQBlock}{1}{}
\kfQStem{순자의 주장이 정치적으로 어떻게 이용되었는지, 그의 핵심 주장의 논리적 흐름을 바탕으로 서술하세요.}
\kfAnswerNote{답안}{4}
\end{kfQBlock}

\kfSecHeader{kfAreaNonlit}{활동}{문장\_독해}
\noindent\textit{\small 다음 문장을 읽고 물음에 답하시오.}\par\medskip
\begin{kfSentenceUnit}{1}{}
\kfSentQ{1.}{전국시대는 어떤 시대였나요?}
\begin{kfChoices}
\kfChoice 여러 나라가 서로 싸우는 혼란한 시대였다.
\kfChoice 백성들이 풍요롭게 살며 걱정이 없는 시대였다.
\kfChoice 모든 나라가 하나로 통일되어 안정된 시대였다.
\kfChoice 왕들이 서로 화합하여 평화롭게 지내는 시대였다.
\end{kfChoices}
\end{kfSentenceUnit}

\begin{kfSentenceUnit}{2}{}
\kfSentQ{2.}{왕들이 끊임없이 전쟁을 벌인 이유는 무엇인가요?}
\begin{kfChoices}
\kfChoice 더 많은 땅을 차지하기 위해서.
\kfChoice 백성들의 고통을 줄여주기 위해서.
\kfChoice 이웃 나라와 친하게 지내기 위해서.
\kfChoice 자신의 왕위를 다른 사람에게 물려주기 위해서.
\end{kfChoices}
\end{kfSentenceUnit}

\begin{kfSentenceUnit}{3}{}
\kfSentQ{3.}{'이 때문에'에서 '이'가 가리키는 것은 무엇인가요?}
\begin{kfChoices}
\kfChoice 이웃 나라와 서로 돕고 지내는 것.
\kfChoice 학자들이 새로운 이론을 만들어내는 것.
\kfChoice 백성들이 평화롭게 농사를 짓고 사는 것.
\kfChoice 왕들이 더 많은 땅을 차지하려고 끊임없이 전쟁을 벌인 것.
\end{kfChoices}
\end{kfSentenceUnit}

\begin{kfSentenceUnit}{4}{}
\kfSentQ{4.}{백성들이 고통받은 이유는 무엇인가요?}
\begin{kfChoices}
\kfChoice 왕들이 끊임없이 전쟁을 벌였기 때문이다.
\kfChoice 나라가 통일되어 할 일이 없어졌기 때문이다.
\kfChoice 왕들이 백성들을 위해 정치를 잘했기 때문이다.
\kfChoice 날씨가 좋지 않아 농사가 잘되지 않았기 때문이다.
\end{kfChoices}
\end{kfSentenceUnit}

\begin{kfSentenceUnit}{5}{}
\kfSentQ{5.}{'이런 상황'은 어떤 상황을 가리키나요?}
\begin{kfChoices}
\kfChoice 학자들이 모두 사라져 버린 상황.
\kfChoice 백성들이 풍요롭고 행복하게 사는 상황.
\kfChoice 왕들이 서로 화해하고 전쟁을 멈춘 상황.
\kfChoice 전쟁 때문에 백성들이 힘들고 고통스럽게 사는 상황.
\end{kfChoices}
\end{kfSentenceUnit}

\begin{kfSentenceUnit}{6}{}
\kfSentQ{6.}{전쟁 때문에 백성들이 힘들고 고통스럽게 사는 상황에서 학자들이 찾으려고 노력한 것은 무엇인가요?}
\begin{kfChoices}
\kfChoice 자신의 이름을 널리 알리고 출세하는 방법.
\kfChoice 왕에게 잘 보여서 높은 벼슬을 얻는 방법.
\kfChoice 더 넓은 땅을 차지하기 위한 효율적인 전쟁 방법.
\kfChoice 세상을 평화롭게 만들고 백성들의 고통을 덜어주는 방법.
\end{kfChoices}
\end{kfSentenceUnit}

\begin{kfSentenceUnit}{7}{}
\kfSentQ{7.}{'그 과정'이 가리키는 것은 무엇인가요?}
\begin{kfChoices}
\kfChoice 우물에 빠질 뻔한 아이를 구해주며 사람의 본성을 증명하는 과정.
\kfChoice 전쟁의 고통을 견디지 못하고 백성들이 다른 나라로 도망치는 과정.
\kfChoice 왕들이 더 넓은 영토를 차지하기 위해 끊임없이 전쟁을 벌이는 과정.
\kfChoice 학자들이 세상을 평화롭게 만들고 백성들의 고통을 덜어주는 방법을 찾으려고 노력하는 과정.
\end{kfChoices}
\end{kfSentenceUnit}

\begin{kfSentenceUnit}{8}{}
\kfSentQ{8.}{사상가들이 깊이 생각하게 된 문제는 무엇인가요?}
\begin{kfChoices}
\kfChoice 사람의 본성은 무엇인가.
\kfChoice 어떤 왕이 가장 힘이 센가.
\kfChoice 세금을 얼마나 더 많이 걷어야 하는가.
\kfChoice 어떻게 하면 전쟁에서 승리할 수 있는가.
\end{kfChoices}
\end{kfSentenceUnit}

\begin{kfSentenceUnit}{9}{}
\kfSentQ{9.}{사람의 본성에 대한 생각은 몇 가지로 나뉘나요?}
\begin{kfChoices}
\kfChoice 두 가지.
\kfChoice 네 가지.
\kfChoice 세 가지.
\kfChoice 한 가지.
\end{kfChoices}
\end{kfSentenceUnit}

\begin{kfSentenceUnit}{10}{}
\kfSentQ{10.}{이 문장을 통해 앞으로 어떤 내용이 나올 것을 예상할 수 있나요?}
\begin{kfChoices}
\kfChoice 고자의 성무선악설에 대한 비판.
\kfChoice 전국시대의 전쟁 역사에 대한 설명.
\kfChoice 각 나라의 지리적 특징에 대한 설명.
\kfChoice 사람의 본성에 대한 세 가지 생각들에 대한 설명.
\end{kfChoices}
\end{kfSentenceUnit}

\begin{kfSentenceUnit}{11}{}
\kfSentQ{11.}{성선설을 주장한 사람은 누구인가요?}
\begin{kfChoices}
\kfChoice 맹자.
\kfChoice 순자.
\kfChoice 고자.
\kfChoice 공자.
\end{kfChoices}
\end{kfSentenceUnit}

\begin{kfSentenceUnit}{12}{}
\kfSentQ{12.}{맹자는 사람이 무엇을 가지고 있다고 믿었나요?}
\begin{kfChoices}
\kfChoice 욕심.
\kfChoice 착한 마음.
\kfChoice 나쁜 마음.
\kfChoice 아무런 마음도 없음.
\end{kfChoices}
\end{kfSentenceUnit}

\begin{kfSentenceUnit}{13}{}
\kfSentQ{13.}{맹자는 사람이 언제부터 착한 마음을 가지고 있다고 믿었나요?}
\begin{kfChoices}
\kfChoice 태어날 때부터.
\kfChoice 어른이 되어서.
\kfChoice 교육을 받은 후부터.
\kfChoice 나쁜 친구를 사귄 후부터.
\end{kfChoices}
\end{kfSentenceUnit}

\begin{kfSentenceUnit}{14}{}
\kfSentQ{14.}{'예를 들어'는 무엇을 설명하기 위한 예시인가요?}
\begin{kfChoices}
\kfChoice 사람의 본성은 교육으로만 만들어진다는 것.
\kfChoice 사람은 태어날 때 아무런 마음도 없다는 것.
\kfChoice 사람이 태어날 때부터 나쁜 마음을 가지고 있다는 것.
\kfChoice 사람이 태어날 때부터 착한 마음을 가지고 있다는 것.
\end{kfChoices}
\end{kfSentenceUnit}

\begin{kfSentenceUnit}{15}{}
\kfSentQ{15.}{맹자는 우물에 빠진 아이를 볼 때, 누구나 어떤 마음이 든다고 했나요?}
\begin{kfChoices}
\kfChoice 불쌍한 마음.
\kfChoice 미워하는 마음.
\kfChoice 질투하는 마음.
\kfChoice 기뻐하는 마음.
\end{kfChoices}
\end{kfSentenceUnit}

\begin{kfSentenceUnit}{16}{}
\kfSentQ{16.}{'이런 마음'이 가리키는 것은 무엇인가요?}
\begin{kfChoices}
\kfChoice 우물에 빠진 아이를 보고 비웃는 마음.
\kfChoice 우물에 빠진 아이를 보고 화내는 마음.
\kfChoice 우물에 빠진 아이를 모른 척하는 마음.
\kfChoice 우물에 빠진 아이를 보면 드는 불쌍한 마음.
\end{kfChoices}
\end{kfSentenceUnit}

\begin{kfSentenceUnit}{17}{}
\kfSentQ{17.}{맹자에 따르면, 착한 마음은 어떻게 나오나요?}
\begin{kfChoices}
\kfChoice 많은 책을 읽어야 나온다.
\kfChoice 법과 규칙을 잘 지켜야 나온다.
\kfChoice 부모님이나 선생님에게 배워서 나온다.
\kfChoice 누가 가르쳐 주지 않아도 자연스럽게 나온다.
\end{kfChoices}
\end{kfSentenceUnit}

\begin{kfSentenceUnit}{18}{}
\kfSentQ{18.}{'이런 착한 본성'이 가리키는 것은 무엇인가요?}
\begin{kfChoices}
\kfChoice 남을 위해 억지로 꾸며낸 착한 마음.
\kfChoice 태어날 때부터 가지고 있는 착한 마음.
\kfChoice 벌을 받지 않으려고 연기하는 착한 마음.
\kfChoice 교육을 통해 후천적으로 만들어진 착한 마음.
\end{kfChoices}
\end{kfSentenceUnit}

\begin{kfSentenceUnit}{19}{}
\kfSentQ{19.}{맹자에 따르면, 착한 본성을 키우려면 어떻게 해야 하나요?}
\begin{kfChoices}
\kfChoice 사람이 스스로 노력해야 한다.
\kfChoice 엄격한 법과 규칙을 따라야 한다.
\kfChoice 아무런 노력도 하지 않아도 된다.
\kfChoice 다른 사람의 도움을 기다려야 한다.
\end{kfChoices}
\end{kfSentenceUnit}

\begin{kfSentenceUnit}{20}{}
\kfSentQ{20.}{성악설을 말한 사람은 누구인가요?}
\begin{kfChoices}
\kfChoice 순자.
\kfChoice 맹자.
\kfChoice 노자.
\kfChoice 고자.
\end{kfChoices}
\end{kfSentenceUnit}

\begin{kfSentenceUnit}{21}{}
\kfSentQ{21.}{앞서 나온 '성선설'과 '성악설'은 사람의 본성을 보는 관점에서 서로 어떤 관계인가요?}
\begin{kfChoices}
\kfChoice 서로 보완해주는 관계.
\kfChoice 아무런 관련이 없는 관계.
\kfChoice 서로 반대되거나 대립하는 관계.
\kfChoice 서로 반대되거나 대립하는 관계.
\end{kfChoices}
\end{kfSentenceUnit}

\begin{kfSentenceUnit}{22}{}
\kfSentQ{22.}{순자는 사람이 본래 어떤 마음을 가지고 있다고 주장했나요?}
\begin{kfChoices}
\kfChoice 착한 마음.
\kfChoice 나쁜 마음.
\kfChoice 불쌍히 여기는 마음.
\kfChoice 착하지도 나쁘지도 않은 마음.
\end{kfChoices}
\end{kfSentenceUnit}

\begin{kfSentenceUnit}{23}{}
\kfSentQ{23.}{사람의 특징은 무엇인가요?}
\begin{kfChoices}
\kfChoice 태어날 때부터 도덕적인 규칙을 알고 있다.
\kfChoice 자기가 원하는 것을 얻으려는 욕심이 많다.
\kfChoice 남을 위해 자신의 것을 양보하는 마음이 크다.
\kfChoice 본래 욕심이 없고 깨끗한 마음을 가지고 있다.
\end{kfChoices}
\end{kfSentenceUnit}

\begin{kfSentenceUnit}{24}{}
\kfSentQ{24.}{현실의 한계는 무엇인가요?}
\begin{kfChoices}
\kfChoice 자연은 모든 것을 무한히 제공해 준다.
\kfChoice 사람의 욕심은 매우 작아서 금방 채워진다.
\kfChoice 모든 사람이 원하는 것을 다 가질 수 있다.
\kfChoice 모든 사람이 원하는 것을 다 가질 수는 없다.
\end{kfChoices}
\end{kfSentenceUnit}

\begin{kfSentenceUnit}{25}{}
\kfSentQ{25.}{'그래서'의 구체적인 내용은 무엇인가요?}
\begin{kfChoices}
\kfChoice 나라의 법과 규칙이 아주 엄격하게 만들어져 있어서.
\kfChoice 사람은 본래 착한 마음을 가지고 태어나서 서로 양보하니까.
\kfChoice 세상에는 물자가 풍부하여 누구나 원하는 것을 가질 수 있어서.
\kfChoice 사람은 자기가 원하는 것을 얻으려는 욕심이 많지만, 모든 사람이 원하는 것을 다 가질 수는 없어서.
\end{kfChoices}
\end{kfSentenceUnit}

\begin{kfSentenceUnit}{26}{}
\kfSentQ{26.}{사람들이 서로 싸우게 되는 이유는 무엇인가요?}
\begin{kfChoices}
\kfChoice 서로 양보하고 배려하기 때문이다.
\kfChoice 법과 규칙이 너무 엄격하기 때문이다.
\kfChoice 원래 싸우는 것을 좋아하기 때문이다.
\kfChoice 모든 사람이 자기가 원하는 것을 다 가질 수 없기 때문이다.
\end{kfChoices}
\end{kfSentenceUnit}

\begin{kfSentenceUnit}{27}{}
\kfSentQ{27.}{'이런 문제'는 무엇을 가리키나요?}
\begin{kfChoices}
\kfChoice 사람들끼리 서로 싸우는 문제.
\kfChoice 학자들이 서로 논쟁하는 문제.
\kfChoice 백성들이 배고픔에 시달리는 문제.
\kfChoice 우물에 빠진 아이를 구하지 않는 문제.
\end{kfChoices}
\end{kfSentenceUnit}

\begin{kfSentenceUnit}{28}{}
\kfSentQ{28.}{순자는 사람들끼리 서로 싸우는 문제를 해결하려면 무엇이 필요하다고 했나요?}
\begin{kfChoices}
\kfChoice 자유로운 생각.
\kfChoice 나라의 법과 규칙.
\kfChoice 왕의 따뜻한 사랑.
\kfChoice 사람들의 착한 본성.
\end{kfChoices}
\end{kfSentenceUnit}

\begin{kfSentenceUnit}{29}{}
\kfSentQ{29.}{성무선악설을 주장한 사람은 누구인가요?}
\begin{kfChoices}
\kfChoice 순자.
\kfChoice 고자.
\kfChoice 맹자.
\kfChoice 공자.
\end{kfChoices}
\end{kfSentenceUnit}

\begin{kfSentenceUnit}{30}{}
\kfSentQ{30.}{고자는 사람의 본성을 어떻게 생각했나요?}
\begin{kfChoices}
\kfChoice 본래 매우 나쁘다.
\kfChoice 본래 매우 착하다.
\kfChoice 본래 정해져 있다.
\kfChoice 본래 착하지도 나쁘지도 않다.
\end{kfChoices}
\end{kfSentenceUnit}

\begin{kfSentenceUnit}{31}{}
\kfSentQ{31.}{고자의 생각은 맹자, 순자와 어떻게 다른가요?}
\begin{kfChoices}
\kfChoice 맹자와 순자는 본성이 없다고 했으나, 고자는 있다고 했다.
\kfChoice 맹자와 순자는 본성이 변한다고 했으나, 고자는 변하지 않는다고 했다.
\kfChoice 맹자와 순자는 본성이 중요하지 않다고 했으나, 고자는 중요하다고 했다.
\kfChoice 맹자는 착하다, 순자는 나쁘다고 했는데, 고자는 착하지도 나쁘지도 않다고 했다.
\end{kfChoices}
\end{kfSentenceUnit}

\begin{kfSentenceUnit}{32}{}
\kfSentQ{32.}{고자는 사람의 본성을 무엇과 같다고 했나요?}
\begin{kfChoices}
\kfChoice 나라에 충성하고 싶은 마음.
\kfChoice 부모님께 효도하고 싶은 마음.
\kfChoice 어려운 사람을 돕고 싶은 마음.
\kfChoice 배고플 때 음식을 먹고 싶어하는 것.
\end{kfChoices}
\end{kfSentenceUnit}

\begin{kfSentenceUnit}{33}{}
\kfSentQ{33.}{고자의 주장에 따르면 '사람의 본성'은 무엇으로 바꾸어 생각할 수 있나요?}
\begin{kfChoices}
\kfChoice 자연스러운 욕구.
\kfChoice 교육을 통해 배우는 후천적인 학습.
\kfChoice 사회를 유지하기 위한 엄격한 규칙.
\kfChoice 부모님으로부터 배우는 도덕적인 판단.
\end{kfChoices}
\end{kfSentenceUnit}

\begin{kfSentenceUnit}{34}{}
\kfSentQ{34.}{'이것'이 가리키는 것은 무엇인가요?}
\begin{kfChoices}
\kfChoice 배고플 때 음식을 먹고 싶어하는 것과 같은 자연스러운 욕구.
\kfChoice 부모님이나 어른들의 가르침을 잘 따르고 예의 바르게 행동하는 태도.
\kfChoice 사회적인 혼란을 막기 위해 나라에서 정한 법과 규칙을 지키려는 의지.
\kfChoice 어려움에 처한 사람을 보았을 때 조건 없이 돕고 싶어하는 도덕적인 마음.
\end{kfChoices}
\end{kfSentenceUnit}

\begin{kfSentenceUnit}{35}{}
\kfSentQ{35.}{고자는 자연스러운 욕구를 어떻게 평가했나요?}
\begin{kfChoices}
\kfChoice 반드시 착하다고 말해야 한다.
\kfChoice 반드시 나쁘다고 말해야 한다.
\kfChoice 착하지만 때로는 나쁠 수도 있다.
\kfChoice 착하다거나 나쁘다고 말할 수 없다.
\end{kfChoices}
\end{kfSentenceUnit}

\begin{kfSentenceUnit}{36}{}
\kfSentQ{36.}{'이런 생각들'이 가리키는 것은 무엇인가요?}
\begin{kfChoices}
\kfChoice 성선설, 성악설, 성무선악설.
\kfChoice 백성들을 배고픔에서 구하기 위한 농업 기술들.
\kfChoice 전쟁을 멈추고 평화를 실현하는 구체적인 정책들.
\kfChoice 나라를 강력하게 만들기 위한 부국강병의 비결들.
\end{kfChoices}
\end{kfSentenceUnit}

\begin{kfSentenceUnit}{37}{}
\kfSentQ{37.}{'단순히 학자들의 깊은 논의에서 끝나지 않았다'는 어떤 의미인가요?}
\begin{kfChoices}
\kfChoice 아무런 쓸모가 없었다는 의미.
\kfChoice 다른 곳에서도 사용되었다는 의미.
\kfChoice 백성들은 전혀 알지 못했다는 의미.
\kfChoice 오직 학자들끼리만 아는 비밀이었다는 의미.
\end{kfChoices}
\end{kfSentenceUnit}

\begin{kfSentenceUnit}{38}{}
\kfSentQ{38.}{누가 이론을 정치에 이용했나요?}
\begin{kfChoices}
\kfChoice 힘을 가진 사람들.
\kfChoice 농사를 짓는 백성들.
\kfChoice 글을 모르는 아이들.
\kfChoice 장사를 하는 상인들.
\end{kfChoices}
\end{kfSentenceUnit}

\begin{kfSentenceUnit}{39}{}
\kfSentQ{39.}{힘을 가진 사람들은 어떤 이론을 선택했나요?}
\begin{kfChoices}
\kfChoice 자신에게 유리한 이론.
\kfChoice 진실에 가장 가까운 이론.
\kfChoice 백성들에게 인기 있는 이론.
\kfChoice 가장 논리적이고 정확한 이론.
\end{kfChoices}
\end{kfSentenceUnit}

\begin{kfSentenceUnit}{40}{}
\kfSentQ{40.}{힘을 가진 사람들은 자신들에게 유리한 이론을 어디에 이용했나요?}
\begin{kfChoices}
\kfChoice 정치.
\kfChoice 장사.
\kfChoice 교육.
\kfChoice 농사.
\end{kfChoices}
\end{kfSentenceUnit}

\begin{kfSentenceUnit}{41}{}
\kfSentQ{41.}{맹자의 성선설은 언제 사용되었나요?}
\begin{kfChoices}
\kfChoice 권력자들이 '우리는 본래 착하니까 간섭하지 말라'는 주장을 펼 때.
\kfChoice 권력자들이 '나라의 살림이 어려우니 세금을 더 내라'고 강요할 때.
\kfChoice 권력자들이 '백성들은 어리석으니 강력한 법으로 다스려야 한다'고 주장할 때.
\kfChoice 권력자들이 '더 넓은 땅이 필요하니 전쟁을 해야 한다'며 백성들을 동원할 때.
\end{kfChoices}
\end{kfSentenceUnit}

\begin{kfSentenceUnit}{42}{}
\kfSentQ{42.}{순자의 성악설은 어떤 목적으로 쓰였나요?}
\begin{kfChoices}
\kfChoice 임금이 백성들에게 봉사해야 한다는 이유.
\kfChoice 임금이 백성들의 자유를 보장해야 한다는 이유.
\kfChoice 임금이 백성들을 강하게 다스려야 한다는 이유.
\kfChoice 임금이 백성들의 세금을 줄여줘야 한다는 이유.
\end{kfChoices}
\end{kfSentenceUnit}

\begin{kfSentenceUnit}{43}{}
\kfSentQ{43.}{성선설과 성악설의 정치적 활용 방식은 어떻게 다른가요?}
\begin{kfChoices}
\kfChoice 성선설은 세금을 줄이는 핑계로 쓰였고, 성악설은 전쟁을 일으키는 구실로만 쓰였다.
\kfChoice 성선설은 간섭하지 말라는 주장에, 성악설은 강하게 다스려야 한다는 주장에 쓰였다.
\kfChoice 성선설은 백성들을 엄격하게 통제하는 데, 성악설은 백성들에게 자유를 주는 데 쓰였다.
\kfChoice 성선설은 학자들의 토론을 위해서만 쓰였고, 성악설은 백성들의 교육을 위해서만 쓰였다.
\end{kfChoices}
\end{kfSentenceUnit}

\begin{kfSentenceUnit}{44}{}
\kfSentQ{44.}{사람의 본성에 대한 생각은 무엇과 깊은 관련이 있었나요?}
\begin{kfChoices}
\kfChoice 문화.
\kfChoice 정치.
\kfChoice 예술.
\kfChoice 경제.
\end{kfChoices}
\end{kfSentenceUnit}



\clearpage

\kfChapterCover{2}{프레게3 (챕터2)}{Chapter 2}{이번 챕터의 학습 내용을 시작합니다.}

\kfStartGrammar{문법 영역 — 분석\_훈련}


\kfSecHeader{kfAreaGrammar}{개념}{분석\_훈련}

 문장 속 품사 찾기 훈련!

 다음 문장을 읽고, 설명에 해당하는 품사를 모두 찾아 써 보세요.

\kfSecHeader{kfAreaGrammar}{개념}{문장 속 단어들의 역할 (품사)}

우리 주변에는 하늘, 매우, 달리다, 예쁘다처럼 수많은 단어가 있다. 이 단어들은 문장 속에서 저마다 맡은 역할이 있는데, 이렇게 성질이 비슷한 단어끼리 나누어 이름을 붙인 것을 ‘품사’라고 한다. 품사는 크게 9가지 종류가 있다.

먼저 문장의 주인공 역할을 하는 단어들이 있다. ‘하늘’이나 ‘사랑’처럼 사람이나 사물의 이름을 나타내는 명사, 명사를 대신해서 ‘나’나 ‘이것’처럼 가리키는 대명사, 그리고 ‘하나’ 또는 ‘첫째’처럼 수량이나 순서를 나타내는 수사가 바로 여기에 속한다.

또한, 이 주인공의 움직임이나 상태를 설명해 주는 단어들도 있다. ‘아이가 달린다’처럼 움직임을 나타내는 말은 동사라고 하고, ‘하늘이 파랗다’처럼 상태나 성질을 설명하는 말은 형용사라고 한다.

다른 말을 꾸며주는 역할을 하는 단어들도 있다. ‘새 책’의 ‘새’처럼 사람이나 사물의 이름을 꾸며주는 관형사가 있고, ‘아주 높다’의 ‘아주’나 ‘멀리 날다’의 ‘멀리’처럼 움직임이나 상태를 설명하는 말을 꾸며주는 부사가 있다.

단어들 사이의 관계를 풀처럼 붙어서 알려주는 단어도 있다. 바로 조사다. 조사는 ‘내가’의 ‘가’, ‘집으로’의 ‘으로’처럼 다른 말 뒤에 붙어서 그 말이 문장 속에서 어떤 역할을 하는지 알려준다.

마지막으로, 문장 속 다른 단어들과 큰 상관없이 혼자서 감정을 표현하는 단어도 있다. ‘우와’, ‘아차’, ‘응’처럼 놀람이나 느낌, 대답을 나타내는 말이 바로 감탄사이다.
\par\medskip\begin{itemize}[leftmargin=18pt, itemsep=2pt, topsep=2pt, label=\textbullet]
\item 체언(명·대·수) — '학교(명)·그것(대)·둘째(수)'
\item 용언(동·형) — '달린다(동)·아름답다(형)'
\item 수식언·관계언·독립언 — '새(관형)·매우(부)·이/가(조사)·아!(감탄)'
\end{itemize}

\kfSecHeader{kfAreaGrammar}{문제}{}

\begin{kfQBlock}{1}{}
\kfQStem{다음 중 품사의 종류가 다른 하나는 무엇인가요?}
\begin{kfChoices}
\kfChoice 하늘
\kfChoice 사랑
\kfChoice 컴퓨터
\kfChoice 우리
\kfChoice 부산
\end{kfChoices}
\end{kfQBlock}
\begin{kfQBlock}{2}{}
\kfQStem{다음 중 움직임을 나타내는 동사는 무엇인가요?}
\begin{kfChoices}
\kfChoice 깨끗하다
\kfChoice 슬프다
\kfChoice 읽는다
\kfChoice 높다
\kfChoice 많다
\end{kfChoices}
\end{kfQBlock}
\begin{kfQBlock}{3}{}
\kfQStem{<보기>의 설명에 해당하는 품사들로만 짝지어진 것은 무엇인가요?}
\begin{kfEvidence}문장에서 주로 다른 말을 꾸며주는 역할을 한다.\end{kfEvidence}
\begin{kfChoices}
\kfChoice 명사, 대명사
\kfChoice 동사, 형용사
\kfChoice 관형사, 부사
\kfChoice 조사, 감탄사
\kfChoice 수사, 명사
\end{kfChoices}
\end{kfQBlock}
\begin{kfQBlock}{4}{}
\kfQStem{다음 중 조사가 사용되지 \uline{않은} 문장은 무엇인가요?}
\begin{kfChoices}
\kfChoice 하늘이 푸르다.
\kfChoice 매우 빠르다.
\kfChoice 아기가 웃는다.
\kfChoice 저것은 나의 책이다.
\kfChoice 나는 학생이다.
\end{kfChoices}
\end{kfQBlock}
\begin{kfQBlock}{5}{}
\kfQStem{다음 중 홀로 쓰일 수 \uline{없는} 품사는 무엇인가요?}
\begin{kfChoices}
\kfChoice 명사
\kfChoice 동사
\kfChoice 부사
\kfChoice 조사
\kfChoice 감탄사
\end{kfChoices}
\end{kfQBlock}
\begin{kfQBlock}{6}{}
\kfQStem{다음 중 품사의 연결이 \uline{틀린} 것은 무엇인가요?}
\begin{kfChoices}
\kfChoice 다섯 - 수사
\kfChoice 헌 - 관형사
\kfChoice 셋 - 조사
\kfChoice 아차 - 감탄사
\kfChoice 뛴다 - 동사
\end{kfChoices}
\end{kfQBlock}
\begin{kfQBlock}{7}{}
\kfQStem{다음 중 상태나 성질을 나타내는 형용사는 무엇인가요?}
\begin{kfChoices}
\kfChoice 간다
\kfChoice 웃는다
\kfChoice 덥다
\kfChoice 본다
\kfChoice 쓴다
\end{kfChoices}
\end{kfQBlock}
\begin{kfQBlock}{1}{}
\kfQStem{다음 문장에서 밑줄 친 단어들의 품사를 각각 쓰세요.}
\begin{kfEvidence}와, 저 새는 빨리 난다!\end{kfEvidence}
\kfAnswerNote{답안}{4}
\end{kfQBlock}

\kfSecHeader{kfAreaGrammar}{활동}{분석\_훈련}
\noindent\textit{\small 다음 활용 사례를 분석하여 빈칸을 채워 보세요.}\par\medskip
\begin{enumerate}[leftmargin=22pt, itemsep=6pt, topsep=4pt, label=\textbf{\arabic*.}]
\item 다음 문장에서 사람이나 사물의 이름을 나타내는 명사를 모두 찾아 써 보세요.
\begin{kfEvidence}하늘에 구름과 해가 떠 있다.\end{kfEvidence}
\item 다음 문장에서 명사를 대신하여 쓰인 대명사를 모두 찾아 써 보세요.
\begin{kfEvidence}이것은 바로 우리가 찾던 책이다.\end{kfEvidence}
\item 다음 문장에서 수량이나 순서를 나타내는 수사를 모두 찾아 써 보세요.
\begin{kfEvidence}엄마는 배 하나를 사셨다.\end{kfEvidence}
\item 다음 문장에서 움직임을 나타내는 동사를 찾아 기본형으로 써 보세요.
\begin{kfEvidence}아기가 방에서 잠을 잔다.\end{kfEvidence}
\item 다음 문장에서 상태나 성질을 나타내는 형용사를 찾아 기본형으로 써 보세요.
\begin{kfEvidence}가을 하늘이 참 푸르고 높다.\end{kfEvidence}
\item 다음 문장에서 명사를 꾸며주는 관형사를 모두 찾아 써 보세요.
\begin{kfEvidence}저 책상 위에 새 연필이 있다.\end{kfEvidence}
\item 다음 문장에서 움직임이나 상태를 꾸며주는 부사를 모두 찾아 써 보세요.
\begin{kfEvidence}자동차가 빨리 달린다.\end{kfEvidence}
\item 다음 문장에서 다른 말에 붙어 관계를 나타내는 조사를 모두 찾아 써 보세요.
\begin{kfEvidence}아빠가 동생에게 선물을 주셨다.\end{kfEvidence}
\item 다음 문장에서 놀람이나 느낌을 나타내는 감탄사를 찾아 써 보세요.
\begin{kfEvidence}와, 정말 아름다운 풍경이다!\end{kfEvidence}
\item 다음 문장에서 사람이나 사물의 이름을 나타내는 명사를 모두 찾아 써 보세요.
\begin{kfEvidence}나의 가방에는 연필과 공책이 있다.\end{kfEvidence}
\end{enumerate}

\kfSecHeader{kfAreaGrammar}{문제}{}

\begin{kfQBlock}{1}{}
\kfQStem{성질이 비슷한 단어끼리 나누어 이름을 붙인 것을 무엇이라고 하나요?}
\kfAnswerLines{1}
\end{kfQBlock}
\begin{kfQBlock}{2}{}
\kfQStem{‘하늘, 우정, 대한민국’처럼 사람이나 사물의 이름을 나타내는 품사는 무엇인가요?}
\kfAnswerLines{1}
\end{kfQBlock}
\begin{kfQBlock}{3}{}
\kfQStem{‘나, 너, 우리, 여기, 저기’처럼 명사를 대신하여 가리키는 품사는 무엇인가요?}
\kfAnswerLines{1}
\end{kfQBlock}
\begin{kfQBlock}{4}{}
\kfQStem{‘하나, 둘, 첫째, 둘째’처럼 수량이나 순서를 나타내는 품사는 무엇인가요?}
\kfAnswerLines{1}
\end{kfQBlock}
\begin{kfQBlock}{5}{}
\kfQStem{‘달리다, 먹다, 공부하다’처럼 움직임을 나타내는 품사는 무엇인가요?}
\kfAnswerLines{1}
\end{kfQBlock}
\begin{kfQBlock}{6}{}
\kfQStem{‘예쁘다, 착하다, 시원하다’처럼 상태나 성질을 나타내는 품사는 무엇인가요?}
\kfAnswerLines{1}
\end{kfQBlock}
\begin{kfQBlock}{7}{}
\kfQStem{‘새, 헌, 이, 그, 저’처럼 명사를 꾸며주는 품사는 무엇인가요?}
\kfAnswerLines{1}
\end{kfQBlock}
\begin{kfQBlock}{8}{}
\kfQStem{‘매우, 빨리, 잘’처럼 동사나 형용사 등을 꾸며주는 품사는 무엇인가요?}
\kfAnswerLines{1}
\end{kfQBlock}
\begin{kfQBlock}{9}{}
\kfQStem{‘이/가, 을/를, 은/는’처럼 다른 말에 붙어 다른 단어와의 관계를 나타내는 품사는 무엇인가요?}
\kfAnswerLines{1}
\end{kfQBlock}
\begin{kfQBlock}{10}{}
\kfQStem{‘와, 네, 아니요, 여보세요’처럼 독립적으로 쓰이는 품사는 무엇인가요?}
\kfAnswerLines{1}
\end{kfQBlock}


\kfStartLit{문학 영역 — 작품}


\kfSecHeader{kfAreaLiterature}{지문}{작품}
\kfPassageNote{작품}{%
우리 집에는 낡고 오래된 행랑채가 있습니다. 그런데 그중 세 칸이 곧 무너질 것 같아서, 어쩔 수 없이 모두 고치게 되었습니다.

이 일이 있기 전에, 그 세 칸 중 두 칸은 오래전부터 비가 새고 있었습니다. 나는 그것을 알면서도 그냥 놔두다가 제때 고치지 못했습니다. 나머지 한 칸은 한 번만 비가 샜을 때 급히 기와를 교체하게 했습니다.

그런데 이번에 고쳐보니, 오래 비가 샌 곳은 서까래와 추녀, 기둥과 들보가 모두 썩어서 쓸 수 없게 되었습니다. 그래서 돈이 많이 들었습니다. 하지만 한 번만 비가 샌 곳은 나무들이 모두 멀쩡해서 다시 쓸 수 있었기 때문에 돈을 아낄 수 있었습니다.

그래서 나는 이런 생각이 들었습니다.

이런 일은 사람의 경우에도 똑같지 않을까요? 잘못을 알고서도 바로 고치지 않는다면, 오래 비를 맞은 나무가 썩어서 못 쓰게 되듯이 자신을 망치게 될 것입니다. 반대로 잘못한 일을 거리낌 없이 고친다면, 비 맞은 나무를 다시 쓸 수 있었던 것처럼 그 잘못은 다시 착한 사람이 되는 데 전혀 방해가 되지 않을 것입니다.

또한 여기서 끝날 일이 아닙니다. 나라의 정치도 이와 같습니다. 모든 일에서 백성에게 큰 해가 되는 것들을 이리저리 임시방편으로만 처리하고 제대로 바꾸지 않다가, 백성이 살기 어려워지고 나라가 위험해지고 나서야 갑자기 고치려 한다면 나라를 지키기 어려운 법입니다. 그러니 깊이 생각하지 않을 수 있겠습니까?
}


\kfSecHeader{kfAreaLiterature}{문제}{}

\begin{kfQBlock}{1}{}
\kfQStem{이 글의 글쓴이가 행랑채를 수리하며 얻은 가장 중심적인 깨달음은 무엇일까요?}
\begin{kfChoices}
\kfChoice 낡은 것은 모두 새것으로 바꾸는 것이 좋다.
\kfChoice 비가 오기 전에 미리 모든 것을 점검해야 한다.
\kfChoice 경험 많은 전문가에게 일을 맡겨야 실수가 없다.
\kfChoice 작은 문제라도 즉시 해결해야 큰 손해를 막을 수 있다.
\kfChoice 돈을 아끼지 말고 좋은 재료를 써야 집을 오래 쓸 수 있다.
\end{kfChoices}
\end{kfQBlock}
\begin{kfQBlock}{2}{}
\kfQStem{글쓴이가 자신의 생각을 전개하기 위해 사용한 주된 방법은 무엇인가요?}
\begin{kfChoices}
\kfChoice 여러 인물의 말을 인용하여 주장을 뒷받침하고 있다.
\kfChoice 어려운 질문을 던져 독자 스스로 답을 찾게 하고 있다.
\kfChoice 시간의 순서를 거슬러 올라가며 문제의 원인을 분석하고 있다.
\kfChoice 집을 고친 구체적인 경험을 다른 대상에 빗대어 설명하고 있다.
\kfChoice 상상 속의 인물을 내세워 자신의 생각을 대신 말하게 하고 있다.
\end{kfChoices}
\end{kfQBlock}
\begin{kfQBlock}{3}{}
\kfQStem{글쓴이가 궁극적으로 비판하고 경고하려는 대상은 누구일까요?}
\begin{kfChoices}
\kfChoice 부모님께 효도하지 않는 불효자들
\kfChoice 이웃과 자주 다투며 사이 나쁘게 지내는 사람들
\kfChoice 작은 이익에 눈이 멀어 큰 것을 보지 못하는 사람들
\kfChoice 백성을 힘들게 하는 문제를 알고도 개혁하지 않는 정치인들
\kfChoice 다른 사람의 충고를 듣지 않고 자기 고집대로만 하는 사람들
\end{kfChoices}
\end{kfQBlock}
\begin{kfQBlock}{4}{}
\kfQStem{글쓴이의 경험에 대한 설명으로 알맞지 \uline{않은} 것은 무엇인가요?}
\begin{kfChoices}
\kfChoice 글쓴이는 결국 돈이 아까워서 행랑채 수리를 포기했다.
\kfChoice 비가 오래 샌 곳은 재목이 썩어 수리 비용이 많이 들었다.
\kfChoice 비가 한 번 샌 곳은 재목이 멀쩡하여 수리비가 적게 들었다.
\kfChoice 글쓴이는 미리 손보지 않은 일을 후회하며 깨달음을 얻었다.
\kfChoice 행랑채 세 칸 가운데 두 칸은 오래전부터 비가 샜다.
\end{kfChoices}
\end{kfQBlock}
\begin{kfQBlock}{5}{}
\kfQStem{이 글의 내용과 관련 깊은 속담으로 가장 알맞은 것은 무엇일까요?}
\begin{kfChoices}
\kfChoice 싼 게 비지떡이다.
\kfChoice 티끌 모아 태산이다.
\kfChoice 소 잃고 외양간 고친다.
\kfChoice 윗물이 맑아야 아랫물이 맑다.
\kfChoice 구슬이 서 말이라도 꿰어야 보배다.
\end{kfChoices}
\end{kfQBlock}
\begin{kfQBlock}{6}{}
\kfQStem{이 글의 마지막 문장, "그러니 깊이 생각하지 않을 수 있겠습니까?"에 담긴 글쓴이의 의도는 무엇일까요?}
\begin{kfChoices}
\kfChoice 문제의 해결 방법을 몰라 독자에게 질문하고 있다.
\kfChoice 자신의 똑똑함을 다른 사람에게 자랑하고 싶어 한다.
\kfChoice 문제를 방치하면 위험하다는 것을 강조하며 경고하고 있다.
\kfChoice 자신의 깨달음이 맞는지 다른 사람에게 확인받고 싶어 한다.
\kfChoice 너무 어려운 문제라 이제 그만 생각하고 싶다는 뜻을 전한다.
\end{kfChoices}
\end{kfQBlock}
\begin{kfQBlock}{7}{}
\kfQStem{글쓴이가 “이런 일은 사람의 경우에도 똑같지 않을까요?”라고 말하며 전달하려는 교훈은 무엇인가요?}
\begin{kfChoices}
\kfChoice 한번 나빠진 건강은 다시 되돌리기 어렵다.
\kfChoice 잘못을 깨달았을 때 바로 고치면 다시 좋은 사람이 될 수 있다.
\kfChoice 몸이 아프면 즉시 병원에 가서 치료를 받아야 한다.
\kfChoice 좋은 사람이 되기 위해서는 항상 몸을 청결하게 해야 한다.
\kfChoice 다른 사람의 잘못을 너그럽게 용서해 주어야 한다.
\end{kfChoices}
\end{kfQBlock}
\begin{kfQBlock}{8}{}
\kfQStem{이 글에 대한 설명으로 가장 알맞은 것은 무엇일까요?}
\begin{kfChoices}
\kfChoice 아름다운 자연의 모습을 묘사하여 감동을 준다.
\kfChoice 상상 속의 인물을 내세워 재미있는 이야기를 들려준다.
\kfChoice 두 인물의 대화를 통해 갈등과 화해의 과정을 보여준다.
\kfChoice 자신의 실제 경험을 바탕으로 중요한 교훈을 이끌어낸다.
\kfChoice 구체적인 통계 자료를 제시하여 주장의 신뢰도를 높인다.
\end{kfChoices}
\end{kfQBlock}
\begin{kfQBlock}{1}{}
\kfQStem{이 글에 드러난 ‘오래 비가 샌 두 칸’과 ‘한 번 비가 샌 한 칸’의 차이점을 서술하세요.}
\kfAnswerNote{답안}{4}
\end{kfQBlock}
\begin{kfQBlock}{2}{}
\kfQStem{글쓴이가 생각하는 바람직한 정치인의 자세는 무엇인지, 집을 고치는 것에 빗대어 서술하세요.}
\kfAnswerNote{답안}{4}
\end{kfQBlock}

\kfSecHeader{kfAreaLiterature}{활동}{문장\_독해}
\noindent\textit{\small 다음 문장을 읽고 물음에 답하시오.}\par\medskip
\begin{kfSentenceUnit}{1}{}
\kfSentQ{1.}{다음 문장에서 서술어를 기준으로 각각의 주어는 무엇인가요?(1) 서술어 '무너질 것 같아서'의 주어는 무엇인가요?}
\begin{kfChoices}
\kfChoice 그가
\kfChoice 내가
\kfChoice 네가
\kfChoice 세 칸이
\end{kfChoices}
\kfSentQ{1.}{(2) 서술어 '고치게 되었습니다'의 주어는 무엇인가요?}
\begin{kfChoices}
\kfChoice 그가
\kfChoice 내가
\kfChoice 네가
\kfChoice 세 칸이
\end{kfChoices}
\end{kfSentenceUnit}

\begin{kfSentenceUnit}{2}{}
\kfSentQ{2.}{다음 문장에서 '그것'이 가리키는 것은 무엇인가요?}
\begin{kfChoices}
\kfChoice 두 칸에서 오래전부터 비가 새고 있다는 사실
\kfChoice 그 세 칸 중 두 칸은 오래전부터 비가 새고 있었다는 사실
\kfChoice 세 칸이 곧 무너질 것 같아서, 어쩔 수 없이 모두 고치게 되었다는 사실
\kfChoice 오래 비가 샌 곳은 서까래와 추녀, 기둥과 들보가 모두 썩어서 쓸 수 없게 되었다는 사실
\end{kfChoices}
\end{kfSentenceUnit}


\kfStartNonlit{비문학 영역 — [사회] 비싸지면 안 사? 꼭 그렇지는 않아!}


\kfSecHeader{kfAreaNonlit}{지문}{[사회] 비싸지면 안 사? 꼭 그렇지는 않아!}
\kfPassageNote{[사회] 비싸지면 안 사? 꼭 그렇지는 않아!}{%
물건의 가격이 오르면 사람들은 그 물건을 적게 사려고 한다. 이를 수요의 법칙이라고 한다. 그런데 물건마다 가격이 변할 때 사려는 양이 변하는 정도가 다르다. 어떤 물건은 가격 변화에 민감해서, 가격이 조금만 올라도 사려는 양이 크게 줄어든다. 이런 경우를 탄력적이라고 한다. 반대로 어떤 물건은 가격 변화에 둔감해서, 가격이 많이 올라도 사려는 양이 조금밖에 줄어들지 않는다. 이런 경우를 비탄력적이라고 한다.

물건의 가격이 변할 때 사려는 양이 얼마나 변하는지에 영향을 주는 요인은 세 가지가 있다. 첫째, 대신 쓸 수 있는 물건이 있는지이다. 버터 가격이 오르면 사람들은 대신 마가린을 살 수 있다. 그래서 버터 가격이 오르면 버터를 사려는 양이 크게 줄어든다. 하지만 달걀은 대신 쓸 수 있는 물건이 없어서 가격이 올라도 사려는 양이 크게 줄어들지 않는다.

둘째, 그 물건이 얼마나 필요한지이다. 휴지처럼 꼭 필요한 물건은 가격이 올라도 사려는 양이 크게 줄어들지 않는다. 하지만 보석처럼 없어도 되는 비싼 물건은 가격이 오르면 사려는 양이 크게 줄어든다.

셋째, 그 물건을 사는 데 드는 돈이 전체 용돈에서 차지하는 비중이다. 전체 용돈에서 많은 돈이 드는 물건일수록 가격이 오르면 사려는 양이 크게 줄어든다. 왜냐하면 그런 물건의 가격이 오르면 다른 물건을 사기 어려워지기 때문이다.

이런 변화 정도는 어떻게 계산할 수 있을까? 사려는 양이 변한 비율을 가격이 변한 비율로 나누면 된다. 예를 들어 아이스크림 가격이 10퍼센트 올랐는데 사려는 양이 20퍼센트 줄었다면, 20을 10으로 나누어 2가 된다. 이 숫자가 1보다 크면 탄력적이고, 1보다 작으면 비탄력적이다.

이런 특성은 가게 주인의 수입에 큰 영향을 준다. 전체 수입은 물건 가격에 팔린 개수를 곱한 것이다. 비탄력적인 물건은 가격을 올리면 전체 수입이 늘어난다. 하지만 탄력적인 물건은 가격을 올리면 전체 수입이 줄어든다. 따라서 가게 주인은 물건의 특성을 잘 파악해야 한다.
}


\kfSecHeader{kfAreaNonlit}{문제}{}

\begin{kfQBlock}{1}{}
\kfQStem{윗글의 내용과 일치하지 않는 것은 무엇인가요?}
\begin{kfChoices}
\kfChoice 용돈에서 차지하는 비중이 작을수록 수요가 탄력적이다.
\kfChoice 보석은 휴지보다 수요가 더 탄력적인 경향이 있다.
\kfChoice 달걀은 버터보다 수요가 더 비탄력적인 경향이 있다.
\kfChoice 가격이 오르면 수요량은 줄어드는 것이 일반적이다.
\kfChoice 수요의 탄력성은 가게의 수입과 관련이 깊다.
\end{kfChoices}
\end{kfQBlock}
\begin{kfQBlock}{2}{}
\kfQStem{윗글을 바탕으로 <보기>의 A, B 상품을 이해한 내용으로 적절하지 \uline{않은} 것은 무엇인가요?}
\begin{kfEvidence}A상품: 가격이 5\% 오르자, 사려는 사람이 1\% 줄었다. B상품: 가격이 5\% 오르자, 사려는 사람이 10\% 줄었다.\end{kfEvidence}
\begin{kfChoices}
\kfChoice B상품의 가격을 올리면 판매 수입이 늘어날 것이다.
\kfChoice B상품은 가격 변화율(5\%)보다 수요 변화율(10\%)이 커서 탄력적이다.
\kfChoice A상품은 B상품보다 수요가 가격 변화에 덜 민감하다.
\kfChoice A상품의 가격을 올리면 판매 수입이 늘어날 것이다.
\kfChoice A상품은 가격 변화율(5\%)보다 수요 변화율(1\%)이 작아 비탄력적이다.
\end{kfChoices}
\end{kfQBlock}
\begin{kfQBlock}{3}{}
\kfQStem{다음 중 수요가 가장 '비탄력적'일 것으로 예상되는 물건은 무엇인가요?}
\begin{kfChoices}
\kfChoice 새로 나온 비싼 게임기
\kfChoice 유행이 지난 명품 가방
\kfChoice 옆집 가게에서도 파는 콜라
\kfChoice 우리 가족에게 꼭 필요한 약
\kfChoice 맛이 비슷한 여러 회사의 과자
\end{kfChoices}
\end{kfQBlock}
\begin{kfQBlock}{4}{}
\kfQStem{<보기>의 상황을 윗글에 근거하여 분석한 내용으로 가장 적절한 것은 무엇인가요?}
\begin{kfEvidence}한 문구점에서 연필 가격을 100원에서 120원으로 올렸더니, 하루에 팔리던 양이 100자루에서 60자루로 줄었다.\end{kfEvidence}
\begin{kfChoices}
\kfChoice 수요가 비탄력적이므로 가격을 올린 것은 좋은 전략이다.
\kfChoice 수요가 탄력적이므로 가격을 더 올려야만 수입이 늘어난다.
\kfChoice 수요가 탄력적이므로 가격을 올린 결과 전체 수입이 줄어들었다.
\kfChoice 수요가 비탄력적이므로 가격을 원래대로 내려야 수입이 늘어난다.
\kfChoice 가격 변화와 수요량 변화가 복잡하여 탄력성을 계산할 수 없다.
\end{kfChoices}
\end{kfQBlock}
\begin{kfQBlock}{5}{}
\kfQStem{다음 <보기>의 일기를 읽고, 밑줄 친 물건 중 수요가 가장 '탄력적'인 것은 무엇인가요?}
\begin{kfEvidence}오늘 용돈을 받았다. 문구점에 가서 준비물을 샀다.

먼저, 학교 미술 시간에 필요한 ㉠\uline{도화지}를 샀다. 가격이 좀 올랐지만 안 살 수가 없었다.

그리고 평소 갖고 싶었던 ㉡\uline{유명 캐릭터 필통}을 집었다가 가격표를 보고 깜짝 놀랐다. 너무 비싸서 그냥 내려놓고 ㉢\uline{저렴한 천 필통}을 샀다.

집에 오는 길에 배가 고파서 ㉣\uline{삼각김밥}을 사 먹고, 목이 말라 ㉤\uline{생수}도 한 병 샀다.\end{kfEvidence}
\begin{kfChoices}
\kfChoice ㉠ 도화지
\kfChoice ㉡ 유명 캐릭터 필통
\kfChoice ㉢ 저렴한 천 필통
\kfChoice ㉣ 삼각김밥
\kfChoice ㉤ 생수
\end{kfChoices}
\end{kfQBlock}
\begin{kfQBlock}{6}{}
\kfQStem{다음 <보기>는 어느 빵집의 가격 변화에 따른 판매량 조사를 나타낸 표입니다. 이를 바르게 분석한 것은 무엇인가요?}
\begin{kfEvidence}| 종류 | 가격 변화 | 판매량(사려는 양) 변화 | | :---- | :---- | :---- | | 우유 식빵 | 10\% 오름 | 2\% 줄어듦 | | 딸기 케이크 | 10\% 오름 | 30\% 줄어듦 |\end{kfEvidence}
\begin{kfChoices}
\kfChoice 식빵은 탄력적이다.
\kfChoice 케이크는 탄력적이다.
\kfChoice 케이크는 꼭 필요한 음식이다.
\kfChoice 식빵은 대신할 음식이 아주 많다.
\kfChoice 식빵의 가격을 올리면 손해를 본다.
\end{kfChoices}
\end{kfQBlock}
\begin{kfQBlock}{7}{}
\kfQStem{이 글을 읽은 가게 주인이 얻을 수 있는 교훈으로 가장 적절한 것은 무엇인가요?}
\begin{kfChoices}
\kfChoice 모든 물건의 가격은 무조건 싸게 팔아야 한다.
\kfChoice 가게의 수입을 늘리려면 물건의 특성을 파악하고 가격을 정해야 한다.
\kfChoice 손님들의 용돈이 얼마나 되는지 파악하는 것이 가장 중요하다.
\kfChoice 물건의 종류와 상관없이 가격을 올리면 항상 수입이 늘어난다.
\kfChoice 물건을 팔 때는 가게의 위치가 수입에 가장 큰 영향을 끼친다.
\end{kfChoices}
\end{kfQBlock}
\begin{kfQBlock}{1}{}
\kfQStem{물건의 가격이 오르면 사람들이 그 물건을 적게 사려고 하는 것을 무슨 법칙이라고 하나요?}
\kfAnswerLines{1}
\end{kfQBlock}
\begin{kfQBlock}{2}{}
\kfQStem{가격이 조금만 올라도 사람들이 사려는 양이 크게 줄어드는 경우를 무엇이라고 하나요?}
\kfAnswerLines{1}
\end{kfQBlock}
\begin{kfQBlock}{3}{}
\kfQStem{가격이 많이 올라도 사람들이 사려는 양이 조금밖에 줄어들지 않는 경우를 무엇이라고 하나요?}
\kfAnswerLines{1}
\end{kfQBlock}
\begin{kfQBlock}{4}{}
\kfQStem{수요량의 변화 정도에 영향을 미치는 요인 세 가지를 모두 쓰세요.}
\kfAnswerLines{1}
\end{kfQBlock}
\begin{kfQBlock}{5}{}
\kfQStem{휴지처럼 꼭 필요한 물건은 가격 변화에 탄력적인가, 비탄력적인가요?}
\kfAnswerLines{1}
\end{kfQBlock}
\begin{kfQBlock}{1}{}
\kfQStem{학생들이 매일 먹어야 하는 학교 급식비가 오른다고 합니다. 이 학교 급식의 수요 특성을 예상하고, 급식비가 올랐을 때 학교의 전체 수입이 어떻게 변할지 지문의 내용을 바탕으로 서술하세요.}
\kfAnswerNote{답안}{4}
\end{kfQBlock}

\kfSecHeader{kfAreaNonlit}{활동}{문장\_독해}
\noindent\textit{\small 다음 문장을 읽고 물음에 답하시오.}\par\medskip
\begin{kfSentenceUnit}{1}{}
\kfSentQ{1.}{이 문장에서 원인은 무엇인가요?}
\begin{kfChoices}
\kfChoice 물건의 가격이 오르는 것.
\kfChoice 물건의 품질이 떨어지는 것.
\kfChoice 가게 주인이 문을 닫는 것.
\kfChoice 사람들이 물건을 사지 않는 것.
\end{kfChoices}
\end{kfSentenceUnit}

\begin{kfSentenceUnit}{2}{}
\kfSentQ{2.}{이 문장에서 결과는 무엇인가요?}
\begin{kfChoices}
\kfChoice 가게의 수입이 늘어나는 것.
\kfChoice 물건의 가격이 급격히 떨어지는 것.
\kfChoice 사람들이 더 많은 물건을 찾는 것.
\kfChoice 사람들이 그 물건을 적게 사려고 하는 것.
\end{kfChoices}
\end{kfSentenceUnit}

\begin{kfSentenceUnit}{3}{}
\kfSentQ{3.}{가격이 오르면 사람들은 어떻게 하나요?}
\begin{kfChoices}
\kfChoice 아무런 변화도 보이지 않는다.
\kfChoice 그 물건을 적게 사려고 한다.
\kfChoice 친구들에게 그 물건을 추천한다.
\kfChoice 그 물건을 더 많이 사려고 한다.
\end{kfChoices}
\end{kfSentenceUnit}

\begin{kfSentenceUnit}{4}{}
\kfSentQ{4.}{'이를'이 가리키는 내용은 무엇인가요?}
\begin{kfChoices}
\kfChoice 물건의 종류가 다양하면 사람들이 더 많이 찾는 것.
\kfChoice 물건의 품질이 좋으면 가격에 상관없이 많이 팔리는 것.
\kfChoice 물건의 가격이 오르면 사람들이 그 물건을 적게 사려고 하는 것.
\kfChoice 물건의 가격이 내리면 사람들이 그 물건을 더 많이 사려고 하는 것.
\end{kfChoices}
\end{kfSentenceUnit}

\begin{kfSentenceUnit}{5}{}
\kfSentQ{5.}{'수요의 법칙'이란 무엇인가요?}
\begin{kfChoices}
\kfChoice 경제가 어려워지면 물건 가격이 내려가는 현상.
\kfChoice 물건이 많이 팔리면 가격이 자연스럽게 올라가는 현상.
\kfChoice 사람들이 돈을 많이 벌면 물건을 더 많이 사는 현상.
\kfChoice 물건의 가격이 오르면 사람들이 그 물건을 적게 사려고 하는 현상.
\end{kfChoices}
\end{kfSentenceUnit}

\begin{kfSentenceUnit}{6}{}
\kfSentQ{6.}{무엇이 물건마다 다르다고 했나요?}
\begin{kfChoices}
\kfChoice 물건의 품질이 변하는 정도.
\kfChoice 가게의 위치가 변하는 정도.
\kfChoice 사람들의 소득이 변하는 정도.
\kfChoice 가격이 변할 때 사려는 양이 변하는 정도.
\end{kfChoices}
\end{kfSentenceUnit}

\begin{kfSentenceUnit}{7}{}
\kfSentQ{7.}{언제 사려는 양이 변하나요?}
\begin{kfChoices}
\kfChoice 가격이 변할 때.
\kfChoice 날씨가 변할 때.
\kfChoice 가게 주인이 바뀔 때.
\kfChoice 사람들의 기분이 변할 때.
\end{kfChoices}
\end{kfSentenceUnit}

\begin{kfSentenceUnit}{8}{}
\kfSentQ{8.}{가격 변화에 민감한 물건의 특징은 무엇인가요?}
\begin{kfChoices}
\kfChoice 가격과는 상관없이 사려는 양이 일정하다.
\kfChoice 가격이 조금만 올라도 사려는 양이 크게 줄어든다.
\kfChoice 가격이 조금 오르면 사려는 양이 오히려 늘어난다.
\kfChoice 가격이 많이 올라도 사려는 양은 거의 변하지 않는다.
\end{kfChoices}
\end{kfSentenceUnit}

\begin{kfSentenceUnit}{9}{}
\kfSentQ{9.}{'이런 경우'가 가리키는 것은 무엇인가요?}
\begin{kfChoices}
\kfChoice 가격 변화에 둔감한 경우.
\kfChoice 가격 변화에 민감한 경우.
\kfChoice 가격 변화와 상관없는 경우.
\kfChoice 가격 변화를 예측할 수 없는 경우.
\end{kfChoices}
\end{kfSentenceUnit}

\begin{kfSentenceUnit}{10}{}
\kfSentQ{10.}{'탄력적'이라는 것은 무엇을 의미하나요?}
\begin{kfChoices}
\kfChoice 가격 변화에 상관없이 항상 일정한 양이 팔리는 경우.
\kfChoice 가격이 내리면 사려는 양이 오히려 줄어드는 특이한 경우.
\kfChoice 가격 변화에 둔감해서, 가격이 올라도 사려는 양이 변하지 않는 경우.
\kfChoice 가격 변화에 민감해서, 가격이 조금만 올라도 사려는 양이 크게 줄어드는 경우.
\end{kfChoices}
\end{kfSentenceUnit}

\begin{kfSentenceUnit}{11}{}
\kfSentQ{11.}{가격 변화에 둔감한 물건의 특징은 무엇인가요?}
\begin{kfChoices}
\kfChoice 가격이 오르면 사려는 양도 같이 늘어난다.
\kfChoice 가격이 내리면 사려는 양이 오히려 줄어든다.
\kfChoice 가격이 조금만 올라도 사려는 양이 아주 크게 줄어든다.
\kfChoice 가격이 많이 올라도 사려는 양이 조금밖에 줄어들지 않는다.
\end{kfChoices}
\end{kfSentenceUnit}

\begin{kfSentenceUnit}{12}{}
\kfSentQ{12.}{'이런 경우'가 가리키는 것은 무엇인가요?}
\begin{kfChoices}
\kfChoice 가격 변화에 둔감한 경우.
\kfChoice 가격이 자주 변하는 경우.
\kfChoice 가격이 전혀 변하지 않는 경우.
\kfChoice 가격 변화에 아주 민감한 경우.
\end{kfChoices}
\end{kfSentenceUnit}

\begin{kfSentenceUnit}{13}{}
\kfSentQ{13.}{'비탄력적'이라는 것은 무엇을 의미하나요?}
\begin{kfChoices}
\kfChoice 가격 변화에 둔감해서, 가격이 많이 올라도 사려는 양이 조금밖에 줄어들지 않는 경우.
\kfChoice 가격이 변화하더라도 수요량이 가격보다는 유행이나 다른 요인에 따라 불규칙하게 변하는 특이한 경우.
\kfChoice 일반적인 수요 법칙과 다르게 가격이 오르면 사람들이 그것을 과시하기 위해 오히려 더 많이 사려는 경우.
\kfChoice 가격의 아주 작은 변화에도 소비자들이 매우 민감하게 반응하여 사려는 양이 급격하게 늘거나 줄어드는 경우.
\end{kfChoices}
\end{kfSentenceUnit}

\begin{kfSentenceUnit}{14}{}
\kfSentQ{14.}{이 문장을 통해 앞으로 어떤 내용이 이어질지 예측한 것으로 옳은 것은 무엇인가요?.}
\begin{kfChoices}
\kfChoice 물건의 가격을 결정하는 방법 세 가지가 설명될 것이다.
\kfChoice 사람들이 물건을 사지 않는 이유 세 가지가 설명될 것이다.
\kfChoice 가게 주인이 돈을 많이 버는 방법 세 가지가 설명될 것이다.
\kfChoice 사려는 양의 변화에 영향을 주는 요인 세 가지가 설명될 것이다.
\end{kfChoices}
\end{kfSentenceUnit}

\begin{kfSentenceUnit}{15}{}
\kfSentQ{15.}{수요량 변화에 영향을 주는 첫 번째 요인은 무엇인가요?}
\begin{kfChoices}
\kfChoice 물건을 홍보하는 방법.
\kfChoice 물건을 파는 가게의 위치.
\kfChoice 물건을 만드는 재료의 종류.
\kfChoice 대신 쓸 수 있는 물건의 유무.
\end{kfChoices}
\end{kfSentenceUnit}

\begin{kfSentenceUnit}{16}{}
\kfSentQ{16.}{버터 가격이 오르면 사람들은 무엇을 살 수 있나요?}
\begin{kfChoices}
\kfChoice 치즈.
\kfChoice 식빵.
\kfChoice 우유.
\kfChoice 마가린.
\end{kfChoices}
\end{kfSentenceUnit}

\begin{kfSentenceUnit}{17}{}
\kfSentQ{17.}{'그래서'의 구체적인 내용은 무엇인가요?}
\begin{kfChoices}
\kfChoice 버터 가격이 오르면 우유 가격도 같이 오르기 때문에.
\kfChoice 버터 가격이 오르면 버터의 품질이 더 좋아지기 때문에.
\kfChoice 버터 가격이 오르면 사람들이 빵을 덜 먹게 되기 때문에.
\kfChoice 버터 가격이 오르면 사람들이 대신 마가린을 살 수 있어서.
\end{kfChoices}
\end{kfSentenceUnit}

\begin{kfSentenceUnit}{18}{}
\kfSentQ{18.}{버터를 사려는 양이 크게 줄어드는 이유는 무엇인가요?}
\begin{kfChoices}
\kfChoice 버터는 건강에 아주 좋기 때문이다.
\kfChoice 대신 마가린을 살 수 있기 때문이다.
\kfChoice 마가린은 버터보다 맛이 없기 때문이다.
\kfChoice 버터가 없으면 빵을 먹을 수 없기 때문이다.
\end{kfChoices}
\end{kfSentenceUnit}

\begin{kfSentenceUnit}{19}{}
\kfSentQ{19.}{버터를 사려는 양은 어떻게 변하나요?}
\begin{kfChoices}
\kfChoice 크게 줄어든다.
\kfChoice 조금만 줄어든다.
\kfChoice 오히려 늘어난다.
\kfChoice 거의 변하지 않는다.
\end{kfChoices}
\end{kfSentenceUnit}

\begin{kfSentenceUnit}{20}{}
\kfSentQ{20.}{버터는 탄력적인가요, 비탄력적인가요?}
\begin{kfChoices}
\kfChoice 규칙적.
\kfChoice 탄력적.
\kfChoice 불규칙적.
\kfChoice 비탄력적.
\end{kfChoices}
\end{kfSentenceUnit}

\begin{kfSentenceUnit}{21}{}
\kfSentQ{21.}{'하지만'을 기준으로, 버터의 사례와 달걀의 사례는 어떤 차이점을 보이나요?}
\begin{kfChoices}
\kfChoice 파는 장소의 차이.
\kfChoice 만드는 재료의 차이.
\kfChoice 가격의 비싸고 싼 차이.
\kfChoice 대체할 수 있는 물건의 유무.
\end{kfChoices}
\end{kfSentenceUnit}

\begin{kfSentenceUnit}{22}{}
\kfSentQ{22.}{달걀은 가격이 오르면 사려는 양이 어떻게 되나요?}
\begin{kfChoices}
\kfChoice 크게 늘어난다.
\kfChoice 조금만 줄어든다.
\kfChoice 아예 사지 않는다.
\kfChoice 크게 줄어들지 않는다.
\end{kfChoices}
\end{kfSentenceUnit}

\begin{kfSentenceUnit}{23}{}
\kfSentQ{23.}{달걀은 왜 가격이 올라도 사려는 양이 크게 줄지 않나요?}
\begin{kfChoices}
\kfChoice 대신 쓸 수 있는 물건이 없기 때문이다.
\kfChoice 달걀은 맛이 좋아서 포기할 수 없기 때문이다.
\kfChoice 달걀은 가격이 저렴해서 부담이 없기 때문이다.
\kfChoice 달걀을 대신할 만한 마땅한 식재료가 없기 때문이다.
\end{kfChoices}
\end{kfSentenceUnit}

\begin{kfSentenceUnit}{24}{}
\kfSentQ{24.}{달걀은 탄력적인가요, 비탄력적인가요?}
\begin{kfChoices}
\kfChoice 탄력적.
\kfChoice 효율적.
\kfChoice 비효율적.
\kfChoice 비탄력적.
\end{kfChoices}
\end{kfSentenceUnit}

\begin{kfSentenceUnit}{25}{}
\kfSentQ{25.}{수요량 변화에 영향을 주는 두 번째 요인은 무엇인가요?}
\begin{kfChoices}
\kfChoice 그 물건의 디자인.
\kfChoice 그 물건의 생산 국가.
\kfChoice 그 물건의 가격 수준.
\kfChoice 그 물건의 필요성 정도.
\end{kfChoices}
\end{kfSentenceUnit}

\begin{kfSentenceUnit}{26}{}
\kfSentQ{26.}{휴지는 어떤 물건인가요?}
\begin{kfChoices}
\kfChoice 꼭 필요한 물건.
\kfChoice 유행에 민감한 물건.
\kfChoice 가격이 매우 비싼 사치품.
\kfChoice 없어도 생활하는 데 지장이 없는 물건.
\end{kfChoices}
\end{kfSentenceUnit}

\begin{kfSentenceUnit}{27}{}
\kfSentQ{27.}{꼭 필요한 물건은 가격이 올라도 어떻게 되나요?}
\begin{kfChoices}
\kfChoice 아예 사지 않게 된다.
\kfChoice 사려는 양이 오히려 늘어난다.
\kfChoice 사려는 양이 아주 크게 줄어든다.
\kfChoice 사려는 양이 크게 줄어들지 않는다.
\end{kfChoices}
\end{kfSentenceUnit}

\begin{kfSentenceUnit}{28}{}
\kfSentQ{28.}{휴지는 탄력적인가요, 비탄력적인가요?}
\begin{kfChoices}
\kfChoice 탄력적.
\kfChoice 가변적.
\kfChoice 고정적.
\kfChoice 비탄력적.
\end{kfChoices}
\end{kfSentenceUnit}

\begin{kfSentenceUnit}{29}{}
\kfSentQ{29.}{'하지만'을 기준으로, 이 문장은 휴지의 사례와 어떤 점에서 대조되나요?}
\begin{kfChoices}
\kfChoice 휴지와 달리 보석은 누구나 가지고 있다는 점.
\kfChoice 휴지와 달리 보석은 쉽게 구할 수 있다는 점.
\kfChoice 휴지와 달리 보석은 가격이 매우 저렴하다는 점.
\kfChoice 휴지와 달리 보석은 꼭 필요하지 않은 물건이라는 점.
\end{kfChoices}
\end{kfSentenceUnit}

\begin{kfSentenceUnit}{30}{}
\kfSentQ{30.}{보석은 어떤 물건인가요?}
\begin{kfChoices}
\kfChoice 없어도 되는 비싼 물건.
\kfChoice 생활에 꼭 필요한 필수품.
\kfChoice 꼭 필요하고 가격도 비싼 물건.
\kfChoice 없어도 되지만 가격이 싼 물건.
\end{kfChoices}
\end{kfSentenceUnit}

\begin{kfSentenceUnit}{31}{}
\kfSentQ{31.}{없어도 되는 물건은 가격이 오르면 어떻게 되나요?}
\begin{kfChoices}
\kfChoice 사려는 양이 크게 줄어든다.
\kfChoice 가격에 상관없이 잘 팔린다.
\kfChoice 사려는 양이 오히려 늘어난다.
\kfChoice 사려는 양이 거의 변하지 않는다.
\end{kfChoices}
\end{kfSentenceUnit}

\begin{kfSentenceUnit}{32}{}
\kfSentQ{32.}{보석은 탄력적인가요, 비탄력적인가요?}
\begin{kfChoices}
\kfChoice 탄력적.
\kfChoice 필수적.
\kfChoice 안정적.
\kfChoice 비탄력적.
\end{kfChoices}
\end{kfSentenceUnit}

\begin{kfSentenceUnit}{33}{}
\kfSentQ{33.}{수요량 변화에 영향을 주는 세 번째 요인은 무엇인가요?}
\begin{kfChoices}
\kfChoice 물건이 만들어진 날짜부터 지금까지의 기간.
\kfChoice 물건을 만든 나라가 어디인지 나타내는 원산지.
\kfChoice 물건이 얼마나 유명한 브랜드인지 나타내는 정도.
\kfChoice 물건 구매 비용이 전체 용돈에서 차지하는 비중.
\end{kfChoices}
\end{kfSentenceUnit}

\begin{kfSentenceUnit}{34}{}
\kfSentQ{34.}{어떤 물건일수록 사려는 양이 크게 줄어드나요?}
\begin{kfChoices}
\kfChoice 친구들이 많이 가지고 있는 물건.
\kfChoice 가격이 매우 싸서 부담 없는 물건.
\kfChoice 전체 용돈에서 많은 돈이 드는 물건.
\kfChoice 전체 용돈에서 아주 적은 돈이 드는 물건.
\end{kfChoices}
\end{kfSentenceUnit}

\begin{kfSentenceUnit}{35}{}
\kfSentQ{35.}{많은 돈이 드는 물건은 가격이 오르면 어떻게 되나요?}
\begin{kfChoices}
\kfChoice 사려는 양이 크게 줄어든다.
\kfChoice 사려는 양이 조금만 줄어든다.
\kfChoice 사려는 양이 오히려 늘어난다.
\kfChoice 사려는 양에 아무런 변화가 없다.
\end{kfChoices}
\end{kfSentenceUnit}

\begin{kfSentenceUnit}{36}{}
\kfSentQ{36.}{용돈에서 많은 돈이 드는 물건은 탄력적인가요, 비탄력적인가요?}
\begin{kfChoices}
\kfChoice 독립적.
\kfChoice 탄력적.
\kfChoice 의존적.
\kfChoice 비탄력적.
\end{kfChoices}
\end{kfSentenceUnit}

\begin{kfSentenceUnit}{37}{}
\kfSentQ{37.}{'왜냐하면'이 나타내는 관계는 무엇인가요?}
\begin{kfChoices}
\kfChoice 앞 문장의 결과.
\kfChoice 앞 문장의 예시.
\kfChoice 앞 문장의 이유.
\kfChoice 앞 문장과 반대되는 내용.
\end{kfChoices}
\end{kfSentenceUnit}

\begin{kfSentenceUnit}{38}{}
\kfSentQ{38.}{비싼 물건의 가격이 오르면 왜 사려는 양이 줄어드나요?}
\begin{kfChoices}
\kfChoice 비싼 물건은 품질이 좋기 때문이다.
\kfChoice 다른 물건을 사기 어려워지기 때문이다.
\kfChoice 비싼 물건은 친구들에게 자랑할 수 있기 때문이다.
\kfChoice 비싼 물건은 나중에 가격이 더 오를 수 있기 때문이다.
\end{kfChoices}
\end{kfSentenceUnit}

\begin{kfSentenceUnit}{39}{}
\kfSentQ{39.}{'이런 변화 정도'는 무엇을 가리키나요?}
\begin{kfChoices}
\kfChoice 가게의 매출이 변하는 정도.
\kfChoice 물건의 품질이 변하는 정도.
\kfChoice 사람들의 유행이 변하는 정도.
\kfChoice 가격이 변할 때 사려는 양이 변하는 정도.
\end{kfChoices}
\end{kfSentenceUnit}

\begin{kfSentenceUnit}{40}{}
\kfSentQ{40.}{이 질문을 통해 앞으로 어떤 내용이 나올 것을 예상할 수 있나요?}
\begin{kfChoices}
\kfChoice 용돈을 아껴 쓰는 방법.
\kfChoice 물건을 싸게 사는 방법.
\kfChoice 가게 주인이 되는 방법.
\kfChoice 변화 정도를 계산하는 방법.
\end{kfChoices}
\end{kfSentenceUnit}

\begin{kfSentenceUnit}{41}{}
\kfSentQ{41.}{변화 정도를 계산하는 방법은 무엇인가요?}
\begin{kfChoices}
\kfChoice (가격이 변한 비율) ÷ (사려는 양이 변한 비율).
\kfChoice (가격이 변한 비율) - (사려는 양이 변한 비율).
\kfChoice (사려는 양이 변한 비율) × (가격이 변한 비율).
\kfChoice (사려는 양이 변한 비율) ÷ (가격이 변한 비율).
\end{kfChoices}
\end{kfSentenceUnit}

\begin{kfSentenceUnit}{42}{}
\kfSentQ{42.}{아이스크림 가격은 몇 퍼센트 올랐나요?}
\begin{kfChoices}
\kfChoice 5퍼센트.
\kfChoice 20퍼센트.
\kfChoice 10퍼센트.
\kfChoice 15퍼센트.
\end{kfChoices}
\end{kfSentenceUnit}

\begin{kfSentenceUnit}{43}{}
\kfSentQ{43.}{사려는 양은 몇 퍼센트 줄었나요?}
\begin{kfChoices}
\kfChoice 10퍼센트.
\kfChoice 40퍼센트.
\kfChoice 30퍼센트.
\kfChoice 20퍼센트.
\end{kfChoices}
\end{kfSentenceUnit}

\begin{kfSentenceUnit}{44}{}
\kfSentQ{44.}{이 예시에서 숫자 '20'과 '10'은 각각 무엇을 대신 나타낸 값인가요?}
\begin{kfChoices}
\kfChoice 5
\kfChoice 10
\kfChoice 20
\kfChoice 30
\end{kfChoices}
\end{kfSentenceUnit}

\begin{kfSentenceUnit}{45}{}
\kfSentQ{45.}{이 예시에서 계산 결과인 2는 어떻게 나온 것인가요?}
\begin{kfChoices}
\kfChoice 변화율을 서로 더해서 평균을 낸 값이다.
\kfChoice 그냥 문제에서 주어진 값을 그대로 옮겨 적은 것이다.
\kfChoice 변화율(20\%)에 2를 곱해서 간단히 계산해 낸 값이다.
\kfChoice 사려는 양 변화율(20\%)을 가격 변화율(10\%)로 나누어서 나왔다.
\end{kfChoices}
\end{kfSentenceUnit}

\begin{kfSentenceUnit}{46}{}
\kfSentQ{46.}{'이 숫자'가 가리키는 것은 무엇인가요?}
\begin{kfChoices}
\kfChoice 원래 물건 가격에서 나중에 바뀐 가격을 뺀 값.
\kfChoice 가격과 수량이 얼마나 변했는지 퍼센트로 나타낸 값.
\kfChoice 물건 하나당 가격에 지금까지 팔린 개수를 곱한 값.
\kfChoice 사려는 양이 변한 비율을 가격이 변한 비율로 나눈 값.
\end{kfChoices}
\end{kfSentenceUnit}

\begin{kfSentenceUnit}{47}{}
\kfSentQ{47.}{계산 결과가 '탄력적'이라고 판단하는 기준은 무엇인가요?}
\begin{kfChoices}
\kfChoice 사려는 양이 변한 비율을 가격이 변한 비율로 나눈 값이 1보다 큰 경우.
\kfChoice 사려는 양이 변한 비율을 가격이 변한 비율로 나눈 값이 0에 가까운 경우.
\kfChoice 사려는 양이 변한 비율을 가격이 변한 비율로 나눈 값이 1보다 훨씬 작은 경우.
\kfChoice 사려는 양이 변한 비율을 가격이 변한 비율로 나눈 값이 1과 정확히 같은 경우.
\end{kfChoices}
\end{kfSentenceUnit}

\begin{kfSentenceUnit}{48}{}
\kfSentQ{48.}{계산 결과가 '비탄력적'이라고 판단하는 기준은 무엇인가요?}
\begin{kfChoices}
\kfChoice 사려는 양이 변한 비율을 가격이 변한 비율로 나눈 값이 1보다 작은 경우.
\kfChoice 사려는 양이 변한 비율을 가격이 변한 비율로 나눈 값이 0보다 작은 경우.
\kfChoice 사려는 양이 변한 비율을 가격이 변한 비율로 나눈 값이 1보다 훨씬 큰 경우.
\kfChoice 사려는 양이 변한 비율을 가격이 변한 비율로 나눈 값이 1과 정확히 같은 경우.
\end{kfChoices}
\end{kfSentenceUnit}

\begin{kfSentenceUnit}{49}{}
\kfSentQ{49.}{위 예시에서 아이스크림은 탄력적인가요, 비탄력적인가요?}
\begin{kfChoices}
\kfChoice 알 수 없음.
\kfChoice 비탄력적 (1보다 작으므로).
\kfChoice 단위 탄력적 (1과 같으므로).
\kfChoice 탄력적 (2는 1보다 크므로).
\end{kfChoices}
\end{kfSentenceUnit}

\begin{kfSentenceUnit}{50}{}
\kfSentQ{50.}{'이런 특성'이 가리키는 것을 무엇인가요?}
\begin{kfChoices}
\kfChoice 탄력적이거나 비탄력적인 특성.
\kfChoice 물건의 겉모습이 얼마나 예쁜지 나타내는 디자인.
\kfChoice 물건이 어느 나라에서 만들어졌는지에 대한 정보.
\kfChoice 물건이 얼마나 유명한 회사 제품인지 알려주는 브랜드.
\end{kfChoices}
\end{kfSentenceUnit}

\begin{kfSentenceUnit}{51}{}
\kfSentQ{51.}{누구의 수입에 영향을 주나요?}
\begin{kfChoices}
\kfChoice 정부.
\kfChoice 가게 주인.
\kfChoice 외국인 관광객.
\kfChoice 물건을 사는 소비자.
\end{kfChoices}
\end{kfSentenceUnit}

\begin{kfSentenceUnit}{52}{}
\kfSentQ{52.}{전체 수입은 어떻게 계산하나요?}
\begin{kfChoices}
\kfChoice 물건 가격에 팔린 개수를 더해서 (전체 수입) = (물건 가격) + (팔린 개수).
\kfChoice 물건 가격에 팔린 개수를 곱해서 (전체 수입) = (물건 가격) × (팔린 개수).
\kfChoice 물건 가격을 팔린 개수로 나누어서 (전체 수입) = (물건 가격) ÷ (팔린 개수).
\kfChoice 팔린 개수를 물건 가격으로 나누어서 (전체 수입) = (팔린 개수) ÷ (물건 가격).
\end{kfChoices}
\end{kfSentenceUnit}

\begin{kfSentenceUnit}{53}{}
\kfSentQ{53.}{어떤 물건은 가격을 올리면 전체 수입이 늘어나나요?}
\begin{kfChoices}
\kfChoice 모든 물건.
\kfChoice 탄력적인 물건.
\kfChoice 비탄력적인 물건.
\kfChoice 단위 탄력적인 물건.
\end{kfChoices}
\end{kfSentenceUnit}

\begin{kfSentenceUnit}{54}{}
\kfSentQ{54.}{비탄력적인 물건의 가격을 올리면 전체 수입이 어떻게 되나요?}
\begin{kfChoices}
\kfChoice 늘어난다.
\kfChoice 줄어든다.
\kfChoice 알 수 없다.
\kfChoice 변하지 않는다.
\end{kfChoices}
\end{kfSentenceUnit}

\begin{kfSentenceUnit}{55}{}
\kfSentQ{55.}{어떤 물건은 가격을 올리면 전체 수입이 줄어드나요?}
\begin{kfChoices}
\kfChoice 모든 물건.
\kfChoice 탄력적인 물건.
\kfChoice 비탄력적인 물건.
\kfChoice 단위 탄력적인 물건.
\end{kfChoices}
\end{kfSentenceUnit}

\begin{kfSentenceUnit}{56}{}
\kfSentQ{56.}{탄력적인 물건의 가격을 올리면 전체 수입이 어떻게 되나요?}
\begin{kfChoices}
\kfChoice 늘어난다.
\kfChoice 줄어든다.
\kfChoice 알 수 없다.
\kfChoice 변하지 않는다.
\end{kfChoices}
\end{kfSentenceUnit}

\begin{kfSentenceUnit}{57}{}
\kfSentQ{57.}{누가 무엇을 잘 파악해야 하나요?}
\begin{kfChoices}
\kfChoice 손님이 물건의 가격을 꼼꼼히 비교해야 한다.
\kfChoice 직원이 손님을 아주 친절하게 맞이해야 한다.
\kfChoice 가게 주인이 물건의 특성을 잘 파악해야 한다.
\kfChoice 정부가 세금을 정확하게 계산해서 걷어야 한다.
\end{kfChoices}
\end{kfSentenceUnit}

\begin{kfSentenceUnit}{58}{}
\kfSentQ{58.}{가게 주인이 물건의 특성을 잘 파악해야 하는 이유는 무엇인가요?}
\begin{kfChoices}
\kfChoice 나라에 내야 할 세금을 정확하게 계산하기 위해서.
\kfChoice 가게 일을 도울 직원을 더 많이 채용하기 위해서.
\kfChoice 가게를 더 넓은 곳으로 확장하여 이사하기 위해서.
\kfChoice 탄력성에 따라 가격을 올렸을 때 수입이 달라지기 때문이다.
\end{kfChoices}
\end{kfSentenceUnit}



\clearpage

\kfChapterCover{3}{프레게3 (챕터3)}{Chapter 3}{이번 챕터의 학습 내용을 시작합니다.}

\kfStartConcept{개념 영역 — 다섯 가지 감각으로 느끼는 심상}


\kfSecHeader{kfAreaConcept}{개념}{다섯 가지 감각으로 느끼는 심상}

시나 소설을 읽을 때 글자만 보는데도 머릿속에서 어떤 장면이 그림처럼 떠오르거나, 소리가 들리고, 냄새가 나는 것 같은 느낌을 받을 때가 있어요. 이처럼 글을 읽으면서 우리 마음속에 떠오르는 구체적이고 생생한 느낌을 심상 또는 이미지라고 해요. 작가들은 우리의 다섯 가지 감각인 시각, 청각, 촉각, 후각, 미각을 자극하는 표현을 사용해서 독자가 작품 속 상황을 직접 체험하는 것처럼 느끼게 만들어요.

시각적 심상은 눈으로 보는 것처럼 느끼게 하는 이미지예요. 색깔, 모양, 움직임 등을 표현해서 마치 한 폭의 그림을 보는 듯한 느낌을 줘요. 김종길의 「성탄제」에서 "어두운 방 안엔 바알간 숯불이 피고"라고 하면 어둠 속에서 빨갛게 타오르는 숯불의 모습이 눈앞에 선명하게 그려져요.

청각적 심상은 귀로 듣는 것처럼 느끼게 하는 이미지예요. 김소월의 「접동새」에서 "접동 접동 아우래비 접동"이라고 하면 접동새의 구슬픈 울음소리가 실제로 들리는 것 같아요.

촉각적 심상은 피부로 만지거나 온도를 느끼는 것처럼 하는 이미지예요. 김종길의 「성탄제」에서 "젊은 아버지의 서느런 옷자락에 열로 상기한 볼을 말없이 부비는 것이었다"라고 하면 차가운 옷자락과 뜨거운 볼의 감촉이 실제로 느껴지는 것 같아요.

후각적 심상은 코로 냄새를 맡는 것처럼 느끼게 하는 이미지예요. 김상옥의 「사향」에서 "어마씨 그리운 솜씨에 향그러운 꽃지짐"이라고 하면 어머니가 부쳐주신 전의 고소한 냄새가 코끝에 와 닿는 것 같아요.

미각적 심상은 혀로 맛을 보는 것처럼 느끼게 하는 이미지예요. 김상옥의 「사향」에서 "집집 끼니마다 봄을 씹고 사는 마을"이라고 하면 봄나물의 상큼한 맛을 느낄 수 있어요.

이처럼 다양한 심상을 사용하면 글이 훨씬 더 생생하고 재미있어져요. 작가들이 단어를 물감 삼아 우리 마음속에 그리는 그림이 바로 심상이에요.
\par\medskip\begin{itemize}[leftmargin=18pt, itemsep=2pt, topsep=2pt, label=\textbullet]
\item 시각 — '바알간 숯불이 피고'(「성탄제」)
\item 청각 — '접동 접동 아우래비 접동'(「접동새」)
\item 공감각(청각의 시각화) — '금으로 타는 태양의 즐거운 울림'
\end{itemize}

\kfSecHeader{kfAreaConcept}{문제}{}

\begin{kfQBlock}{1}{}
\kfQStem{시나 소설을 읽을 때, 언어를 통해 마음속에 떠오르는 구체적이고 감각적인 인상을 무엇이라고 하나요?}
\begin{kfChoices}
\kfChoice 주제
\kfChoice 심상
\kfChoice 배경
\kfChoice 갈등
\kfChoice 운율
\end{kfChoices}
\end{kfQBlock}
\begin{kfQBlock}{2}{}
\kfQStem{다음 중 시각적 심상이 가장 두드러지게 나타나는 구절은 무엇인가요?}
\begin{kfChoices}
\kfChoice 파아란 하늘빛 그리워
\kfChoice 발목이 시리도록 밟아도 보고
\kfChoice 머리맡에 찬물을 솨아 퍼붓고는
\kfChoice 접동 / 접동 / 아우래비 접동
\kfChoice 어마씨 그리운 솜씨에 향그러운 꽃지짐
\end{kfChoices}
\end{kfQBlock}
\begin{kfQBlock}{3}{}
\kfQStem{다음 중 청각적 심상이 가장 두드러지게 나타나는 구절은 무엇인가요?}
\begin{kfChoices}
\kfChoice 방범대원의 호각 소리, 메밀묵 사려 소리
\kfChoice 젊은 아버지의 서느런 옷자락에
\kfChoice 술 익는 마을마다 타는 저녁놀
\kfChoice 어두운 방 안에는 바알간 숯불이 피고
\kfChoice 물새알은 간간하고 짭조름한
\end{kfChoices}
\end{kfQBlock}
\begin{kfQBlock}{4}{}
\kfQStem{다음 중 후각적 심상이 가장 두드러지게 나타나는 구절은 무엇인가요?}
\begin{kfChoices}
\kfChoice 메마른 입술이 쓰디쓰다
\kfChoice 썩달라무 썩는 냄새 유달리 향그러웠다
\kfChoice 좁은 들길에 들장미 열매 붉어
\kfChoice 유리에 차고 슬픈 그 무엇이 어른거린다
\kfChoice 산산이 부서진 이름이여!
\end{kfChoices}
\end{kfQBlock}
\begin{kfQBlock}{5}{}
\kfQStem{다음 중 미각적 심상이 가장 두드러지게 나타나는 구절은 무엇인가요?}
\begin{kfChoices}
\kfChoice 끼니마다 봄을 씹고 사는 마을
\kfChoice 따가운 한낮의 햇살을 등에 지고
\kfChoice 뒷문 밖 정원에는 갈잎의 노래
\kfChoice 뜰에는 반짝반짝 빛나는 금모래빛
\kfChoice 점점 어두워져 가는 하늘 밑에
\end{kfChoices}
\end{kfQBlock}
\begin{kfQBlock}{6}{}
\kfQStem{다음 중 촉각적 심상이 가장 두드러지게 나타나는 구절은 무엇인가요?}
\begin{kfChoices}
\kfChoice 님의 침묵
\kfChoice 네 입술의 뜨거움
\kfChoice 내 마음은 호수요
\kfChoice 돌담에 속삭이는 햇발같이
\kfChoice 산은 자하산 봄눈 녹으면
\end{kfChoices}
\end{kfQBlock}
\begin{kfQBlock}{7}{}
\kfQStem{심상에 대한 설명으로 알맞지 \uline{않은} 것은 무엇인가요?}
\begin{kfChoices}
\kfChoice 눈에 보이는 시각적 표현만 심상으로 인정된다.
\kfChoice 독자가 작품을 오감으로 체험하게 한다.
\kfChoice 추상적인 감정을 구체적인 모습으로 느끼게 한다.
\kfChoice 시각, 청각, 후각, 미각, 촉각적 심상이 모두 있다.
\kfChoice '마음의 그림'이라고도 불린다.
\end{kfChoices}
\end{kfQBlock}
\begin{kfQBlock}{8}{}
\kfQStem{다음 중 후각적 심상이 가장 두드러지게 나타나는 구절은 무엇인가요?}
\begin{kfChoices}
\kfChoice 하늘에는 섣근 별
\kfChoice 발목이 시리도록 밟아도 보고
\kfChoice 어마씨 그리운 솜씨에 향그러운 꽃지짐
\kfChoice 길은 한 줄기 구겨진 넥타이처럼 풀어져
\kfChoice 방범대원의 호각 소리, 메밀묵 사려 소리
\end{kfChoices}
\end{kfQBlock}
\begin{kfQBlock}{9}{}
\kfQStem{다음 <보기>에 사용된 심상에 대한 설명으로 적절하지 \uline{않은} 것은 무엇인가요?}
\begin{kfEvidence}접동 / 접동 / 아우래비 접동\end{kfEvidence}
\begin{kfChoices}
\kfChoice 새의 울음소리를 흉내 낸 청각적 심상이다.
\kfChoice 음식의 단맛을 떠올리게 하는 미각적 심상이다.
\kfChoice 슬픔의 정서를 청각으로 환기한다.
\kfChoice 의성어를 활용하여 감각을 자극한다.
\kfChoice 같은 소리를 반복하여 운율을 만든다.
\end{kfChoices}
\end{kfQBlock}
\begin{kfQBlock}{10}{}
\kfQStem{다음 중 촉각적 심상이 사용되었다고 보기 \uline{어려운} 것은 무엇인가요?}
\begin{kfChoices}
\kfChoice 산산이 부서진 이름이여!
\kfChoice 젊은 아버지의 서느런 옷자락
\kfChoice 발목이 시리도록 밟아도 보고
\kfChoice 유리에 차고 슬픈 것이 어른거린다.
\kfChoice 내 볼에 와 닿던 네 입술의 뜨거움
\end{kfChoices}
\end{kfQBlock}
\begin{kfQBlock}{1}{}
\kfQStem{글을 읽으면서 마음속에 그림처럼 떠오르는 구체적이고 생생한 느낌을 무엇이라고 하나요?}
\kfAnswerLines{1}
\end{kfQBlock}
\begin{kfQBlock}{2}{}
\kfQStem{눈으로 보는 것처럼 색깔, 모양, 움직임을 느끼게 하는 이미지를 무슨 심상이라고 하나요?}
\kfAnswerLines{1}
\end{kfQBlock}
\begin{kfQBlock}{3}{}
\kfQStem{귀로 듣는 것처럼 소리를 느끼게 하는 이미지를 무슨 심상이라고 하나요?}
\kfAnswerLines{1}
\end{kfQBlock}
\begin{kfQBlock}{4}{}
\kfQStem{피부로 만지거나 온도를 느끼는 것처럼 하는 이미지를 무슨 심상이라고 하나요?}
\kfAnswerLines{1}
\end{kfQBlock}
\begin{kfQBlock}{5}{}
\kfQStem{코로 냄새를 맡는 것처럼 느끼게 하는 이미지를 무슨 심상이라고 하나요?}
\kfAnswerLines{1}
\end{kfQBlock}
\begin{kfQBlock}{6}{}
\kfQStem{혀로 맛을 보는 것처럼 느끼게 하는 이미지를 무슨 심상이라고 하나요?}
\kfAnswerLines{1}
\end{kfQBlock}
\begin{kfQBlock}{7}{}
\kfQStem{김종길의 「성탄제」에 나오는 "어두운 방 안엔 바알간 숯불이 피고"는 어떤 심상을 사용했나요?}
\kfAnswerLines{1}
\end{kfQBlock}
\begin{kfQBlock}{1}{}
\kfQStem{‘심상’이란 무엇이며, 시에서 심상을 사용했을 때의 효과는 무엇인지 서술하세요.}
\kfAnswerNote{답안}{4}
\end{kfQBlock}


\kfStartLit{문학 영역 — 작품}


\kfSecHeader{kfAreaLiterature}{지문}{작품}
\kfPassageNote{작품}{%
높은 가지를 흔드는 매미 소리에 묻혀 \\[4pt]
내 울음 아직은 노래가 아니다. \\[4pt]
차가운 바닥 위에 토해내는 울음, \\[4pt]
풀잎 없고 이슬 한 방울 내리지 않는 \\[4pt]
지하도 콘크리트 벽 좁은 틈에서 \\[4pt]
숨 막힐 듯, 그러나 나 여기 살아 있다 \\[4pt]
귀뚜르르 뚜르르 보내는 타전 소리가 \\[4pt]
누구의 마음 하나 울릴 수 있을까. \\[14pt]
지금은 매미들이 하늘을 찌르는 시절 \\[4pt]
그 소리 걷히고 맑은 가을이 \\[4pt]
어린 풀숲 위에 내려와 뒤척이기도 하고 \\[4pt]
계단을 타고 이 땅 밑까지 내려오는 날 \\[4pt]
발길에 눌려 우는 내 울음도 \\[4pt]
누군가의 가슴에 실려 가는 노래일 수 있을까. \\[14pt]
*타전: 전보나 무전을 침.
}


\kfSecHeader{kfAreaLiterature}{문제}{}

\begin{kfQBlock}{1}{}
\kfQStem{이 시에서 말하는 ‘나’는 누구를 의인화한 것인가요?}
\begin{kfChoices}
\kfChoice 귀뚜라미
\kfChoice 시인
\kfChoice 가을
\kfChoice 풀잎
\kfChoice 매미
\end{kfChoices}
\end{kfQBlock}
\begin{kfQBlock}{2}{}
\kfQStem{말하는 이가 현재 처해 있는 공간의 분위기로 가장 알맞은 것은 무엇인가요?}
\begin{kfChoices}
\kfChoice 높고 탁 트인 푸른 숲속
\kfChoice 따뜻하고 아늑한 보금자리
\kfChoice 화려하고 시끄러운 잔칫집
\kfChoice 춥고 삭막하며 비좁은 공간
\kfChoice 친구들과 함께 있는 즐거운 장소
\end{kfChoices}
\end{kfQBlock}
\begin{kfQBlock}{3}{}
\kfQStem{이 시에서 ‘매미’와 말하는 이의 관계를 바르게 설명한 것은 무엇인가요?}
\begin{kfChoices}
\kfChoice 서로 의지하며 다정하게 살아가는 친구 관계
\kfChoice 힘이 센 ‘매미’와 그에 억눌려 있는 ‘나’의 대립 관계
\kfChoice ‘매미’를 무조건 존경하고 따르는 스승과 제자 관계
\kfChoice 누가 더 노래를 잘하는지 다투는 경쟁의 라이벌 관계
\kfChoice 서로의 어려운 처지를 안타까워하는 연민의 관계
\end{kfChoices}
\end{kfQBlock}
\begin{kfQBlock}{4}{}
\kfQStem{말하는 이가 자신의 소리를 ‘울음’이라고 표현한 이유는 무엇인가요?}
\begin{kfChoices}
\kfChoice 노래를 부르는 방법을 아직 배우지 못해서
\kfChoice 자신의 처지가 너무나 고통스럽고 슬프기 때문에
\kfChoice 목소리가 너무 작아 다른 사람에게 들리지 않아서
\kfChoice 다른 사람의 마음을 울리고 싶다는 소망을 담아서
\kfChoice 아직 어리기 때문에 미숙한 소리를 낼 수밖에 없어서
\end{kfChoices}
\end{kfQBlock}
\begin{kfQBlock}{5}{}
\kfQStem{힘든 상황 속에서도 드러나는 말하는 이의 태도로 적절하지 \uline{않은} 것은 무엇인가요?}
\begin{kfChoices}
\kfChoice 자신의 삶을 포기하고 현실을 원망하는 체념적·부정적 태도
\kfChoice 시련을 견디며 자신의 자리를 지키려는 인내의 태도
\kfChoice 가을이라는 더 나은 시간을 꿈꾸는 긍정적 태도
\kfChoice 자신의 처지를 ‘울음’으로 솔직히 드러내는 진솔한 태도
\kfChoice 희망을 잃지 않고 미래를 기다리는 의지적인 태도
\end{kfChoices}
\end{kfQBlock}
\begin{kfQBlock}{6}{}
\kfQStem{이 시에서 말하는 이가 간절히 기다리는 ‘가을’은 어떤 시간인가요?}
\begin{kfChoices}
\kfChoice 자신의 존재를 알리는 것을 포기하는 시간
\kfChoice 모든 소리가 사라져 고요해지는 침묵의 시간
\kfChoice 자신의 진정한 목소리가 의미를 갖게 되는 희망의 시간
\kfChoice 날씨가 점점 더 추워져서 견디기 힘들고 몸과 마음이 모두 얼어붙는 시간
\kfChoice 매미들이 높은 가지 위에서 하늘을 찌를 듯이 더 힘차게 울어대는 시끄러운 시간
\end{kfChoices}
\end{kfQBlock}
\begin{kfQBlock}{7}{}
\kfQStem{“그 소리 걷히고”라는 표현은 시끄러운 소리가 사라지는 것을 눈에 보이는 것처럼 나타냈습니다. 이처럼 하나의 감각을 다른 감각으로 바꾸어 표현하는 방식의 효과는 무엇인가요?}
\begin{kfChoices}
\kfChoice 슬픈 내용을 재미있게 만들어 웃음을 유발한다.
\kfChoice 말하는 이의 불안한 심리를 감추는 역할을 한다.
\kfChoice 어려운 내용을 더 어렵게 만들어 비밀스럽게 한다.
\kfChoice 추상적인 느낌을 더 구체적이고 생생하게 느끼게 한다.
\kfChoice 독자가 직접 선택할 수 있도록 여러 의미를 제시한다.
\end{kfChoices}
\end{kfQBlock}
\begin{kfQBlock}{8}{}
\kfQStem{이 시의 전체 주제에 대한 설명으로 적절하지 \uline{않은} 것은 무엇인가요?}
\begin{kfChoices}
\kfChoice 짝사랑 상대가 자신의 마음을 알아주기를 간절히 바라는 사랑 노래
\kfChoice 힘든 현실 속에서도 가을(희망)을 기다리는 의지
\kfChoice ‘울음’이라는 말로 자기 처지를 솔직히 드러내는 자기 표현
\kfChoice 작은 존재의 노래로 세상과 소통하고자 하는 바람
\kfChoice 소외된 존재가 자신의 가치를 인정받고자 하는 소망
\end{kfChoices}
\end{kfQBlock}
\begin{kfQBlock}{1}{}
\kfQStem{말하는 이가 자신의 소리를 처음에는 ‘울음’이라고 했다가 나중에는 ‘노래’가 되기를 소망합니다. ‘울음’과 ‘노래’의 의미 차이를 <조건>에 맞게 서술하세요.}
\kfAnswerNote{답안}{4}
\end{kfQBlock}
\begin{kfQBlock}{2}{}
\kfQStem{말하는 이는 자신의 힘든 처지를 설명한 뒤, "그러나 나 여기 살아 있다"라고 말합니다. 이 말에 담긴 태도를 <조건>에 맞게 서술하세요.}
\kfAnswerNote{답안}{4}
\end{kfQBlock}

\kfSecHeader{kfAreaLiterature}{활동}{문장\_독해}
\noindent\textit{\small 다음 문장을 읽고 물음에 답하시오.}\par\medskip
\begin{kfSentenceUnit}{1}{}
\kfSentQ{1.}{다음 문장에서 서술어를 기준으로 각각의 주어는 무엇인가요?(1) 서술어 '찌르는'의 주어는 무엇인가요?}
\begin{kfChoices}
\kfChoice 내가
\kfChoice 매미들이
\kfChoice 그 소리가
\kfChoice 맑은 가을이
\end{kfChoices}
\kfSentQ{1.}{(2) 서술어 '걷히고'의 주어는 무엇인가요?}
\begin{kfChoices}
\kfChoice 내가
\kfChoice 매미들이
\kfChoice 그 소리가
\kfChoice 맑은 가을이
\end{kfChoices}
\kfSentQ{1.}{(3) 서술어 '내려와'의 주어는 무엇인가요?}
\begin{kfChoices}
\kfChoice 내가
\kfChoice 매미들이
\kfChoice 그 소리가
\kfChoice 맑은 가을이
\end{kfChoices}
\end{kfSentenceUnit}

\begin{kfSentenceUnit}{2}{}
\kfSentQ{2.}{다음 문장에서 '그 소리'가 가리키는 것은 무엇인가요?}
\begin{kfChoices}
\kfChoice 맑은 가을이 내려오는 소리
\kfChoice 매미들이 하늘을 찌르는 소리
\kfChoice 귀뚜라미가 보내는 타전 소리
\kfChoice 내 울음이 가슴에 실려 가는 소리
\end{kfChoices}
\end{kfSentenceUnit}

\begin{kfSentenceUnit}{3}{}
\kfSentQ{3.}{화자에게 '노래'의 의미는 무엇인가요?}
\begin{kfChoices}
\kfChoice 단순히 가을이 왔음을 알리는 계절의 신호이다.
\kfChoice 매미 소리에 묻혀 사라지는 의미 없는 소리이다.
\kfChoice 차가운 바닥에서 혼자 토해내는 고통의 외침이다.
\kfChoice 누군가의 마음을 울리고 가슴에 실려 가는 울음이다.
\end{kfChoices}
\end{kfSentenceUnit}


\kfStartNonlit{비문학 영역 — [과학] 제로 칼로리 식품의 비밀}


\kfSecHeader{kfAreaNonlit}{지문}{[과학] 제로 칼로리 식품의 비밀}
\kfPassageNote{[과학] 제로 칼로리 식품의 비밀}{%
살이 찌는 것을 걱정하지 않고 맛있는 음식을 먹으며 건강을 관리하려는 사람들이 늘어나고 있다. 그래서 최근 제로 칼로리 식품이 인기를 끌고 있다. 제로 칼로리 식품은 설탕 대신 인공으로 만든 감미료를 사용해서 단맛을 내면서도 열량은 낮춘 제품이다.

제로 칼로리라는 이름 때문에 실제 열량도 0이라고 생각할 수 있다. 하지만 우리나라 식품 관리 기관의 기준에 따르면, 식품 100밀리리터당 열량이 4칼로리보다 적으면 0칼로리로 표기할 수 있다. 따라서 제로 칼로리 식품의 실제 열량은 0칼로리가 아니라 매우 낮은 편이라고 보는 것이 맞다.

제로 칼로리 식품에 사용되는 인공 감미료는 음식에 단맛을 내기 위해 화학적으로 만든 재료이다. 이 재료는 설탕보다 훨씬 강한 단맛을 가지고 있다. 제로 칼로리 음료에 쓰이는 수크랄로스라는 재료의 단맛은 설탕보다 600배나 더 강하다.

이처럼 강한 단맛인데도 인공 감미료를 사용한 제품의 열량이 낮은 이유는 무엇일까? 에너지를 만들기 위해 당분으로 분해되어 우리 몸에 흡수되는 설탕과 달리, 인공 감미료는 우리 몸에서 에너지로 바뀌지 않는다. 또한 다른 형태로 변하지 않은 상태로 몸 밖으로 나간다. 수크랄로스의 경우, 우리 몸이 당분으로 알아보지 못해서 85퍼센트는 위와 장에서 흡수되지 않은 채 배출된다. 일부 흡수된 양도 소변을 통해 빠르게 나간다. 이와 같은 인공 감미료의 특성 때문에 제로 칼로리 식품을 만들 수 있는 것이다.

그렇다면 인공 감미료로 단맛을 낸 제로 칼로리 식품은 마음 놓고 먹어도 좋은 것일까? 전문가들은 열량 부담 없이 단맛을 즐길 수 있는 인공 감미료의 장점이 우리 몸에 나쁜 영향을 줄 수도 있다고 경고한다. 우리 몸은 단맛을 에너지가 들어오는 것으로 여겨 왔기 때문에 단맛과 열량이 맞지 않으면 몸의 활동에 혼란을 줄 수 있다. 또한 인공 감미료가 들어간 제로 칼로리 식품을 자주 먹으면 단맛에 대한 감각이 둔해져서 더 강한 단맛을 찾는 단맛 중독 상태가 될 수 있다.

제로 칼로리 식품은 열량 섭취를 줄이고 싶거나 혈액 속 당분을 관리해야 하는 사람들에게 하나의 방법이 될 수도 있다. 하지만 인공 감미료가 들어간 제로 칼로리 식품을 지나치게 먹을 경우 우리 몸에 나쁜 영향이 있을 수 있으므로 주의해야 한다. 인공 감미료가 건강에 미치는 영향은 개인의 건강 상태와 먹는 양에 따라 다르므로 이 점을 생각하여 제로 칼로리 식품을 이용하는 지혜가 필요하다.
}


\kfSecHeader{kfAreaNonlit}{문제}{}

\begin{kfQBlock}{1}{}
\kfQStem{윗글의 중심 내용으로 가장 적절한 것은 무엇인가요?}
\begin{kfChoices}
\kfChoice 제로 칼로리 식품의 의미와 장단점
\kfChoice 인공 감미료를 만드는 화학적 과정
\kfChoice 설탕이 우리 몸에 미치는 긍정적 영향
\kfChoice 단맛 중독을 해결하기 위한 구체적인 방법
\kfChoice 올바른 식습관을 통한 건강 관리의 중요성
\end{kfChoices}
\end{kfQBlock}
\begin{kfQBlock}{2}{}
\kfQStem{윗글의 내용과 일치하지 않는 것은 무엇인가요?}
\begin{kfChoices}
\kfChoice 제로 칼로리 식품도 실제로는 약간의 열량을 포함할 수 있다.
\kfChoice 전문가들은 제로 칼로리 식품이 건강에 아무런 해가 없다고 말한다.
\kfChoice 인공 감미료는 설탕보다 적은 양으로도 강한 단맛을 낼 수 있다.
\kfChoice 수크랄로스는 대부분 우리 몸에 흡수되지 않고 몸 밖으로 나간다.
\kfChoice 우리 몸은 단맛을 느끼면 자연스럽게 에너지가 들어온다고 여긴다.
\end{kfChoices}
\end{kfQBlock}
\begin{kfQBlock}{3}{}
\kfQStem{우리나라에서 식품에 '0칼로리'라고 표시하기 위한 기준에 대한 설명으로 적절하지 \uline{않은} 것은 무엇인가요?}
\begin{kfChoices}
\kfChoice 식품 100ml당 4칼로리보다 적은 음료에는 ‘0칼로리’ 표시를 할 수 있다.
\kfChoice 음료의 경우 100ml 기준으로 열량 기준이 정해져 있다.
\kfChoice 식품 100g당 4칼로리 미만이면 ‘0칼로리’로 표시할 수 있다.
\kfChoice 열량이 정확히 0칼로리여야만 ‘0칼로리’ 표시를 허용한다.
\kfChoice 식약처가 정한 기준에 따라 ‘0칼로리’ 표시 여부가 결정된다.
\end{kfChoices}
\end{kfQBlock}
\begin{kfQBlock}{4}{}
\kfQStem{윗글에서 설명하는 '제로 칼로리 식품'에 대한 설명으로 적절하지 \uline{않은} 것은 무엇인가요?}
\begin{kfChoices}
\kfChoice 비타민과 미네랄이 풍부해 의약품으로도 분류되는 건강 식품이다.
\kfChoice ‘0칼로리’ 표시 기준 안에 들어가는 저열량 음료·식품이다.
\kfChoice 설탕 대신 인공 감미료를 사용하므로 단맛은 충분히 느껴진다.
\kfChoice 식약처 기준상 100ml당 4칼로리 미만이면 ‘0칼로리’ 표기가 가능하다.
\kfChoice 인공 감미료로 단맛을 내지만 실제 열량은 매우 낮은 식품이다.
\end{kfChoices}
\end{kfQBlock}
\begin{kfQBlock}{5}{}
\kfQStem{다음 <보기>의 상황을 보고, 민지가 마신 음료수의 실제 열량에 대해 가장 바르게 추측한 것은 무엇인가요?}
\begin{kfEvidence}민지는 편의점에서 200ml(밀리리터)짜리 음료수를 샀습니다. 음료수 병에는 크게 '0칼로리(Zero Calorie)'라고 적혀 있었습니다. 민지는 "와, 이건 열량이 아예 없구나!"라고 생각했습니다.\end{kfEvidence}
\begin{kfChoices}
\kfChoice 0칼로리다.
\kfChoice 10칼로리다.
\kfChoice 8칼로리 미만이다.
\kfChoice 살이 많이 찌는 물이다.
\kfChoice 설탕이 아주 많이 들어있다.
\end{kfChoices}
\end{kfQBlock}
\begin{kfQBlock}{6}{}
\kfQStem{다음 <보기>의 철수에게 일어난 현상을 설명하는 말로 가장 알맞은 것은 무엇인가요?}
\begin{kfEvidence}철수: "나 요즘 다이어트하려고 제로 콜라만 마시고 있어." 영희: "그래? 그럼 단걸 덜 먹게 되니?" 철수: "아니, 이상하게 귤이나 사과를 먹으면 예전만큼 달게 느껴지지가 않아. 그래서 자꾸 더 달콤한 사탕이나 초콜릿이 먹고 싶어져."\end{kfEvidence}
\begin{kfChoices}
\kfChoice 배가 부른 상태다.
\kfChoice 건강이 아주 좋아졌다.
\kfChoice 단맛 중독 상태가 되었다.
\kfChoice 혀의 감각이 매우 예민해졌다.
\kfChoice 에너지가 넘쳐서 운동이 필요하다.
\end{kfChoices}
\end{kfQBlock}
\begin{kfQBlock}{7}{}
\kfQStem{윗글을 읽은 후, <보기>의 질문에 해 줄 수 있는 조언으로 가장 적절한 것은 무엇인가요?}
\begin{kfEvidence}질문: "저는 단 음료수를 너무 좋아해서 살이 많이 쪘어요. 앞으로는 일반 음료수 대신 제로 칼로리 음료수만 하루에 10병씩 마시려고 해요. 그럼 건강해질까요?"\end{kfEvidence}
\begin{kfChoices}
\kfChoice 물 대신 마셔도 된다.
\kfChoice 무조건 많이 먹는 게 좋다.
\kfChoice 건강 상태에 따라 조절해야 한다.
\kfChoice 다이어트에 도움이 되니 마음껏 먹어라.
\kfChoice 뇌의 활동을 도와주니 공부할 때 필수다.
\end{kfChoices}
\end{kfQBlock}
\begin{kfQBlock}{8}{}
\kfQStem{전문가들이 제로 칼로리 식품을 먹을 때 주의하라고 경고하는 이유로 적절하지 \uline{않은} 것은 무엇인가요?}
\begin{kfChoices}
\kfChoice 단맛과 실제 열량의 불일치로 뇌가 혼란을 느낄 수 있기 때문이다.
\kfChoice 단맛이 너무 약해 식욕을 잃게 만들기 때문이다.
\kfChoice 인공 감미료의 안전성이 완벽하게 입증된 것은 아니기 때문이다.
\kfChoice ‘0칼로리’라는 표시 때문에 마음 놓고 과다 섭취할 수 있기 때문이다.
\kfChoice 단맛에 익숙해지면 더 단 음식을 찾게 되어 결국 살이 찔 수 있기 때문이다.
\end{kfChoices}
\end{kfQBlock}
\begin{kfQBlock}{1}{}
\kfQStem{우리나라 기준상 '0칼로리'로 표기하려면, 식품 100ml당 열량이 몇 칼로리보다 적어야 하나요?}
\kfAnswerLines{1}
\end{kfQBlock}
\begin{kfQBlock}{2}{}
\kfQStem{지문에 예시로 나온 '수크랄로스'는 설탕보다 몇 배 더 강한 단맛을 가지고 있나요?}
\kfAnswerLines{1}
\end{kfQBlock}
\begin{kfQBlock}{3}{}
\kfQStem{인공 감미료가 열량이 낮은 이유는 우리 몸에서 무엇으로 바뀌지 않기 때문인가요?}
\kfAnswerLines{1}
\end{kfQBlock}
\begin{kfQBlock}{4}{}
\kfQStem{우리 몸이 수크랄로스를 대부분 흡수하지 못하는 이유는 무엇인가요?}
\kfAnswerLines{1}
\end{kfQBlock}
\begin{kfQBlock}{5}{}
\kfQStem{제로 칼로리 식품을 자주 먹으면 단맛에 대한 감각이 둔해져서 더 강한 단맛을 찾게 되는 상태를 무엇이라고 하나요?}
\kfAnswerLines{1}
\end{kfQBlock}
\begin{kfQBlock}{1}{}
\kfQStem{제로 칼로리 식품을 긍정적으로 보는 입장과 부정적으로 보는 입장의 근거를 지문에서 각각 찾아 서술하세요.}
\kfAnswerNote{답안}{4}
\end{kfQBlock}

\kfSecHeader{kfAreaNonlit}{활동}{문장\_독해}
\noindent\textit{\small 다음 문장을 읽고 물음에 답하시오.}\par\medskip
\begin{kfSentenceUnit}{1}{}
\kfSentQ{1.}{최근 어떤 사람들이 늘어나고 있나요?}
\begin{kfChoices}
\kfChoice 살이 찌는 것을 걱정하지 않고 맛있는 음식을 먹으며 건강을 관리하려는 사람들.
\kfChoice 운동은 하기 싫어하고 음식 섭취량만 극단적으로 줄여서 쉽고 빠르게 살을 빼려는 사람들.
\kfChoice 주머니 사정이 가벼워서 맛이나 영양보다는 가격이 싸고 양이 많은 음식을 선호하는 사람들.
\kfChoice 건강을 위해서라면 맛있는 음식을 꾹 참고 오직 몸에 좋은 유기농 채소만 찾아 먹는 사람들.
\end{kfChoices}
\end{kfSentenceUnit}

\begin{kfSentenceUnit}{2}{}
\kfSentQ{2.}{이 문장에 나타난, 현대인들이 바라는 것이 \uline{아닌} 것은 무엇인가요?}
\begin{kfChoices}
\kfChoice 살이 찌지 않는 것.
\kfChoice 건강을 관리하는 것.
\kfChoice 돈을 많이 버는 것.
\kfChoice 맛있는 음식을 먹는 것.
\end{kfChoices}
\end{kfSentenceUnit}

\begin{kfSentenceUnit}{3}{}
\kfSentQ{3.}{'그래서'의 구체적인 내용은 무엇인가요?}
\begin{kfChoices}
\kfChoice 살이 찌는 것을 걱정하지 않고 맛있는 음식을 먹으며 건강을 관리하려는 사람들이 늘어나고 있어서.
\kfChoice 해외에서 큰 인기를 끈 다양한 종류의 다이어트 보조 식품들이 국내 시장에 본격적으로 수입되고 있어서.
\kfChoice 경제적인 이유로 가격이 저렴하면서도 양이 푸짐한 가성비 좋은 식품을 찾는 소비자들이 늘어나고 있어서.
\kfChoice 건강을 챙기기보다는 자극적이고 맛있는 음식을 먹는 즐거움을 더 중요하게 생각하는 사람들이 늘어나고 있어서.
\end{kfChoices}
\end{kfSentenceUnit}

\begin{kfSentenceUnit}{4}{}
\kfSentQ{4.}{살이 찌는 것을 걱정하지 않고 맛있는 음식을 먹으며 건강을 관리하려는 사람들이 늘어나고 있어서 무엇이 인기를 끌고 있나요?}
\begin{kfChoices}
\kfChoice 발효 식품.
\kfChoice 유기농 채소.
\kfChoice 고단백질 식품.
\kfChoice 제로 칼로리 식품.
\end{kfChoices}
\end{kfSentenceUnit}

\begin{kfSentenceUnit}{5}{}
\kfSentQ{5.}{제로 칼로리 식품은 설탕 대신 무엇을 사용하나요?}
\begin{kfChoices}
\kfChoice 유기농 설탕.
\kfChoice 천연 과일 농축액.
\kfChoice 인공으로 만든 감미료.
\kfChoice 전통적인 방식의 조미료.
\end{kfChoices}
\end{kfSentenceUnit}

\begin{kfSentenceUnit}{6}{}
\kfSentQ{6.}{감미료는 어떻게 만든 것인가요?}
\begin{kfChoices}
\kfChoice 인공으로 만든 것.
\kfChoice 자연에서 얻은 것.
\kfChoice 동물의 뼈를 고아서 만든 것.
\kfChoice 전통적인 방법으로 발효시킨 것.
\end{kfChoices}
\end{kfSentenceUnit}

\begin{kfSentenceUnit}{7}{}
\kfSentQ{7.}{감미료를 사용하는 목적은 무엇인가요?}
\begin{kfChoices}
\kfChoice 단맛을 내기 위해서.
\kfChoice 음식의 짠맛을 줄이기 위해서.
\kfChoice 음식을 오랫동안 보관하기 위해서.
\kfChoice 음식의 색깔을 예쁘게 만들기 위해서.
\end{kfChoices}
\end{kfSentenceUnit}

\begin{kfSentenceUnit}{8}{}
\kfSentQ{8.}{제로 칼로리 식품은 열량을 어떻게 했나요?}
\begin{kfChoices}
\kfChoice 낮췄다.
\kfChoice 높였다.
\kfChoice 없앴다.
\kfChoice 유지했다.
\end{kfChoices}
\end{kfSentenceUnit}

\begin{kfSentenceUnit}{9}{}
\kfSentQ{9.}{'제로 칼로리 식품'을 정의하는 세 가지 특징은 무엇인가요?}
\begin{kfChoices}
\kfChoice 천연 재료 사용, 짠맛을 냄, 열량은 높임.
\kfChoice 전통 방식 제조, 신맛을 냄, 열량은 약간 낮춤.
\kfChoice 화학 조미료 사용, 매운맛을 냄, 열량은 그대로임.
\kfChoice 설탕 대신 인공 감미료 사용, 단맛을 냄, 열량은 낮춤.
\end{kfChoices}
\end{kfSentenceUnit}

\begin{kfSentenceUnit}{10}{}
\kfSentQ{10.}{'제로 칼로리'라는 이름이 사람들에게 어떤 오해를 불러일으킬 수 있나요?}
\begin{kfChoices}
\kfChoice 맛이 없을 것이라는 오해.
\kfChoice 가격이 매우 비쌀 것이라는 오해.
\kfChoice 건강에 무조건 좋을 것이라는 오해.
\kfChoice 실제 열량이 0칼로리일 것이라는 오해.
\end{kfChoices}
\end{kfSentenceUnit}

\begin{kfSentenceUnit}{11}{}
\kfSentQ{11.}{사람들이 무엇 때문에 잘못 생각할 수 있나요?}
\begin{kfChoices}
\kfChoice 가격이 저렴해서 잘못 생각할 수 있다.
\kfChoice 포장이 화려해서 잘못 생각할 수 있다.
\kfChoice 광고 모델이 유명해서 잘못 생각할 수 있다.
\kfChoice 제로 칼로리라는 이름 때문에 잘못 생각할 수 있다.
\end{kfChoices}
\end{kfSentenceUnit}

\begin{kfSentenceUnit}{12}{}
\kfSentQ{12.}{사람들은 제로 칼로리의 실제 열량이 얼마라고 생각할 수 있나요?}
\begin{kfChoices}
\kfChoice 0.
\kfChoice 10.
\kfChoice 50.
\kfChoice 100.
\end{kfChoices}
\end{kfSentenceUnit}

\begin{kfSentenceUnit}{13}{}
\kfSentQ{13.}{무엇의 기준에 따른 내용인가요?}
\begin{kfChoices}
\kfChoice 학교 급식 영양사.
\kfChoice 동네 슈퍼마켓 주인.
\kfChoice 유명한 요리 연구가.
\kfChoice 우리나라 식품 관리 기관.
\end{kfChoices}
\end{kfSentenceUnit}

\begin{kfSentenceUnit}{14}{}
\kfSentQ{14.}{우리나라에서 '0칼로리'로 표기하기 위한 조건은 무엇인가요?}
\begin{kfChoices}
\kfChoice 식품 100밀리리터당 열량이 0칼로리여야 한다.
\kfChoice 식품 100밀리리터당 열량이 4칼로리보다 적어야 한다.
\kfChoice 식품 100밀리리터당 열량이 50칼로리보다 적어야 한다.
\kfChoice 식품 100밀리리터당 열량이 10칼로리보다 적어야 한다.
\end{kfChoices}
\end{kfSentenceUnit}

\begin{kfSentenceUnit}{15}{}
\kfSentQ{15.}{'따라서'라는 말을 볼 때, 이 문장은 앞선 '0칼로리 표기 기준'으로부터 어떤 결론을 내리고 있나요?}
\begin{kfChoices}
\kfChoice 제로 칼로리 식품은 맛이 없을 수도 있다는 결론.
\kfChoice 제로 칼로리 식품은 가격이 비쌀 수도 있다는 결론.
\kfChoice 제로 칼로리 식품은 건강에 해로울 수도 있다는 결론.
\kfChoice 제로 칼로리 식품의 열량은 0이 아닐 수도 있다는 결론.
\end{kfChoices}
\end{kfSentenceUnit}

\begin{kfSentenceUnit}{16}{}
\kfSentQ{16.}{인공 감미료는 어떻게 만든 재료인가요?}
\begin{kfChoices}
\kfChoice 화학적으로 만든 재료.
\kfChoice 동물성 단백질로 만든 재료.
\kfChoice 자연에서 추출한 천연 재료.
\kfChoice 전통적인 방식으로 발효시킨 재료.
\end{kfChoices}
\end{kfSentenceUnit}

\begin{kfSentenceUnit}{17}{}
\kfSentQ{17.}{인공 감미료를 만드는 목적은 무엇인가요?}
\begin{kfChoices}
\kfChoice 음식에 단맛을 내기 위해.
\kfChoice 음식의 색깔을 내기 위해.
\kfChoice 음식의 향을 좋게 하기 위해.
\kfChoice 음식의 유통 기한을 늘리기 위해.
\end{kfChoices}
\end{kfSentenceUnit}

\begin{kfSentenceUnit}{18}{}
\kfSentQ{18.}{인공 감미료는 어디에 사용되나요?}
\begin{kfChoices}
\kfChoice 전통 발효 식품.
\kfChoice 고단백질 영양제.
\kfChoice 제로 칼로리 식품.
\kfChoice 유기농 채소 샐러드.
\end{kfChoices}
\end{kfSentenceUnit}

\begin{kfSentenceUnit}{19}{}
\kfSentQ{19.}{'이 재료'가 가리키는 것은 무엇인가요?}
\begin{kfChoices}
\kfChoice 발효 식초.
\kfChoice 천연 조미료.
\kfChoice 유기농 설탕.
\kfChoice 인공 감미료.
\end{kfChoices}
\end{kfSentenceUnit}

\begin{kfSentenceUnit}{20}{}
\kfSentQ{20.}{인공 감미료는 설탕과 비교했을 때 어떤 특징이 있나요?}
\begin{kfChoices}
\kfChoice 설탕보다 신맛이 강한 특징이 있다.
\kfChoice 설탕보다 단맛이 훨씬 강한 특징이 있다.
\kfChoice 설탕보다 단맛이 훨씬 약한 특징이 있다.
\kfChoice 설탕과 단맛은 비슷하지만 짠맛이 나는 특징이 있다.
\end{kfChoices}
\end{kfSentenceUnit}

\begin{kfSentenceUnit}{21}{}
\kfSentQ{21.}{수크랄로스는 어디에 쓰이나요?}
\begin{kfChoices}
\kfChoice 전통 차.
\kfChoice 유기농 우유.
\kfChoice 천연 과일 주스.
\kfChoice 제로 칼로리 음료.
\end{kfChoices}
\end{kfSentenceUnit}

\begin{kfSentenceUnit}{22}{}
\kfSentQ{22.}{수크랄로스의 단맛은 설탕보다 몇 배나 강한가요?}
\begin{kfChoices}
\kfChoice 설탕의 600배이다.
\kfChoice 설탕의 50배 수준이다.
\kfChoice 설탕의 2배에 불과하다.
\kfChoice 설탕의 10배 정도이다.
\end{kfChoices}
\end{kfSentenceUnit}

\begin{kfSentenceUnit}{23}{}
\kfSentQ{23.}{'이처럼'이 가리키는 것은 무엇인가요?}
\begin{kfChoices}
\kfChoice 수크랄로스가 설탕보다 가격이 600배나 비싼 것.
\kfChoice 수크랄로스가 설탕보다 인기가 600배나 많은 것.
\kfChoice 수크랄로스가 설탕보다 생산량이 600배나 많은 것.
\kfChoice 수크랄로스가 설탕보다 600배나 더 강한 단맛을 가지는 것.
\end{kfChoices}
\end{kfSentenceUnit}

\begin{kfSentenceUnit}{24}{}
\kfSentQ{24.}{이 문장은 글의 흐름에서 어떤 역할을 하나요?}
\begin{kfChoices}
\kfChoice 이미 알려진 사실을 다시 한번 강조하는 역할.
\kfChoice 전문가의 의견을 인용하여 신뢰성을 높이는 역할.
\kfChoice 글의 결론을 미리 제시하여 독자를 설득하는 역할.
\kfChoice 독자의 궁금증을 유발하고, 앞으로 이어질 내용을 안내하는 역할.
\end{kfChoices}
\end{kfSentenceUnit}

\begin{kfSentenceUnit}{25}{}
\kfSentQ{25.}{설탕은 우리 몸에서 무엇으로 분해되나요?}
\begin{kfChoices}
\kfChoice 지방.
\kfChoice 당분.
\kfChoice 비타민.
\kfChoice 단백질.
\end{kfChoices}
\end{kfSentenceUnit}

\begin{kfSentenceUnit}{26}{}
\kfSentQ{26.}{설탕과 인공 감미료의 차이점은 무엇인가요?}
\begin{kfChoices}
\kfChoice 설탕은 짠맛을 내지만, 인공 감미료는 단맛을 낸다.
\kfChoice 설탕은 차갑게 먹지만, 인공 감미료는 뜨겁게 먹는다.
\kfChoice 설탕은 공장에서 만들지만, 인공 감미료는 자연에서 얻는다.
\kfChoice 설탕은 에너지로 바뀌지만, 인공 감미료는 에너지로 바뀌지 않는다.
\end{kfChoices}
\end{kfSentenceUnit}

\begin{kfSentenceUnit}{27}{}
\kfSentQ{27.}{'또한'이 나타내는 관계는 무엇인가요?}
\begin{kfChoices}
\kfChoice 앞 내용과 반대되는 내용.
\kfChoice 앞 내용에 추가되는 내용.
\kfChoice 앞 내용의 원인이 되는 내용.
\kfChoice 앞 내용의 예시가 되는 내용.
\end{kfChoices}
\end{kfSentenceUnit}

\begin{kfSentenceUnit}{28}{}
\kfSentQ{28.}{무엇이 몸 밖으로 나가나요?}
\begin{kfChoices}
\kfChoice 지방.
\kfChoice 설탕.
\kfChoice 단백질.
\kfChoice 인공 감미료.
\end{kfChoices}
\end{kfSentenceUnit}

\begin{kfSentenceUnit}{29}{}
\kfSentQ{29.}{인공 감미료는 어떤 상태로 몸 밖으로 나가나요?}
\begin{kfChoices}
\kfChoice 에너지로 변한 상태.
\kfChoice 지방으로 축적된 상태.
\kfChoice 근육으로 흡수된 상태.
\kfChoice 다른 형태로 변하지 않은 상태.
\end{kfChoices}
\end{kfSentenceUnit}

\begin{kfSentenceUnit}{30}{}
\kfSentQ{30.}{우리 몸은 수크랄로스를 무엇으로 알아보지 못하나요?}
\begin{kfChoices}
\kfChoice 당분.
\kfChoice 지방.
\kfChoice 단백질.
\kfChoice 비타민.
\end{kfChoices}
\end{kfSentenceUnit}

\begin{kfSentenceUnit}{31}{}
\kfSentQ{31.}{무엇이 위와 장에서 흡수되지 않은 채 배출되나요?}
\begin{kfChoices}
\kfChoice 설탕.
\kfChoice 포도당.
\kfChoice 비타민.
\kfChoice 수크랄로스.
\end{kfChoices}
\end{kfSentenceUnit}

\begin{kfSentenceUnit}{32}{}
\kfSentQ{32.}{우리 몸이 수크랄로스를 흡수하지 못하는 이유는 무엇인가요?}
\begin{kfChoices}
\kfChoice 우리 몸이 수크랄로스를 독성 물질로 인식하기 때문이다.
\kfChoice 우리 몸이 수크랄로스를 당분으로 알아보지 못하기 때문이다.
\kfChoice 우리 몸이 수크랄로스를 이미 충분히 가지고 있기 때문이다.
\kfChoice 우리 몸에 수크랄로스를 분해하는 효소가 너무 많기 때문이다.
\end{kfChoices}
\end{kfSentenceUnit}

\begin{kfSentenceUnit}{33}{}
\kfSentQ{33.}{우리 몸에 흡수된 일부 수크랄로스는 어떻게 되나요?}
\begin{kfChoices}
\kfChoice 뇌에 에너지를 공급한다.
\kfChoice 근육을 만드는 데 사용된다.
\kfChoice 지방으로 바뀌어 몸속에 쌓인다.
\kfChoice 소변을 통해 빠르게 몸 밖으로 나간다.
\end{kfChoices}
\end{kfSentenceUnit}

\begin{kfSentenceUnit}{34}{}
\kfSentQ{34.}{'이와 같은 인공 감미료의 특성'이란 구체적으로 무엇을 말하나요?}
\begin{kfChoices}
\kfChoice 소화 기관에서 빠르게 흡수되어 혈액 속의 당분 수치를 급격하게 높이는 특성.
\kfChoice 위와 장에서 빠르게 분해되어 우리 몸에 즉각적인 에너지를 공급해 주는 특성.
\kfChoice 체내에 흡수되어 지방으로 축적됨으로써 체중 증가의 직접적인 원인이 되는 특성.
\kfChoice 몸에서 에너지로 바뀌지 않고, 대부분 흡수되지 않은 채 몸 밖으로 나가는 특성.
\end{kfChoices}
\end{kfSentenceUnit}

\begin{kfSentenceUnit}{35}{}
\kfSentQ{35.}{인공 감미료가 에너지로 바뀌지 않고 몸 밖으로 배출되는 특성 때문에 무엇을 만들 수 있나요?}
\begin{kfChoices}
\kfChoice 고지방 식품.
\kfChoice 고칼슘 식품.
\kfChoice 고단백질 식품.
\kfChoice 제로 칼로리 식품.
\end{kfChoices}
\end{kfSentenceUnit}

\begin{kfSentenceUnit}{36}{}
\kfSentQ{36.}{'그렇다면'이 나타내는 관계는 무엇인가요?}
\begin{kfChoices}
\kfChoice 앞 내용을 요약하는 질문.
\kfChoice 앞 내용을 반박하는 질문.
\kfChoice 앞 내용의 예시를 묻는 질문.
\kfChoice 앞 내용을 바탕으로 한 새로운 질문.
\end{kfChoices}
\end{kfSentenceUnit}

\begin{kfSentenceUnit}{37}{}
\kfSentQ{37.}{이 질문은 앞으로 나올 내용을 안내하는 역할을 합니다. 이 질문을 통해 앞으로 무엇에 대해 이야기할 것인지 알 수 있나요?}
\begin{kfChoices}
\kfChoice 제로 칼로리 식품이 처음 개발된 역사적 배경과 다양한 종류에 대해 설명할 것.
\kfChoice 제로 칼로리 식품이 시중에서 판매되는 가격과 소비자들의 반응에 대해 설명할 것.
\kfChoice 제로 칼로리 식품의 긍정적인 면이 아닌, 부정적인 면이나 문제점에 대해 설명할 것.
\kfChoice 제로 칼로리 식품을 만드는 데 사용되는 인공 감미료의 제조 과정에 대해 설명할 것.
\end{kfChoices}
\end{kfSentenceUnit}

\begin{kfSentenceUnit}{38}{}
\kfSentQ{38.}{누가 경고하나요?}
\begin{kfChoices}
\kfChoice 요리사.
\kfChoice 전문가.
\kfChoice 정치인.
\kfChoice 운동선수.
\end{kfChoices}
\end{kfSentenceUnit}

\begin{kfSentenceUnit}{39}{}
\kfSentQ{39.}{인공 감미료의 장점은 무엇인가요?}
\begin{kfChoices}
\kfChoice 포장이 예뻐서 선물하기 좋은 것.
\kfChoice 열량 부담 없이 단맛을 즐길 수 있는 것.
\kfChoice 가격이 저렴해서 누구나 쉽게 살 수 있는 것.
\kfChoice 유통 기한이 길어서 오래 보관할 수 있는 것.
\end{kfChoices}
\end{kfSentenceUnit}

\begin{kfSentenceUnit}{40}{}
\kfSentQ{40.}{전문가들이 경고하는 내용은 무엇인가요?}
\begin{kfChoices}
\kfChoice 인공 감미료는 맛이 없어서 사람들이 싫어한다는 것.
\kfChoice 인공 감미료는 가격이 너무 비싸서 경제적이지 않다는 것.
\kfChoice 인공 감미료는 만들기가 너무 어려워서 생산량이 적다는 것.
\kfChoice 인공 감미료의 장점이 오히려 몸에 나쁜 영향을 줄 수 있다는 것.
\end{kfChoices}
\end{kfSentenceUnit}

\begin{kfSentenceUnit}{41}{}
\kfSentQ{41.}{우리 몸은 단맛을 무엇으로 여겨 왔나요?}
\begin{kfChoices}
\kfChoice 잠을 자야 하는 신호.
\kfChoice 에너지가 들어오는 것.
\kfChoice 운동을 해야 하는 신호.
\kfChoice 물을 마셔야 하는 신호.
\end{kfChoices}
\end{kfSentenceUnit}

\begin{kfSentenceUnit}{42}{}
\kfSentQ{42.}{우리 몸에 혼란을 주는 조건은 무엇과 무엇이 맞지 않는 경우인가요?}
\begin{kfChoices}
\kfChoice 맛과 냄새.
\kfChoice 가격과 품질.
\kfChoice 단맛과 열량.
\kfChoice 디자인과 기능.
\end{kfChoices}
\end{kfSentenceUnit}

\begin{kfSentenceUnit}{43}{}
\kfSentQ{43.}{인공 감미료가 우리 몸의 활동에 혼란을 줄 수 있는 이유는 무엇인가요?}
\begin{kfChoices}
\kfChoice 우리 몸은 자연적인 맛을 좋아하는데, 인공 감미료는 화학적인 맛을 내어 미각을 둔하게 만들기 때문이다.
\kfChoice 우리 몸은 단맛을 에너지로 여기는데, 인공 감미료는 단맛만 있고 열량이 없어서 몸이 혼란스러워하기 때문이다.
\kfChoice 우리 몸은 에너지를 필요로 하는데, 인공 감미료는 에너지를 과도하게 공급하여 호르몬 균형을 깨뜨리기 때문이다.
\kfChoice 우리 몸은 단맛을 느끼면 소화 효소를 분비하는데, 인공 감미료는 소화가 되지 않아 위장에 큰 부담을 주기 때문이다.
\end{kfChoices}
\end{kfSentenceUnit}

\begin{kfSentenceUnit}{44}{}
\kfSentQ{44.}{어떤 경우에 단맛 중독 상태가 될 수 있나요?}
\begin{kfChoices}
\kfChoice 규칙적인 식습관을 유지할 경우.
\kfChoice 천연 재료로 만든 음식만 먹을 경우.
\kfChoice 인공 감미료가 들어간 식품을 전혀 먹지 않을 경우.
\kfChoice 인공 감미료가 들어간 제로 칼로리 식품을 자주 먹을 경우.
\end{kfChoices}
\end{kfSentenceUnit}

\begin{kfSentenceUnit}{45}{}
\kfSentQ{45.}{'단맛 중독'이 발생하는 과정은 어떠한가요?}
\begin{kfChoices}
\kfChoice 1단계
\kfChoice 원인, 결과
\kfChoice 초기, 말기
\kfChoice 1단계, 2단계
\end{kfChoices}
\end{kfSentenceUnit}

\begin{kfSentenceUnit}{46}{}
\kfSentQ{46.}{제로 칼로리 식품이 도움이 될 수 있는 사람들은 어떤 사람들인가요?}
\begin{kfChoices}
\kfChoice 소화 기능이 약한 노인들.
\kfChoice 성장기 어린이나 청소년들.
\kfChoice 운동을 전문적으로 하는 선수들.
\kfChoice 열량 섭취를 줄이고 싶거나 혈당 관리가 필요한 사람들.
\end{kfChoices}
\end{kfSentenceUnit}

\begin{kfSentenceUnit}{47}{}
\kfSentQ{47.}{이 문장에서 제로 칼로리 식품의 긍정적인 면은 무엇인가요?}
\begin{kfChoices}
\kfChoice 근육을 키우고 체력을 증진하는 데 도움이 될 수 있다.
\kfChoice 피부를 맑게 하고 노화를 방지하는 데 도움이 될 수 있다.
\kfChoice 소화를 돕고 위장을 튼튼하게 하는 데 도움이 될 수 있다.
\kfChoice 열량 섭취를 줄이고 혈액 속 당분을 관리하는 데 도움이 될 수 있다.
\end{kfChoices}
\end{kfSentenceUnit}

\begin{kfSentenceUnit}{48}{}
\kfSentQ{48.}{이 문장에 따르면, 어떤 경우에 우리 몸에 나쁜 영향이 있을 수 있나요?}
\begin{kfChoices}
\kfChoice 운동을 규칙적으로 할 경우.
\kfChoice 천연 재료로 만든 음식만 먹을 경우.
\kfChoice 인공 감미료가 들어간 식품을 전혀 먹지 않을 경우.
\kfChoice 인공 감미료가 들어간 제로 칼로리 식품을 지나치게 먹을 경우.
\end{kfChoices}
\end{kfSentenceUnit}

\begin{kfSentenceUnit}{49}{}
\kfSentQ{49.}{인공 감미료가 건강에 미치는 영향은 무엇에 따라 다른가요?}
\begin{kfChoices}
\kfChoice 개인의 건강 상태와 먹는 양.
\kfChoice 개인이 처한 주변의 환경과 날씨.
\kfChoice 각자가 사는 지역과 직업의 종류.
\kfChoice 사람마다 다른 나이와 성별의 차이.
\end{kfChoices}
\end{kfSentenceUnit}

\begin{kfSentenceUnit}{50}{}
\kfSentQ{50.}{무엇을 생각해야 한다고 하나요?}
\begin{kfChoices}
\kfChoice 인공 감미료는 무조건 몸에 좋다고 생각해야 한다는 점.
\kfChoice 인공 감미료가 건강에 미치는 영향이 개인차가 있다는 점.
\kfChoice 인공 감미료는 무조건 몸에 나쁘다고 생각해야 한다는 점.
\kfChoice 인공 감미료는 가격이 가장 중요하다고 생각해야 한다는 점.
\end{kfChoices}
\end{kfSentenceUnit}



\clearpage

\kfChapterCover{4}{프레게3 (챕터4)}{Chapter 4}{이번 챕터의 학습 내용을 시작합니다.}

\kfStartGrammar{문법 영역 — 분석\_훈련}


\kfSecHeader{kfAreaGrammar}{개념}{분석\_훈련}

 문장 속 품사 찾기 훈련!

다음 문장을 읽고, 설명에 해당하는 품사를 모두 찾아 써 보세요.

\kfSecHeader{kfAreaGrammar}{개념}{문장의 핵심과 연결고리 (명사, 대명사, 수사, 조사)}

우리말에는 수많은 단어가 있다. 이 단어들은 문장에서 저마다 하는 역할이 정해져 있다. 이 중에서 문장의 뼈대나 주인공 역할을 하는 말들이 있다. 바로 명사, 대명사, 수사이다. 그리고 이 주인공들 뒤에 꼭 붙어서 도와주는 말이 있다. 이것이 바로 조사이다.

먼저 명사는 세상에 있는 모든 것의 이름을 나타낸다. 우리 눈에 보이는 사람, 물건, 장소의 이름은 모두 명사이다. 예를 들면 학교, 친구, 강아지, 컴퓨터, 대한민국 같은 말이다. 또 사랑, 우정, 행복, 꿈처럼 우리 눈에 보이지 않지만 마음으로 느끼거나 생각할 수 있는 것들의 이름도 명사이다.

대명사는 명사를 대신해서 부르는 말이다. 명사의 이름을 매번 부르면 불편하기 때문에 대명사를 사용한다. 예를 들어 사람 이름 ‘김철수’ 대신 ‘그’라고 한다. 물건을 가리킬 때도 '책가방' 대신 '이것'이라고 부른다. 나, 너, 우리처럼 사람을 대신하는 말도 있고, 여기, 저기처럼 장소를 대신하는 말, 무엇처럼 사물을 대신하는 말도 모두 대명사다. 수사는 숫자나 순서를 가리키는 말이다. 하나, 둘, 셋, 넷처럼 수를 셀 때 쓰거나 첫째, 둘째, 셋째처럼 순서를 나타낼 때 쓴다.

조사는 혼자서는 쓰일 수 없고, 명사, 대명사, 수사 뒤에 붙어서 쓰인다. 조사는 주인공 말들이 문장에서 어떤 역할을 하는지 알려주는 이름표와 같다. 예를 들어 '영희'라는 명사에 '가'가 붙으면 '영희가'처럼 문장의 주인이 된다. '를'이 붙으면 '영희를'처럼 행동의 대상이 된다. 조사는 특별한 의미를 더해주기도 한다. '나'라는 대명사에 '는'이 붙으면 '나는'처럼 다른 사람과 구별하는 느낌을 준다. '도'가 붙으면 '나도'처럼 함께한다는 뜻이 된다. '만'이 붙으면 '나만'처럼 오직 그것 하나라는 뜻이 된다.

이처럼 명사, 대명사, 수사는 문장의 중심이 되고, 조사는 그 뒤에 붙어 문장을 더 정확하고 풍부하게 만들어 준다.
\par\medskip\begin{itemize}[leftmargin=18pt, itemsep=2pt, topsep=2pt, label=\textbullet]
\item 명사 — '책상·사랑'(구체/추상 모두 이름을 나타냄)
\item 대명사 — '김철수'→'그', '책'→'이것'
\item 조사(격조사·보조사) — '영희가/영희를(격) · 나는/나도/나만(보조)'
\end{itemize}

\kfSecHeader{kfAreaGrammar}{문제}{}

\begin{kfQBlock}{1}{}
\kfQStem{다음 중 품사의 종류가 \uline{다른} 하나는 무엇인가요?}
\begin{kfChoices}
\kfChoice 하늘
\kfChoice 컴퓨터
\kfChoice 우리
\kfChoice 사랑
\kfChoice 대한민국
\end{kfChoices}
\end{kfQBlock}
\begin{kfQBlock}{2}{}
\kfQStem{다음 중 품사의 종류가 \uline{다른} 하나는 무엇인가요?}
\begin{kfChoices}
\kfChoice 나
\kfChoice 이것
\kfChoice 저기
\kfChoice 첫째
\kfChoice 무엇
\end{kfChoices}
\end{kfQBlock}
\begin{kfQBlock}{3}{}
\kfQStem{다음 중 품사의 종류가 \uline{다른} 하나는 무엇인가요?}
\begin{kfChoices}
\kfChoice 하나
\kfChoice 셋
\kfChoice 다섯
\kfChoice 영희
\kfChoice 둘째
\end{kfChoices}
\end{kfQBlock}
\begin{kfQBlock}{4}{}
\kfQStem{조사에 대한 설명으로 바르지 \uline{않은} 것은 무엇인가요?}
\begin{kfChoices}
\kfChoice 홀로 쓰일 수 없다.
\kfChoice 문장에서 주인공 역할을 한다.
\kfChoice 주로 명사, 대명사, 수사 뒤에 붙는다.
\kfChoice 문법적인 관계를 알려주거나 특별한 뜻을 더해준다.
\kfChoice 단어와 단어를 서로 이어주는 역할을 하기도 한다.
\end{kfChoices}
\end{kfQBlock}
\begin{kfQBlock}{5}{}
\kfQStem{문장 “아기가 밥을 먹는다.”에 사용된 조사를 모두 고른 것은 무엇인가요?}
\begin{kfChoices}
\kfChoice 이, 을
\kfChoice 가, 을
\kfChoice 이, 은
\kfChoice 가, 은
\kfChoice 이, 가
\end{kfChoices}
\end{kfQBlock}
\begin{kfQBlock}{6}{}
\kfQStem{<보기>의 밑줄 친 부분과 품사가 같은 것은 무엇인가요?}
\begin{kfEvidence}바나나 \uline{하나}를 주세요.\end{kfEvidence}
\begin{kfChoices}
\kfChoice 너
\kfChoice 하늘
\kfChoice 둘째
\kfChoice 빨리
\kfChoice 예쁘다
\end{kfChoices}
\end{kfQBlock}
\begin{kfQBlock}{7}{}
\kfQStem{<보기>의 밑줄 친 부분과 품사가 같은 것은 무엇인가요?}
\begin{kfEvidence}\uline{이것}은 내 책이다.\end{kfEvidence}
\begin{kfChoices}
\kfChoice 저기
\kfChoice 영희
\kfChoice 가방
\kfChoice 새
\kfChoice 먹다
\end{kfChoices}
\end{kfQBlock}
\begin{kfQBlock}{1}{}
\kfQStem{문장에서 뼈대나 주인공 역할을 하는 말 3가지는 무엇인가요?}
\kfAnswerLines{1}
\end{kfQBlock}
\begin{kfQBlock}{2}{}
\kfQStem{‘학교, 친구, 사랑’처럼 세상 모든 것의 이름을 나타내는 품사는 무엇인가요?}
\kfAnswerLines{1}
\end{kfQBlock}
\begin{kfQBlock}{3}{}
\kfQStem{‘나, 너, 우리, 이것’처럼 명사를 대신해서 부르는 품사는 무엇인가요?}
\kfAnswerLines{1}
\end{kfQBlock}
\begin{kfQBlock}{4}{}
\kfQStem{‘하나, 둘, 첫째’처럼 수량이나 순서를 나타내는 품사는 무엇인가요?}
\kfAnswerLines{1}
\end{kfQBlock}
\begin{kfQBlock}{5}{}
\kfQStem{주인공 말 뒤에 찰싹 붙어서 도와주는 품사는 무엇인가요?}
\kfAnswerLines{1}
\end{kfQBlock}
\begin{kfQBlock}{6}{}
\kfQStem{‘우정, 행복, 꿈’처럼 눈에 보이지 않는 생각이나 느낌의 이름도 속하는 품사는 무엇인가요?}
\kfAnswerLines{1}
\end{kfQBlock}
\begin{kfQBlock}{7}{}
\kfQStem{‘이것, 저것, 여기, 저기’처럼 사물이나 장소를 대신 가리키는 품사는 무엇인가요?}
\kfAnswerLines{1}
\end{kfQBlock}
\begin{kfQBlock}{8}{}
\kfQStem{‘첫째, 둘째, 셋째’처럼 순서를 나타내는 품사는 무엇인가요?}
\kfAnswerLines{1}
\end{kfQBlock}
\begin{kfQBlock}{9}{}
\kfQStem{문장에서 절대로 홀로 쓰일 수 없고, 항상 다른 말 뒤에 붙어 쓰이는 품사는 무엇인가요?}
\kfAnswerLines{1}
\end{kfQBlock}
\begin{kfQBlock}{10}{}
\kfQStem{문장 “나도 너만 믿는다.”에서 ‘역시’, ‘오직 그것’이라는 특별한 의미를 더해주는 조사 두 가지는 무엇인가요?}
\kfAnswerLines{1}
\end{kfQBlock}
\begin{kfQBlock}{1}{}
\kfQStem{다음 문장에서 명사, 대명사, 수사, 조사를 각각 모두 찾아 쓰세요.}
\begin{kfEvidence}우리는 공원에서 사과 하나를 보았다.\end{kfEvidence}
\kfAnswerNote{답안}{4}
\end{kfQBlock}

\kfSecHeader{kfAreaGrammar}{활동}{분석\_훈련}
\noindent\textit{\small 다음 활용 사례를 분석하여 빈칸을 채워 보세요.}\par\medskip
\begin{enumerate}[leftmargin=22pt, itemsep=6pt, topsep=4pt, label=\textbf{\arabic*.}]
\item 다음 문장에서 사람이나 사물의 이름을 나타내는 명사를 모두 찾아 써 보세요.
\begin{kfEvidence}토끼가 풀을 먹는다.\end{kfEvidence}
\item 다음 문장에서 명사를 대신하여 쓰인 대명사를 모두 찾아 써 보세요.
\begin{kfEvidence}그는 나의 친구이다.\end{kfEvidence}
\item 다음 문장에서 수량이나 순서를 나타내는 수사를 모두 찾아 써 보세요.
\begin{kfEvidence}하나부터 열까지 차근차근 세어 보아라.\end{kfEvidence}
\item 다음 문장에서 다른 말에 붙어 관계를 나타내는 조사를 모두 찾아 써 보세요.
\begin{kfEvidence}하늘이 푸르고 바다가 넓다.\end{kfEvidence}
\item 다음 문장에서 명사와 대명사를 모두 찾아 써 보세요.
\begin{kfEvidence}이것은 연필이고 저것은 지우개이다.\end{kfEvidence}
\item 다음 문장에서 수사와 조사를 모두 찾아 써 보세요.
\begin{kfEvidence}첫째는 건강이다.\end{kfEvidence}
\item 다음 문장에서 사람이나 사물의 이름을 나타내는 명사를 모두 찾아 써 보세요.
\begin{kfEvidence}영희가 학교에서 공부를 한다.\end{kfEvidence}
\item 다음 문장에서 다른 말에 붙어 관계를 나타내는 조사를 모두 찾아 써 보세요.
\begin{kfEvidence}나도 너처럼 할 수 있다.\end{kfEvidence}
\item 다음 문장에서 명사를 대신하여 쓰인 대명사를 모두 찾아 써 보세요.
\begin{kfEvidence}너는 여기에서 무엇을 했니?\end{kfEvidence}
\item 다음 문장에서 명사, 대명사, 수사, 조사를 모두 찾아 써 보세요.
\begin{kfEvidence}그가 우리에게 사과 하나를 주었다.\end{kfEvidence}
\end{enumerate}


\kfStartLit{문학 영역 — 작품}


\kfSecHeader{kfAreaLiterature}{지문}{작품}
\kfPassageNote{작품}{%
처음에 북부여 왕 해부루가 동부여로 자리를 피하고 나서 해부루가 죽으니 금와가 왕이 되었습니다. 이때에 왕은 태백산 남쪽 우발수에서 한 여자를 만나서 사정을 물어보니 그가 말했습니다. "나는 물의 신 하백의 딸로서 이름은 유화인데, 여러 동생들과 함께 나와서 놀던 중에 어떤 남자가 나타나서 하늘의 제왕 아들 해모수라고 하면서 나를 꾀어서 웅신산 아래 압록강 가에서 함께 지내고는 가서 돌아오지 않았습니다. 부모님은 내가 중매도 없이 다른 남자를 따랐다고 꾸짖었습니다. 그래서 지금 이곳에서 쫓겨나 살고 있습니다." 금와왕은 유화의 이야기를 듣고 그녀를 궁으로 데려갔습니다. 궁에서 지내던 유화는 어느 날 햇빛을 받게 되었습니다. 금와왕이 이를 이상하게 여겨서 방 속에 깊이 가두었더니 햇빛이 그녀를 비추었습니다. 그녀가 이를 피했으나 햇빛은 또 쫓아와서 비추곤 했습니다.

그렇게 임신하여 알 한 개를 낳으니 크기가 다섯 되 정도 되었습니다. 왕이 이것을 버려서 개와 돼지에게 주니 모두 먹지 않았습니다. 다시 이것을 길바닥에 버렸더니 소와 말이 피해갔습니다. 이것을 들에 버렸더니 새와 짐승이 덮어주었습니다. 왕이 이것을 깨뜨리려 해도 깨뜨릴 수가 없어서 그만 그 어미에게 돌려주었습니다. 어미가 이것을 천으로 감싸서 따뜻한 곳에 두었더니 아이 하나가 껍질을 깨고서 나왔는데, 몸과 외모가 영리하고 신기하게 생겼습니다. 나이 겨우 일곱 살에 뛰어나게 자라서 제 손으로 활과 화살을 만들어 백 번 쏘면 백 번 맞혔습니다. 이 나라 풍속에 활 잘 쏘는 사람을 주몽이라 하므로 이로써 ‘주몽’이라 이름을 지었습니다.

금와왕에게는 아들 일곱이 있어서 언제나 주몽과 함께 놀았는데 그의 재주를 따를 수 없었습니다. 맏아들 대소가 왕에게 말했습니다. "주몽은 사람의 자식이 아닙니다. 만일 빨리 처리하지 않는다면 나중에 문제가 될 것입니다." 왕은 이 말을 듣지 않고 그를 시켜서 말을 기르게 하였습니다. 주몽은 그중 빠른 말에게는 먹이를 적게 주어서 마르도록 만들고, 느린 말은 잘 먹여서 살이 찌도록 하였습니다. 왕은 살찐 말을 자신이 타고 마른 말을 주몽에게 주었습니다. 왕자들과 여러 신하들이 장차 그를 해치려는 것을 주몽의 어머니가 알고 그에게 말했습니다. "나라 사람들이 장차 너를 해치려고 하는데, 너 같은 재주를 가지고 어디로 가든지 못 살겠느냐? 빨리 떠나는 것이 좋을 것이다."

이에 주몽은 오이 등 세 사람과 친구가 되어서 엄수까지 와서 물에게 말했습니다. "나는 하늘 제왕의 아들이고, 물의 신 하백의 손자다. 오늘 도망가는 길에 뒤따르는 사람이 쫓아 오니 이 일을 어떻게 할 것인가?" 이때에 물고기와 자라들이 나와서 다리가 되어 물을 건너게 해주었습니다. 그러고 나서 다리가 풀려버려서 쫓아오던 말 탄 사람들은 물을 건널 수가 없었습니다. 그는 졸본주까지 와서 드디어 여기에 도읍을 정하였습니다. 아직 궁궐을 지을 시간도 없어서 그저 비류수 가에 초가집을 짓고 살면서 나라 이름을 고구려라 하고, 고 씨로 성을 삼았습니다. 당시의 그의 나이가 열두 살이었습니다.
}


\kfSecHeader{kfAreaLiterature}{문제}{}

\begin{kfQBlock}{1}{}
\kfQStem{이 글의 내용과 일치하지 않는 것은 무엇인가요?}
\begin{kfChoices}
\kfChoice 금와왕은 처음부터 주몽이 큰 인물이 될 것을 알고 정성껏 보살폈다.
\kfChoice 주몽은 사람의 모습이 아닌 알의 형태로 태어났다.
\kfChoice 금와왕의 아들 대소는 주몽의 뛰어난 능력을 질투했다.
\kfChoice 주몽은 12살이라는 어린 나이에 고구려를 세우고 고 씨 성을 썼다.
\kfChoice 주몽의 어머니 유화는 물의 신 하백의 딸이다.
\end{kfChoices}
\end{kfQBlock}
\begin{kfQBlock}{2}{}
\kfQStem{주몽이 ‘주몽’이라는 이름을 갖게 된 이유는 무엇인가요?}
\begin{kfChoices}
\kfChoice 어머니 유화가 지어준 아명이라서
\kfChoice 위대한 하늘의 제왕이 직접 내려준 이름이라서
\kfChoice 그 나라 말로 ‘활 잘 쏘는 사람’을 뜻해서
\kfChoice 온몸이 금빛으로 빛나는 개구리 모양으로 태어났기 때문에
\kfChoice 알을 깨고 밖으로 나올 때 스스로 ‘주몽’이라고 크게 외쳤기 때문에
\end{kfChoices}
\end{kfQBlock}
\begin{kfQBlock}{3}{}
\kfQStem{금와왕의 아들 대소가 주몽을 없애야 한다고 주장한 가장 큰 이유는 무엇인가요?}
\begin{kfChoices}
\kfChoice 주몽이 말을 제대로 기르지 못해서
\kfChoice 주몽이 아버지 해모수를 그리워해서
\kfChoice 주몽의 재주가 자신들보다 훨씬 뛰어나서
\kfChoice 주몽이 평소에 너무 게으르고 욕심이 많았기 때문에
\kfChoice 주몽이 금와왕의 왕위를 몰래 넘보고 있었기 때문에
\end{kfChoices}
\end{kfQBlock}
\begin{kfQBlock}{4}{}
\kfQStem{주몽이 금와왕으로부터 좋은 말을 얻기 위해 사용한 지혜로운 방법은 무엇인가요?}
\begin{kfChoices}
\kfChoice 날랜 말은 굶기고, 둔한 말은 잘 먹였다.
\kfChoice 둔한 말을 훔쳐서 날랜 말과 몰래 바꾸었다.
\kfChoice 금와왕에게 날랜 말을 달라고 눈물로 호소했다.
\kfChoice 밤에 몰래 마굿간에 들어가 말들을 훈련시켰다.
\kfChoice 사냥 대회에서 이겨 상품으로 날랜 말을 받았다.
\end{kfChoices}
\end{kfQBlock}
\begin{kfQBlock}{5}{}
\kfQStem{부여를 탈출하던 주몽이 다리 없는 강을 건널 수 있었던 이유는 무엇인가요?}
\begin{kfChoices}
\kfChoice 친구들과 함께 뗏목을 만들어 건넜다.
\kfChoice 말을 타고 강물을 단숨에 뛰어넘었다.
\kfChoice 물고기와 자라들이 떠올라 다리를 만들어주었다.
\kfChoice 갑자기 하늘에서 거대한 용이 내려와 그를 태워주었다.
\kfChoice 어머니 유화 부인이 신비한 마법으로 강물을 멈춰주었다.
\end{kfChoices}
\end{kfQBlock}
\begin{kfQBlock}{6}{}
\kfQStem{이 글에 나타난 주몽의 부모님에 대한 설명으로 옳은 것은 무엇인가요?}
\begin{kfChoices}
\kfChoice 아버지는 알 수 없고, 어머니는 부여의 공주이다.
\kfChoice 아버지는 부여의 왕 해부루, 어머니는 하백의 딸 유화이다.
\kfChoice 아버지는 물의 신 하백이고, 어머니는 하늘 제왕의 고귀한 딸이다.
\kfChoice 아버지는 하늘 제왕의 아들 해모수, 어머니는 하백의 딸 유화이다.
\kfChoice 아버지는 동부여의 왕 금와, 어머니는 인간 세상의 평범한 여자이다.
\end{kfChoices}
\end{kfQBlock}
\begin{kfQBlock}{7}{}
\kfQStem{주몽이 알의 상태로 버려졌을 때, 알을 보호해 준 존재가 \uline{아닌} 것은 무엇인가요?}
\begin{kfChoices}
\kfChoice 자라
\kfChoice 소
\kfChoice 새
\kfChoice 돼지
\kfChoice 개
\end{kfChoices}
\end{kfQBlock}
\begin{kfQBlock}{8}{}
\kfQStem{이 글에 대한 설명으로 가장 알맞은 것은 무엇인가요?}
\begin{kfChoices}
\kfChoice 한 인물의 평범한 일상을 재미있게 그린 이야기이다.
\kfChoice 어떤 물건이 어떻게 만들어졌는지 설명하는 이야기이다.
\kfChoice 나라를 세운 영웅의 신비로운 탄생과 활약을 다룬 신화이다.
\kfChoice 동물들의 세계를 통해 인간 세상을 비판하는 재미있는 우화이다.
\kfChoice 실제 역사적 사실을 바탕으로 작가의 상상력을 풍부하게 더해 만든 역사 소설이다.
\end{kfChoices}
\end{kfQBlock}
\begin{kfQBlock}{1}{}
\kfQStem{주몽이 어릴 때부터 평범한 아이가 아니었음을 보여주는 모습 두 가지를 서술하세요.}
\kfAnswerNote{답안}{4}
\end{kfQBlock}
\begin{kfQBlock}{2}{}
\kfQStem{주몽이 여러 위기에서 벗어날 수 있도록 도와준 존재들을 두 종류로 나누어 서술하세요.}
\kfAnswerNote{답안}{4}
\end{kfQBlock}

\kfSecHeader{kfAreaLiterature}{활동}{문장\_독해}
\noindent\textit{\small 다음 문장을 읽고 물음에 답하시오.}\par\medskip
\begin{kfSentenceUnit}{1}{}
\kfSentQ{1.}{다음 문장에서 서술어를 기준으로 각각의 주어는 무엇인가요?(1) 서술어 '버려서'의 주어는 무엇인가요?}
\begin{kfChoices}
\kfChoice 내가
\kfChoice 네가
\kfChoice 왕이
\kfChoice 개와 돼지가
\end{kfChoices}
\kfSentQ{1.}{(2) 서술어 '주니'의 주어는 무엇인가요?}
\begin{kfChoices}
\kfChoice 내가
\kfChoice 네가
\kfChoice 왕이
\kfChoice 개와 돼지가
\end{kfChoices}
\kfSentQ{1.}{(3) 서술어 '먹지 않았습니다'의 주어는 무엇인가요?}
\begin{kfChoices}
\kfChoice 내가
\kfChoice 네가
\kfChoice 왕이
\kfChoice 개와 돼지가
\end{kfChoices}
\end{kfSentenceUnit}

\begin{kfSentenceUnit}{2}{}
\kfSentQ{2.}{다음 문장에서 '이것'이 가리키는 것은 무엇인가요?}
\begin{kfChoices}
\kfChoice 알
\kfChoice 활
\kfChoice 말
\kfChoice 햇빛
\end{kfChoices}
\end{kfSentenceUnit}

\begin{kfSentenceUnit}{3}{}
\kfSentQ{3.}{다음 문장에서 '어미'와 ‘이것’이 각각 가리키는 것은 무엇인가요?}
\begin{kfChoices}
\kfChoice 어미 - 유화, 이것 - 알
\kfChoice 어미 - 유화, 이것 - 햇빛
\kfChoice 어미 - 하백의 딸, 이것 - 활
\kfChoice 어미 - 평범한 여자, 이것 - 말
\end{kfChoices}
\end{kfSentenceUnit}


\kfStartNonlit{비문학 영역 — [기술] 조상들의 지혜, 천연 냉장고 석빙고}


\kfSecHeader{kfAreaNonlit}{지문}{[기술] 조상들의 지혜, 천연 냉장고 석빙고}
\kfPassageNote{[기술] 조상들의 지혜, 천연 냉장고 석빙고}{%
우리 조상들은 더운 여름에도 얼음을 사용할 수 있었다. 겨울에 만든 얼음을 가을까지 보관하는 창고가 있었기 때문이다. 이 창고를 석빙고라고 한다. 석빙고는 돌로 만든 얼음 창고라는 뜻이다.

석빙고에서 얼음이 오래 보관될 수 있는 이유는 무엇일까? 먼저 얼음이 녹는 과정을 살펴보자. 얼음이 물로 변할 때는 주변의 열을 빨아들인다. 그래서 얼음 주변의 공기가 차가워진다. 이렇게 차가워진 공기 때문에 다른 얼음들은 쉽게 녹지 않는다. 하지만 녹은 물은 빨리 밖으로 내보내야 한다. 그래서 조상들은 석빙고 바닥을 비스듬하게 만들어 물이 잘 흘러나가도록 했다.

석빙고를 아무리 꽁꽁 막아도 밖의 더운 공기가 조금씩 들어온다. 이 문제를 해결하기 위해 조상들은 천장에 구멍을 뚫었다. 따뜻한 공기는 가벼워서 위로 올라가고, 차가운 공기는 무거워서 아래로 내려온다. 천장의 구멍으로 따뜻한 공기가 밖으로 나가면 석빙고 안이 계속 시원하게 유지된다. 이 구멍에는 비나 햇빛이 들어오지 않도록 돌로 덮개를 만들었다.

얼음을 더 오래 보관하기 위해 얼음 사이사이에 짚을 넣었다. 짚은 열이 잘 전달되지 않는 성질이 있다. 얼음끼리 직접 닿아 있으면 열이 쉽게 옮겨가지만, 짚이 가운데 있으면 열이 천천히 이동한다. 또한 짚에는 작은 공기 구멍이 많이 있어서 단열 효과가 크다. 석빙고 주변에도 특별한 장치들이 있었다. 건물 위에 흙을 덮어서 밖의 열이 안으로 들어오는 것을 막았다. 그리고 풀을 심어서 태양열이 직접 닿지 않도록 했다. 얼음을 저장하는 방은 땅속 깊이 만들어서 온도가 더 낮게 유지되도록 했다.

석빙고는 조상들의 과학적 지혜가 담긴 천연 냉장고였다. 전기도 없던 시대에 이런 기술을 만들어낸 것은 정말 놀라운 일이다.
}


\kfSecHeader{kfAreaNonlit}{문제}{}

\begin{kfQBlock}{1}{}
\kfQStem{다음 중 석빙고에서 얼음이 오래 보관될 수 있는 원리로 가장 적절하지 \uline{않은} 것은 무엇인가요?}
\begin{kfChoices}
\kfChoice 천장 구멍을 통해 공기가 순환된다.
\kfChoice 바닥을 비스듬하게 만들어 물을 배출한다.
\kfChoice 얼음이 녹을 때 주변의 열을 빨아들인다.
\kfChoice 전기를 사용하여 인공적으로 온도를 낮춘다.
\kfChoice 얼음 사이에 짚을 넣어서 단열 효과를 높인다.
\end{kfChoices}
\end{kfQBlock}
\begin{kfQBlock}{2}{}
\kfQStem{윗글의 내용과 일치하지 않는 것은 무엇인가요?}
\begin{kfChoices}
\kfChoice 석빙고는 돌을 주된 재료로 사용하여 만들어졌다.
\kfChoice 석빙고 지붕 위에 흙과 풀을 덮어 외부 열을 차단했다.
\kfChoice 얼음이 녹으면서 나오는 물은 얼음 보관에 방해가 될 수 있다.
\kfChoice 짚은 얼음과 얼음이 직접 닿아 열이 더 빨리 전달되도록 돕는다.
\kfChoice 차가운 공기는 따뜻한 공기보다 무거워 아래로 내려오는 성질이 있다.
\end{kfChoices}
\end{kfQBlock}
\begin{kfQBlock}{3}{}
\kfQStem{'짚'의 역할로 적절하지 \uline{않은} 것은 무엇인가요?}
\begin{kfChoices}
\kfChoice 얼음 사이의 열 전달을 느리게 한다.
\kfChoice 작은 공기 구멍을 통해 단열 효과를 높인다.
\kfChoice 녹은 물을 흡수하여 바닥이 축축해지는 것을 막는다.
\kfChoice 얼음끼리 직접 닿는 것을 막아 열의 이동을 방해한다.
\kfChoice 얼음이 녹는 것을 막아주는 보온병과 비슷한 원리로 작용한다.
\end{kfChoices}
\end{kfQBlock}
\begin{kfQBlock}{4}{}
\kfQStem{석빙고의 원리가 적용된 사례로 보기 어려운 것은 무엇인가요?}
\begin{kfChoices}
\kfChoice 겨울철에 창문에 에어캡을 붙여 집을 따뜻하게 한다.
\kfChoice 에어컨을 천장에 설치하여 찬 공기가 아래로 내려오게 한다.
\kfChoice 검은색 옷을 입어 햇빛을 잘 흡수하여 몸을 따뜻하게 한다.
\kfChoice 아이스박스 안에 음식과 얼음을 함께 넣어 시원하게 보관한다.
\kfChoice 컵라면 용기를 스티로폼으로 만들어 뜨거운 물이 잘 식지 않게 한다.
\end{kfChoices}
\end{kfQBlock}
\begin{kfQBlock}{5}{}
\kfQStem{윗글을 통해 알 수 있는 우리 조상들의 모습으로 가장 적절한 것은 무엇인가요?}
\begin{kfChoices}
\kfChoice 편리함보다는 예술적인 아름다움을 더 우선시했다.
\kfChoice 자연의 힘에 순응하며 모든 것을 운명으로 받아들였다.
\kfChoice 다른 나라의 기술을 그대로 모방하여 빠르게 발전했다.
\kfChoice 자연 현상을 과학적으로 이해하고 생활에 적극적으로 활용했다.
\kfChoice 실용적인 기술보다는 이론적인 학문 탐구를 더 중요하게 여겼다.
\end{kfChoices}
\end{kfQBlock}
\begin{kfQBlock}{6}{}
\kfQStem{다음 <보기>는 석빙고의 원리를 이해한 학생들이 나눈 대화입니다. 잘못 이해한 학생은 누구일까요?}
\begin{kfEvidence}철수: "석빙고는 얼음이 녹으면서 열을 뺏어가는 성질을 이용했어."

영희: "맞아. 그리고 천장에 구멍을 뚫어서 뜨거운 공기를 밖으로 보냈지."

민수: "바닥을 기울게 만들어서 녹은 물이 밖으로 빨리 빠지도록 한 것도 중요해."

지혜: "얼음 사이에 짚을 넣은 건 얼음끼리 붙여서 더 꽝꽝 얼게 하기 위해서야."\end{kfEvidence}
\begin{kfChoices}
\kfChoice 철수
\kfChoice 영희
\kfChoice 민수
\kfChoice 지혜
\kfChoice 모두 옳다
\end{kfChoices}
\end{kfQBlock}
\begin{kfQBlock}{7}{}
\kfQStem{다음 <보기>의 발명품 중 석빙고의 '천장 구멍'과 가장 비슷한 원리를 이용한 것은 무엇인가요?}
\begin{kfEvidence}석빙고의 천장 구멍은 가벼운 더운 공기는 위로 올라가고, 무거운 찬 공기는 아래로 내려오는 성질을 이용하여 내부를 시원하게 유지한다.\end{kfEvidence}
\begin{kfChoices}
\kfChoice 보온병
\kfChoice 전기장판
\kfChoice 굴뚝이 있는 난로
\kfChoice 회전하는 선풍기
\kfChoice 빛을 모으는 돋보기
\end{kfChoices}
\end{kfQBlock}
\begin{kfQBlock}{1}{}
\kfQStem{겨울에 만든 얼음을 여름까지 보관하던 돌로 만든 창고의 이름은 무엇인가요?}
\kfAnswerLines{1}
\end{kfQBlock}
\begin{kfQBlock}{2}{}
\kfQStem{얼음이 물로 변할 때 무엇을 빨아들이나요?}
\kfAnswerLines{1}
\end{kfQBlock}
\begin{kfQBlock}{3}{}
\kfQStem{녹은 물이 잘 빠져나가도록 하기 위해 석빙고 바닥을 어떻게 만들었나요?}
\kfAnswerLines{1}
\end{kfQBlock}
\begin{kfQBlock}{4}{}
\kfQStem{석빙고 안의 더워진 공기를 밖으로 내보내기 위해 어느 곳에 구멍을 뚫었나요?}
\kfAnswerLines{1}
\end{kfQBlock}
\begin{kfQBlock}{5}{}
\kfQStem{따뜻한 공기와 차가운 공기 중 어느 것이 더 가벼운가?}
\kfAnswerLines{1}
\end{kfQBlock}
\begin{kfQBlock}{1}{}
\kfQStem{석빙고가 외부의 더운 열을 차단하기 위해 사용한 방법 두 가지를 지문에서 찾아 서술하세요.}
\kfAnswerNote{답안}{4}
\end{kfQBlock}

\kfSecHeader{kfAreaNonlit}{활동}{문장\_독해}
\noindent\textit{\small 다음 문장을 읽고 물음에 답하시오.}\par\medskip
\begin{kfSentenceUnit}{1}{}
\kfSentQ{1.}{우리 조상들은 더운 여름에도 무엇을 사용할 수 있었나요?}
\begin{kfChoices}
\kfChoice 불.
\kfChoice 얼음.
\kfChoice 나무.
\kfChoice 바람.
\end{kfChoices}
\end{kfSentenceUnit}

\begin{kfSentenceUnit}{2}{}
\kfSentQ{2.}{이 문장에서 설명하는, 우리 조상들이 가졌던 놀라운 능력은 무엇인가요?}
\begin{kfChoices}
\kfChoice 엄청난 양의 식량을 생산할 수 있었던 능력.
\kfChoice 더운 여름에도 얼음을 사용할 수 있었던 능력.
\kfChoice 먼 거리를 아주 빠르게 이동할 수 있었던 능력.
\kfChoice 다른 나라와 활발하게 무역을 할 수 있었던 능력.
\end{kfChoices}
\end{kfSentenceUnit}

\begin{kfSentenceUnit}{3}{}
\kfSentQ{3.}{우리 조상들은 언제 얼음을 만들었나요?}
\begin{kfChoices}
\kfChoice 봄.
\kfChoice 겨울.
\kfChoice 여름.
\kfChoice 가을.
\end{kfChoices}
\end{kfSentenceUnit}

\begin{kfSentenceUnit}{4}{}
\kfSentQ{4.}{우리 조상들은 얼음을 언제까지 보관했나요?}
\begin{kfChoices}
\kfChoice 봄.
\kfChoice 가을.
\kfChoice 여름.
\kfChoice 겨울.
\end{kfChoices}
\end{kfSentenceUnit}

\begin{kfSentenceUnit}{5}{}
\kfSentQ{5.}{우리 조상들이 여름에도 얼음을 사용할 수 있었던 이유는 무엇인가요?}
\begin{kfChoices}
\kfChoice 얼음을 만드는 아주 특별한 기계가 있었기 때문이다.
\kfChoice 여름에도 눈이 내리는 아주 추운 지역에 살았기 때문이다.
\kfChoice 다른 나라에서 얼음을 싼 가격에 수입할 수 있었기 때문이다.
\kfChoice 겨울에 만든 얼음을 가을까지 보관하는 창고가 있었기 때문이다.
\end{kfChoices}
\end{kfSentenceUnit}

\begin{kfSentenceUnit}{6}{}
\kfSentQ{6.}{'이 창고'가 가리키는 것은 무엇인가요?}
\begin{kfChoices}
\kfChoice 식량을 오랫동안 싱싱하게 보관하는 창고.
\kfChoice 왕의 보물을 안전하게 숨겨두는 비밀 창고.
\kfChoice 전쟁에 필요한 무기를 많이 모아두는 창고.
\kfChoice 겨울에 만든 얼음을 가을까지 보관하는 창고.
\end{kfChoices}
\end{kfSentenceUnit}

\begin{kfSentenceUnit}{7}{}
\kfSentQ{7.}{석빙고는 무엇으로 만든 창고인가요?}
\begin{kfChoices}
\kfChoice 돌.
\kfChoice 흙.
\kfChoice 철.
\kfChoice 나무.
\end{kfChoices}
\end{kfSentenceUnit}

\begin{kfSentenceUnit}{8}{}
\kfSentQ{8.}{석빙고는 무엇을 보관하는 창고인가요?}
\begin{kfChoices}
\kfChoice 얼음.
\kfChoice 곡식.
\kfChoice 무기.
\kfChoice 보물.
\end{kfChoices}
\end{kfSentenceUnit}

\begin{kfSentenceUnit}{9}{}
\kfSentQ{9.}{'석빙고'라는 이름의 뜻은 무엇인가요?}
\begin{kfChoices}
\kfChoice 돌로 만든 얼음 창고.
\kfChoice 철로 만든 튼튼한 창고.
\kfChoice 나무로 만든 서늘한 창고.
\kfChoice 흙으로 빚은 도자기 창고.
\end{kfChoices}
\end{kfSentenceUnit}

\begin{kfSentenceUnit}{10}{}
\kfSentQ{10.}{이 질문을 통해 앞으로 어떤 내용이 나올 것을 예상할 수 있나요?}
\begin{kfChoices}
\kfChoice 우리나라의 전통 건축 양식에 대한 자세한 설명.
\kfChoice 얼음을 만드는 방법에 대한 과학적인 원리 설명.
\kfChoice 조선 시대 사람들의 생활 모습에 대한 전반적인 설명.
\kfChoice 석빙고에서 얼음이 오래 보관될 수 있는 이유에 대한 설명.
\end{kfChoices}
\end{kfSentenceUnit}

\begin{kfSentenceUnit}{11}{}
\kfSentQ{11.}{이 문장을 통해 앞으로 무엇에 대해 설명할 것인지 알 수 있나요?}
\begin{kfChoices}
\kfChoice 얼음이 녹는 과정.
\kfChoice 얼음이 어는 과정.
\kfChoice 소금이 녹는 과정.
\kfChoice 물이 증발하는 과정.
\end{kfChoices}
\end{kfSentenceUnit}

\begin{kfSentenceUnit}{12}{}
\kfSentQ{12.}{얼음이 무엇으로 변하나요?}
\begin{kfChoices}
\kfChoice 물.
\kfChoice 구름.
\kfChoice 안개.
\kfChoice 수증기.
\end{kfChoices}
\end{kfSentenceUnit}

\begin{kfSentenceUnit}{13}{}
\kfSentQ{13.}{얼음이 물로 변할 때 무엇을 빨아들이나요?}
\begin{kfChoices}
\kfChoice 주변의 열.
\kfChoice 땅속의 흙.
\kfChoice 차가운 공기.
\kfChoice 주변의 수분.
\end{kfChoices}
\end{kfSentenceUnit}

\begin{kfSentenceUnit}{14}{}
\kfSentQ{14.}{'그래서'의 구체적인 내용은 무엇인가요?}
\begin{kfChoices}
\kfChoice 얼음이 물로 변할 때 아주 큰 소리를 내서.
\kfChoice 얼음이 물로 변할 때 주변의 열을 빨아들여서.
\kfChoice 얼음이 물로 변할 때 주변으로 물방울을 튀겨서.
\kfChoice 얼음이 물로 변할 때 주변의 열을 밖으로 내보내서.
\end{kfChoices}
\end{kfSentenceUnit}

\begin{kfSentenceUnit}{15}{}
\kfSentQ{15.}{'그래서'라는 표현을 볼 때, 얼음 주변 공기가 차가워지는 원인은 무엇인가요?}
\begin{kfChoices}
\kfChoice 얼음이 녹으면서 주변의 열을 빨아들였기 때문이다.
\kfChoice 겨울바람이 창고 안으로 계속 불어 들어왔기 때문이다.
\kfChoice 차가운 지하수가 흘러들어와 주변을 식혀주었기 때문이다.
\kfChoice 얼음이 꽁꽁 얼면서 주변의 공기를 차갑게 만들었기 때문이다.
\end{kfChoices}
\end{kfSentenceUnit}

\begin{kfSentenceUnit}{16}{}
\kfSentQ{16.}{'이렇게 차가워진 공기'는 무엇을 가리키나요?}
\begin{kfChoices}
\kfChoice 밖에서 들어온 뜨거운 공기.
\kfChoice 사람들이 내뿜는 따뜻한 입김.
\kfChoice 땅속 깊은 곳에서 올라오는 열기.
\kfChoice 얼음이 열을 빨아들여서 차가워진 공기.
\end{kfChoices}
\end{kfSentenceUnit}

\begin{kfSentenceUnit}{17}{}
\kfSentQ{17.}{다른 얼음이 쉽게 녹지 않는 이유는 무엇인가요?}
\begin{kfChoices}
\kfChoice 얼음이 녹아 차가워진 공기 때문이다.
\kfChoice 창고 문을 자주 열어두었기 때문이다.
\kfChoice 얼음이 서로 꽉 붙어 있기 때문이다.
\kfChoice 창고 안의 온도가 매우 높기 때문이다.
\end{kfChoices}
\end{kfSentenceUnit}

\begin{kfSentenceUnit}{18}{}
\kfSentQ{18.}{무엇을 밖으로 내보내야 하나요?}
\begin{kfChoices}
\kfChoice 녹은 물.
\kfChoice 얼음 조각.
\kfChoice 차가운 바람.
\kfChoice 뜨거운 수증기.
\end{kfChoices}
\end{kfSentenceUnit}

\begin{kfSentenceUnit}{19}{}
\kfSentQ{19.}{얼음 보관을 위해 꼭 해야 할 일은 무엇인가요?}
\begin{kfChoices}
\kfChoice 얼음을 밖으로 꺼내서 말리는 것.
\kfChoice 뜨거운 바람을 안으로 불어넣는 것.
\kfChoice 창고 문을 활짝 열어 환기하는 것.
\kfChoice 녹은 물을 빨리 밖으로 내보내는 것.
\end{kfChoices}
\end{kfSentenceUnit}

\begin{kfSentenceUnit}{20}{}
\kfSentQ{20.}{'그래서'의 구체적인 내용은 무엇인가요?}
\begin{kfChoices}
\kfChoice 녹은 물을 빨리 밖으로 내보내야 해서.
\kfChoice 녹은 물이 다시 얼음으로 변해야 해서.
\kfChoice 녹은 물이 얼음을 깨끗하게 씻어줘야 해서.
\kfChoice 녹은 물이 고여 있으면 얼음이 더 빨리 녹아서.
\end{kfChoices}
\end{kfSentenceUnit}

\begin{kfSentenceUnit}{21}{}
\kfSentQ{21.}{조상들이 석빙고 바닥을 어떻게 만들었나요?}
\begin{kfChoices}
\kfChoice 평평하게.
\kfChoice 비스듬하게.
\kfChoice 계단식으로.
\kfChoice 울퉁불퉁하게.
\end{kfChoices}
\end{kfSentenceUnit}

\begin{kfSentenceUnit}{22}{}
\kfSentQ{22.}{조상들이 석빙고 바닥을 비스듬하게 만든 이유는 무엇인가요?}
\begin{kfChoices}
\kfChoice 녹은 물이 잘 흘러나가도록 하기 위해.
\kfChoice 얼음을 차곡차곡 쌓기 편하게 하기 위해.
\kfChoice 뜨거운 공기가 바닥에 머물도록 하기 위해.
\kfChoice 사람들이 걸어 다닐 때 미끄러지지 않게 하기 위해.
\end{kfChoices}
\end{kfSentenceUnit}

\begin{kfSentenceUnit}{23}{}
\kfSentQ{23.}{석빙고를 꽁꽁 막아도 무엇이 들어오나요?}
\begin{kfChoices}
\kfChoice 차가운 빗물.
\kfChoice 작은 벌레들.
\kfChoice 시원한 바람.
\kfChoice 밖의 더운 공기.
\end{kfChoices}
\end{kfSentenceUnit}

\begin{kfSentenceUnit}{24}{}
\kfSentQ{24.}{석빙고를 관리할 때 발생하는 문제점은 무엇인가요?}
\begin{kfChoices}
\kfChoice 밖의 더운 공기가 조금씩 들어온다는 문제점.
\kfChoice 창고 안에 물이 고여서 얼음이 썩는 문제점.
\kfChoice 얼음이 너무 꽁꽁 얼어서 깨지지 않는 문제점.
\kfChoice 창고 안이 너무 어두워서 앞이 보이지 않는 문제점.
\end{kfChoices}
\end{kfSentenceUnit}

\begin{kfSentenceUnit}{25}{}
\kfSentQ{25.}{'이 문제'가 가리키는 것은 무엇인가요?}
\begin{kfChoices}
\kfChoice 밖의 더운 공기가 들어오는 것.
\kfChoice 창고 안의 얼음이 너무 빨리 녹는 것.
\kfChoice 창고 바닥에 물이 흥건하게 고이는 것.
\kfChoice 창고 문이 고장 나서 열리지 않는 것.
\end{kfChoices}
\end{kfSentenceUnit}

\begin{kfSentenceUnit}{26}{}
\kfSentQ{26.}{조상들은 밖의 더운 공기가 들어오는 문제를 어떻게 해결했나요?}
\begin{kfChoices}
\kfChoice 천장에 구멍을 뚫었다.
\kfChoice 바닥에 차가운 물을 뿌렸다.
\kfChoice 얼음 사이에 소금을 뿌렸다.
\kfChoice 창고 문을 두 개 더 만들었다.
\end{kfChoices}
\end{kfSentenceUnit}

\begin{kfSentenceUnit}{27}{}
\kfSentQ{27.}{따뜻한 공기는 어떤 특징이 있나요?}
\begin{kfChoices}
\kfChoice 가벼워서 위로 올라간다.
\kfChoice 차가워서 바닥에 깔린다.
\kfChoice 무거워서 아래로 내려온다.
\kfChoice 색깔이 있어서 눈에 보인다.
\end{kfChoices}
\end{kfSentenceUnit}

\begin{kfSentenceUnit}{28}{}
\kfSentQ{28.}{차가운 공기는 어떤 특징이 있나요?}
\begin{kfChoices}
\kfChoice 가벼워서 위로 올라간다.
\kfChoice 뜨거워서 천장에 붙는다.
\kfChoice 무거워서 아래로 내려온다.
\kfChoice 냄새가 나서 코를 찌른다.
\end{kfChoices}
\end{kfSentenceUnit}

\begin{kfSentenceUnit}{29}{}
\kfSentQ{29.}{따뜻한 공기와 차가운 공기의 차이점은 무엇인가요?}
\begin{kfChoices}
\kfChoice 따뜻한 공기는 가볍고, 차가운 공기는 무겁다는 점.
\kfChoice 따뜻한 공기는 색깔이 붉고, 차가운 공기는 푸르다는 점.
\kfChoice 따뜻한 공기는 맛이 달고, 차가운 공기는 맛이 쓰다는 점.
\kfChoice 따뜻한 공기는 냄새가 나고, 차가운 공기는 냄새가 없다는 점.
\end{kfChoices}
\end{kfSentenceUnit}

\begin{kfSentenceUnit}{30}{}
\kfSentQ{30.}{무엇이 천장의 구멍으로 나가나요?}
\begin{kfChoices}
\kfChoice 빗물.
\kfChoice 작은 새들.
\kfChoice 따뜻한 공기.
\kfChoice 차가운 공기.
\end{kfChoices}
\end{kfSentenceUnit}

\begin{kfSentenceUnit}{31}{}
\kfSentQ{31.}{따뜻한 공기가 나가면 석빙고 안이 어떻게 되나요?}
\begin{kfChoices}
\kfChoice 아주 깜깜해진다.
\kfChoice 점점 더 더워진다.
\kfChoice 습기가 가득 찬다.
\kfChoice 계속 시원하게 유지된다.
\end{kfChoices}
\end{kfSentenceUnit}

\begin{kfSentenceUnit}{32}{}
\kfSentQ{32.}{석빙고 안이 시원하게 유지되는 이유는 무엇인가요?}
\begin{kfChoices}
\kfChoice 얼음이 밖의 열을 계속 빨아들이기 때문이다.
\kfChoice 차가운 공기가 계속 밖으로 빠져나가기 때문이다.
\kfChoice 바닥에 흐르는 지하수가 열을 식혀주기 때문이다.
\kfChoice 천장 구멍으로 따뜻한 공기가 빠져나가기 때문이다.
\end{kfChoices}
\end{kfSentenceUnit}

\begin{kfSentenceUnit}{33}{}
\kfSentQ{33.}{'이 구멍'은 어떤 구멍인가요?}
\begin{kfChoices}
\kfChoice 벽에 뚫은 환기 구멍.
\kfChoice 문에 뚫은 열쇠 구멍.
\kfChoice 바닥에 뚫은 배수 구멍.
\kfChoice 석빙고 천장에 뚫은 구멍.
\end{kfChoices}
\end{kfSentenceUnit}

\begin{kfSentenceUnit}{34}{}
\kfSentQ{34.}{구멍에 무엇이 들어오지 않도록 했나요?}
\begin{kfChoices}
\kfChoice 비나 햇빛.
\kfChoice 신선한 공기.
\kfChoice 작은 동물들.
\kfChoice 사람들의 시선.
\end{kfChoices}
\end{kfSentenceUnit}

\begin{kfSentenceUnit}{35}{}
\kfSentQ{35.}{무엇으로 덮개를 만들었나요?}
\begin{kfChoices}
\kfChoice 돌.
\kfChoice 철.
\kfChoice 나무.
\kfChoice 유리.
\end{kfChoices}
\end{kfSentenceUnit}

\begin{kfSentenceUnit}{36}{}
\kfSentQ{36.}{천장 구멍에 덮개를 만든 이유는 무엇인가요?}
\begin{kfChoices}
\kfChoice 새들이 둥지를 틀지 못하게 하기 위해.
\kfChoice 비나 햇빛이 들어오지 않도록 하기 위해.
\kfChoice 사람들이 구멍으로 들어가지 못하게 하기 위해.
\kfChoice 창고 안의 찬 공기가 밖으로 새지 않게 하기 위해.
\end{kfChoices}
\end{kfSentenceUnit}

\begin{kfSentenceUnit}{37}{}
\kfSentQ{37.}{얼음을 더 오래 보관하기 위해 어떻게 했나요?}
\begin{kfChoices}
\kfChoice 얼음을 천으로 감쌌다.
\kfChoice 얼음 사이에 소금을 뿌렸다.
\kfChoice 얼음을 나무 상자에 넣었다.
\kfChoice 얼음 사이사이에 짚을 넣었다.
\end{kfChoices}
\end{kfSentenceUnit}

\begin{kfSentenceUnit}{38}{}
\kfSentQ{38.}{짚은 어떤 성질이 있나요?}
\begin{kfChoices}
\kfChoice 불에 아주 잘 타는 성질.
\kfChoice 아주 단단하고 튼튼한 성질.
\kfChoice 열이 잘 전달되지 않는 성질.
\kfChoice 물이 아주 잘 스며드는 성질.
\end{kfChoices}
\end{kfSentenceUnit}

\begin{kfSentenceUnit}{39}{}
\kfSentQ{39.}{얼음끼리 직접 닿아 있으면 열이 어떻게 되나요?}
\begin{kfChoices}
\kfChoice 쉽게 옮겨간다.
\kfChoice 사라져 버린다.
\kfChoice 그대로 멈춘다.
\kfChoice 소멸해 버린다.
\end{kfChoices}
\end{kfSentenceUnit}

\begin{kfSentenceUnit}{40}{}
\kfSentQ{40.}{짚이 가운데 있으면 열이 어떻게 되나요?}
\begin{kfChoices}
\kfChoice 천천히 이동한다.
\kfChoice 빠르게 이동한다.
\kfChoice 갑자기 사라진다.
\kfChoice 폭발적으로 늘어난다.
\end{kfChoices}
\end{kfSentenceUnit}

\begin{kfSentenceUnit}{41}{}
\kfSentQ{41.}{'짚이 있을 때'와 '없을 때'의 열 전달 속도는 어떻게 다른가요?}
\begin{kfChoices}
\kfChoice 짚이 있거나 없거나 열의 이동 속도는 똑같다.
\kfChoice 짚이 있으면 열이 아예 이동하지 않고, 없으면 폭발한다.
\kfChoice 짚이 있으면 열이 천천히 이동하고, 없으면 쉽게 옮겨간다.
\kfChoice 짚이 있으면 열이 빨리 이동하고, 없으면 천천히 이동한다.
\end{kfChoices}
\end{kfSentenceUnit}

\begin{kfSentenceUnit}{42}{}
\kfSentQ{42.}{짚에는 무엇이 많이 있나요?}
\begin{kfChoices}
\kfChoice 작은 씨앗.
\kfChoice 단단한 껍질.
\kfChoice 가느다란 실.
\kfChoice 작은 공기 구멍.
\end{kfChoices}
\end{kfSentenceUnit}

\begin{kfSentenceUnit}{43}{}
\kfSentQ{43.}{공기 구멍이 많으면 어떤 효과가 있나요?}
\begin{kfChoices}
\kfChoice 보기에 아름답다.
\kfChoice 단열 효과가 크다.
\kfChoice 물을 잘 흡수한다.
\kfChoice 소리를 잘 막아준다.
\end{kfChoices}
\end{kfSentenceUnit}

\begin{kfSentenceUnit}{44}{}
\kfSentQ{44.}{어디에 특별한 장치들이 있었나요?}
\begin{kfChoices}
\kfChoice 석빙고 주변.
\kfChoice 석빙고 내부.
\kfChoice 석빙고 천장.
\kfChoice 석빙고 바닥.
\end{kfChoices}
\end{kfSentenceUnit}

\begin{kfSentenceUnit}{45}{}
\kfSentQ{45.}{이 문장은 앞으로 어떤 내용이 나올 것을 알려주나요?}
\begin{kfChoices}
\kfChoice 석빙고를 관리하는 규칙들.
\kfChoice 석빙고에 보관했던 얼음의 종류.
\kfChoice 석빙고를 만든 사람들에 대한 소개.
\kfChoice 석빙고 건물 자체 외에 주변에 설치한 장치들.
\end{kfChoices}
\end{kfSentenceUnit}

\begin{kfSentenceUnit}{46}{}
\kfSentQ{46.}{어디에 흙을 덮었나요?}
\begin{kfChoices}
\kfChoice 건물 위.
\kfChoice 건물 안.
\kfChoice 건물 문.
\kfChoice 건물 바닥.
\end{kfChoices}
\end{kfSentenceUnit}

\begin{kfSentenceUnit}{47}{}
\kfSentQ{47.}{무엇을 막으려고 했나요?}
\begin{kfChoices}
\kfChoice 비가 올 때 물이 새는 것.
\kfChoice 밖의 열이 안으로 들어오는 것.
\kfChoice 안의 냉기가 밖으로 나가는 것.
\kfChoice 사람들이 건물 안을 들여다보는 것.
\end{kfChoices}
\end{kfSentenceUnit}

\begin{kfSentenceUnit}{48}{}
\kfSentQ{48.}{건물 위에 흙을 덮은 이유는 무엇인가요?}
\begin{kfChoices}
\kfChoice 건물을 예쁘게 꾸미기 위해.
\kfChoice 건물을 튼튼하게 만들기 위해.
\kfChoice 동물들이 접근하지 못하게 하기 위해.
\kfChoice 밖의 열이 안으로 들어오는 것을 막기 위해.
\end{kfChoices}
\end{kfSentenceUnit}

\begin{kfSentenceUnit}{49}{}
\kfSentQ{49.}{건물 위에 무엇을 심었나요?}
\begin{kfChoices}
\kfChoice 풀.
\kfChoice 꽃.
\kfChoice 나무.
\kfChoice 농작물.
\end{kfChoices}
\end{kfSentenceUnit}

\begin{kfSentenceUnit}{50}{}
\kfSentQ{50.}{흙 위에 풀을 심은 이유는 무엇인가요?}
\begin{kfChoices}
\kfChoice 보기에 아름답게 꾸미기 위해.
\kfChoice 벌레가 꼬이지 않게 하기 위해.
\kfChoice 태양열이 직접 닿지 않도록 하기 위해.
\kfChoice 뿌리가 흙을 단단하게 잡도록 하기 위해.
\end{kfChoices}
\end{kfSentenceUnit}

\begin{kfSentenceUnit}{51}{}
\kfSentQ{51.}{얼음을 저장하는 방은 어디에 만들었나요?}
\begin{kfChoices}
\kfChoice 강가 옆.
\kfChoice 땅속 깊이.
\kfChoice 산 꼭대기.
\kfChoice 건물 2층.
\end{kfChoices}
\end{kfSentenceUnit}

\begin{kfSentenceUnit}{52}{}
\kfSentQ{52.}{얼음 저장실을 땅속 깊이 만든 이유는 무엇인가요?}
\begin{kfChoices}
\kfChoice 지진에 대비하기 위해.
\kfChoice 적의 침입을 막기 위해.
\kfChoice 비밀스럽게 보관하기 위해.
\kfChoice 온도를 더 낮게 유지하기 위해.
\end{kfChoices}
\end{kfSentenceUnit}

\begin{kfSentenceUnit}{53}{}
\kfSentQ{53.}{석빙고를 무엇이라고 표현했나요?}
\begin{kfChoices}
\kfChoice 천연 냉장고.
\kfChoice 거대한 보물 창고.
\kfChoice 신비한 지하 동굴.
\kfChoice 위험한 비밀 요새.
\end{kfChoices}
\end{kfSentenceUnit}

\begin{kfSentenceUnit}{54}{}
\kfSentQ{54.}{석빙고는 무엇이 담긴 것인가요?}
\begin{kfChoices}
\kfChoice 조상들의 과학적 지혜.
\kfChoice 조상들의 예술적 재능.
\kfChoice 조상들의 문학적 소양.
\kfChoice 조상들의 종교적 믿음.
\end{kfChoices}
\end{kfSentenceUnit}

\begin{kfSentenceUnit}{55}{}
\kfSentQ{55.}{'이런 기술'은 무엇을 가리키나요?}
\begin{kfChoices}
\kfChoice 석빙고에 사용된, 얼음을 오랫동안 보관하기 위한 과학적 기술.
\kfChoice 건축에 사용된, 아름답고 웅장한 궁궐을 짓기 위해 발전시킨 건축 기술.
\kfChoice 농사에 사용된, 많은 양의 곡식을 빠르고 편하게 수확하기 위한 농업 기술.
\kfChoice 전쟁에 사용된, 적의 공격으로부터 나라를 튼튼하게 지키기 위한 무기 기술.
\end{kfChoices}
\end{kfSentenceUnit}

\begin{kfSentenceUnit}{56}{}
\kfSentQ{56.}{글쓴이가 석빙고 기술을 '놀라운 일'이라고 평가한 이유는 무엇인가요?}
\begin{kfChoices}
\kfChoice 많은 돈을 들여서 만들었기 때문이다.
\kfChoice 전기도 없던 시대에 만들어냈기 때문이다.
\kfChoice 아주 많은 사람들이 동원되었기 때문이다.
\kfChoice 외국의 기술을수입해서 만들었기 때문이다.
\end{kfChoices}
\end{kfSentenceUnit}



\end{document}
